% Leçon 18 — Expression écrite : protection de l'environnement (你应该怎么保护学校的环境) — Chinois 1ère A — Cours signé OnBuch+
% Programme officiel MINESEC. Compiler avec : tectonic ZHA18-ecrit-protection-de-l-environnement.tex
\newcommand{\DOCMATIERE}{Chinois}
\newcommand{\DOCNIVEAU}{1ère A}
\newcommand{\DOCMODULE}{Module 2 : Environnement et santé}
\newcommand{\DOCLECON}{ZHA18}
\newcommand{\DOCTITRE}{Leçon 18 — Expression écrite : protection de l'environnement (你应该怎么保护学校的环境)}
% OnBuch+ — préambule « cours signés » ENRICHI (compilé avec tectonic / XeLaTeX)
% Système visuel « Pop » d'OnBuch+ : fond crème (jamais blanc), bordures encre
% épaisses, ombres « dures » décalées, titres Archivo Black en capitales.
% Version MATHÉMATIQUES : graphes (pgfplots/tikz), boîtes pédagogiques
% (définition, propriété, méthode, attention, le savais-tu, exercice résolu,
% exercice, TP, bilan), page de garde et signature OnBuch+ ancrée en bas.
%
% Macros attendues dans main.tex AVANT \input{preamble} :
%   \DOCMATIERE  \DOCNIVEAU  \DOCMODULE  \DOCLECON (numéro)  \DOCTITRE
\documentclass[11pt]{article}

\usepackage{fontspec}
\usepackage[french]{babel} % français = langue principale ; le chinois via \zh{..} / textebox (xeCJK)
\usepackage{xeCJK}
\usepackage{pifont} % \ding{51} \ding{55} utilisés par les modèles
\usepackage[a4paper,margin=2cm,top=3.4cm,bottom=2.8cm,headheight=1.4cm,footskip=1.3cm]{geometry}
\usepackage{amsmath,amssymb,amsfonts}
\usepackage{mathtools} % \xmapsto, \coloneqq... souvent utilisés par les modèles
\usepackage{cancel} % simplifications barrées (bilans de forces)
% \abs{z} (module, valeur absolue) et \norm{u} (norme), utilisés par les modèles
\providecommand{\abs}{}\let\abs\relax\DeclarePairedDelimiter{\abs}{\lvert}{\rvert}
\providecommand{\norm}{}\let\norm\relax\DeclarePairedDelimiter{\norm}{\lVert}{\rVert}
\usepackage[table]{xcolor} % [table] : fournit aussi \rowcolors (alternance de lignes)
\usepackage{pagecolor}
\usepackage{tikz}
\usetikzlibrary{arrows.meta,positioning,calc,shapes.geometric,shapes.misc,shapes.arrows,decorations.pathmorphing,decorations.pathreplacing,decorations.markings,patterns,shadows,mindmap,trees,fit,backgrounds,matrix,intersections,angles,quotes}
\usepackage{pgfplots}
\usepackage[version=4]{mhchem} % équations nucléaires \ce{^{235}_{92}U}
\usepackage[european,siunitx]{circuitikz} % schémas électriques
\pgfplotsset{compat=1.18}
\usepgfplotslibrary{fillbetween,groupplots} % aires entre courbes (calcul intégral)
\usepackage{enumitem}
\usepackage{fancyhdr}
% [most] charge toutes les bibliothèques tcolorbox (listings, minted, raster,
% documentation...) et gonfle inutilement la mémoire TeX sur un document long
% (~300 boîtes) : on ne charge que ce qui est réellement utilisé.
\usepackage[skins,breakable]{tcolorbox}
\usepackage{graphicx}
\usepackage{colortbl}
\usepackage{booktabs}
\usepackage{tabularx}
\usepackage{diagbox} % case d'en-tête diagonale des tableaux à double entrée
% Colonnes extensibles centrée / gauche / droite (C, L, R), souvent utilisées par les modèles
\newcolumntype{C}{>{\centering\arraybackslash}X}
\newcolumntype{L}{>{\raggedright\arraybackslash}X}
\newcolumntype{R}{>{\raggedleft\arraybackslash}X}
\usepackage{array}
\usepackage{multicol}
\usepackage{multirow}
\usepackage{siunitx}
% parse-numbers=false : accepte aussi \SI{2,5 \times 10^{-3}}{...}
\sisetup{locale=FR,output-decimal-marker={,},group-minimum-digits=4,parse-numbers=false}
\DeclareSIUnit\ohm{\ensuremath{\symup{\Omega}}} % U+2126 absent de la police
\DeclareSIUnit\division{div} % réglages d'oscilloscope (ms/div, V/div)
\DeclareSIUnit\tr{tr} % tours (vitesse de rotation, ex. \SI{1500}{\tr\per\minute})
\DeclareSIUnit\molar{M} % concentration molaire (ex. \SI{3}{\molar}, \SI{1}{\micro\molar})
\DeclareSIUnit\an{an} % année (le modèle confond parfois avec \year, un compteur TeX) — ex. \SI{5}{\kilo\meter\per\an}
\DeclareSIUnit\FCFA{FCFA} % franc CFA (ex. \SI{500}{\FCFA\per\kilo\gram})
\DeclareSIUnit\degreeBrix{\textdegree Brix} % teneur en sucre d'un sirop/jus (réfractométrie)
\DeclareSIUnit\month{mois} % le modèle utilise parfois \month, un compteur TeX, comme unité de durée
\usepackage{qrcode}
% \degree : commande fréquemment utilisée par les modèles pour un angle ou une
% température en dehors de \SI{}{} (non fournie par les paquets chargés).
\providecommand{\degree}{^{\circ}}
% \qed (fin de démonstration), utilisé par les modèles sans amsthm
\providecommand{\qed}{\hfill$\square$}
% \barre : double barre des valeurs interdites dans un tableau de variations (tkz-tab)
\providecommand{\barre}{\ensuremath{\|}}
% environnement proof (amsthm non chargé) utilisé par les modèles
\ifdefined\proof\else\newenvironment{proof}[1][Démonstration]{\par\noindent\textit{#1.}\ }{\qed\par}\fi
\usepackage{lastpage}
\usepackage{hyperref}
\hypersetup{hidelinks}

% ============================================================
% Palette « Pop » OnBuch+ — mêmes hex que la classe P (plus_ui.dart)
% ============================================================
\definecolor{poporange}{HTML}{F59321}
\definecolor{popdark}{HTML}{E07A0C}
\definecolor{popcream}{HTML}{FFF6E8}
\definecolor{popink}{HTML}{1C1714}
\definecolor{poppurple}{HTML}{7A5AE0}
\definecolor{popgreen}{HTML}{2E9E6A}
\definecolor{poppink}{HTML}{E0457B}
\definecolor{popblue}{HTML}{2F7DE1}
\definecolor{popgold}{HTML}{C98A12}
\definecolor{popmuted}{HTML}{8A7F76}
\definecolor{popline}{HTML}{EADFD2}
\definecolor{popgray}{HTML}{8A7F76}
\colorlet{popgrey}{popgray}
% Alias tolérants (le modèle nomme parfois les couleurs différemment)
\colorlet{popwhite}{white}
\colorlet{popblack}{popink}
\colorlet{popred}{poppink}
\colorlet{popyellow}{popgold}
% Teintes claires (fonds de tableaux/encadrés) — le modèle nomme parfois
% les couleurs avec un suffixe L (ex. popblueL) sans les avoir définies.
\colorlet{poporangeL}{poporange!20}
\colorlet{popdarkL}{popdark!20}
\colorlet{poppurpleL}{poppurple!20}
\colorlet{popgreenL}{popgreen!20}
\colorlet{poppinkL}{poppink!20}
\colorlet{popblueL}{popblue!20}
\colorlet{popgoldL}{popgold!20}
\colorlet{popredL}{popred!20}
\colorlet{popgrayL}{popgray!20}
% Teintes claires pour fonds de tableaux / zones
\colorlet{popblueL}{popblue!10!white}
\colorlet{poporangeL}{poporange!14!white}
\colorlet{popgreenL}{popgreen!12!white}
\colorlet{poppurpleL}{poppurple!12!white}
\colorlet{poppinkL}{poppink!12!white}
\colorlet{popgoldL}{popgold!14!white}

\pagecolor{popcream}

% ============================================================
% Typographie
% ============================================================
\defaultfontfeatures{Ligatures=TeX}
\setmainfont{PlusJakartaSans-Medium.ttf}[Path=../../../fonts/,BoldFont=PlusJakartaSans-Bold.ttf]
% Symboles, API et emojis absents de la police principale : repli sur DejaVu Sans (caractères actifs)
\newfontface\symfont{DejaVuSans.ttf}[Path=../../../fonts/]
\makeatletter
\newcommand\symactive[1]{\catcode#1=13 \begingroup\lccode`\~=#1 \lowercase{\endgroup\def~}{{\symfont\char#1}}}
\newcommand\symblank[1]{\catcode#1=13 \begingroup\lccode`\~=#1 \lowercase{\endgroup\def~}{}}
\newcommand\symhyph[1]{\catcode#1=13 \begingroup\lccode`\~=#1 \lowercase{\endgroup\def~}{\mbox{-}}}
\newcount\symc
\newcommand\symrange[2]{\symc=#1 \@whilenum\symc<#2 \do{\expandafter\symactive\expandafter{\the\symc}\advance\symc by 1 }}
\newcommand\symblankrange[2]{\symc=#1 \@whilenum\symc<#2 \do{\expandafter\symblank\expandafter{\the\symc}\advance\symc by 1 }}
\makeatother
\symrange{"0250}{"02FF}\symactive{"2713}\symactive{"2714}\symactive{"2717}\symactive{"2718}\symactive{"2610}\symactive{"2611}\symactive{"2612}
\symrange{"2600}{"27BF}
\catcode"2705=13 \begingroup\lccode`\~="2705 \lowercase{\endgroup\def~}{{\symfont\char"2713}}\symblank{"26BD}\symblank{"26A1}
\symrange{"2460}{"257F}\symactive{"223C}\symactive{"2248}\symactive{"2260}
\symactive{"01F8}\symactive{"01F9}\symactive{"1E3E}\symactive{"1E3F}
\symhyph{"2011}\symhyph{"2E1A}\symblankrange{"1F300}{"1FAFF}\symblank{"FE0F}
% Caractères chinois : WenQuanYi Zen Hei (sous-ensemble GB2312 + pinyin), gras/italique simulés
\setCJKmainfont{WenQuanYiZenHei-subset.ttf}[Path=../../../fonts/,AutoFakeBold=3,AutoFakeSlant=0.2]
\xeCJKsetup{CJKspace=false}
% Pinyin : voyelles à caron (ǎ ǐ ǒ ǔ ǖ ǘ ǚ ǜ) absentes de la police principale -> caractères actifs
\newfontface\pyfont{WenQuanYiZenHei-subset.ttf}[Path=../../../fonts/]
\newcommand\pyactive[1]{\catcode#1=13 \begingroup\lccode`\~=#1 \lowercase{\endgroup\def~}{{\pyfont\char#1}}}
\pyactive{"01CD}\pyactive{"01CE}\pyactive{"01CF}\pyactive{"01D0}\pyactive{"01D1}\pyactive{"01D2}\pyactive{"01D3}\pyactive{"01D4}\pyactive{"01D5}\pyactive{"01D6}\pyactive{"01D7}\pyactive{"01D8}\pyactive{"01D9}\pyactive{"01DA}\pyactive{"01DB}\pyactive{"01DC}
% Maths en sans-serif (Fira Math) avec chiffres et lettres droites de Plus
% Jakarta, pour que les formules se fondent dans le texte.
\usepackage[math-style=ISO,bold-style=ISO]{unicode-math}
\setmathfont{FiraMath-Regular.otf}[Path=../../../fonts/]
\setmathfont{PlusJakartaSans-Medium.ttf}[Path=../../../fonts/,range={up/num,up/{Latin,latin},\mathup}]
\usepackage{icomma} % virgule décimale sans espace en mode math
% Caractères mathématiques Unicode absents des polices de texte (Plus Jakarta,
% Archivo Black) : rendus par leur équivalent mathématique, en texte comme en
% formule — sinon ils s'affichent en carré vide (ex. « Arithmétique dans ℤ »).
\usepackage{newunicodechar}
\newunicodechar{ℝ}{\ensuremath{\mathbb{R}}}
\newunicodechar{ℤ}{\ensuremath{\mathbb{Z}}}
\newunicodechar{ℂ}{\ensuremath{\mathbb{C}}}
\newunicodechar{ℕ}{\ensuremath{\mathbb{N}}}
\newunicodechar{ℚ}{\ensuremath{\mathbb{Q}}}
\newunicodechar{𝔻}{\ensuremath{\mathbb{D}}}
\newunicodechar{α}{\ensuremath{\alpha}}
\newunicodechar{β}{\ensuremath{\beta}}
\newunicodechar{γ}{\ensuremath{\gamma}}
\newunicodechar{δ}{\ensuremath{\delta}}
\newunicodechar{ε}{\ensuremath{\varepsilon}}
\newunicodechar{θ}{\ensuremath{\theta}}
\newunicodechar{λ}{\ensuremath{\lambda}}
\newunicodechar{μ}{\ensuremath{\mu}}
\newunicodechar{σ}{\ensuremath{\sigma}}
\newunicodechar{τ}{\ensuremath{\tau}}
\newunicodechar{φ}{\ensuremath{\varphi}}
\newunicodechar{ψ}{\ensuremath{\psi}}
\newunicodechar{ω}{\ensuremath{\omega}}
\newunicodechar{Σ}{\ensuremath{\Sigma}}
% majuscules produites par \MakeUppercase dans les titres (α-aminés -> Α)
\newunicodechar{Α}{\ensuremath{\alpha}}
\newunicodechar{Β}{\ensuremath{\beta}}
\newunicodechar{Γ}{\ensuremath{\Gamma}}
\newunicodechar{Δ}{\ensuremath{\Delta}}
\newunicodechar{Ε}{\ensuremath{\varepsilon}}
\newunicodechar{Θ}{\ensuremath{\Theta}}
\newunicodechar{Λ}{\ensuremath{\Lambda}}
\newunicodechar{Μ}{\ensuremath{\mu}}
\newunicodechar{Τ}{\ensuremath{\tau}}
\newunicodechar{Φ}{\ensuremath{\Phi}}
\newunicodechar{Ψ}{\ensuremath{\Psi}}
\newunicodechar{Ω}{\ensuremath{\Omega}}
\newunicodechar{↦}{\ensuremath{\mapsto}}
\newunicodechar{∈}{\ensuremath{\in}}
\newunicodechar{∉}{\ensuremath{\notin}}
\newunicodechar{∧}{\ensuremath{\wedge}}
\newunicodechar{⇔}{\ensuremath{\Leftrightarrow}}
\newunicodechar{⟺}{\ensuremath{\Longleftrightarrow}}
\newunicodechar{⇒}{\ensuremath{\Rightarrow}}
\newunicodechar{⟹}{\ensuremath{\Longrightarrow}}
\newunicodechar{∘}{\ensuremath{\circ}}
\newunicodechar{∪}{\ensuremath{\cup}}
\newunicodechar{∩}{\ensuremath{\cap}}
\newunicodechar{≡}{\ensuremath{\equiv}}
\newunicodechar{⊂}{\ensuremath{\subset}}
\newunicodechar{⊥}{\ensuremath{\perp}}
\newunicodechar{∃}{\ensuremath{\exists}}
\newunicodechar{∀}{\ensuremath{\forall}}
\newunicodechar{∛}{\ensuremath{\sqrt[3]{\,}}}
\newunicodechar{∜}{\ensuremath{\sqrt[4]{\,}}}
\newunicodechar{⃗}{\ensuremath{{}^{\to}}}
\newunicodechar{ⁿ}{\ifmmode^{n}\else\textsuperscript{\ensuremath{n}}\fi}
\newunicodechar{ˣ}{\ifmmode^{x}\else\textsuperscript{\ensuremath{x}}\fi}
\newunicodechar{ᵇ}{\ifmmode^{b}\else\textsuperscript{\ensuremath{b}}\fi}
\newunicodechar{ʳ}{\ifmmode^{r}\else\textsuperscript{\ensuremath{r}}\fi}
\newunicodechar{ᵘ}{\ifmmode^{u}\else\textsuperscript{\ensuremath{u}}\fi}
\newunicodechar{ʸ}{\ifmmode^{y}\else\textsuperscript{\ensuremath{y}}\fi}
\newunicodechar{ᵃ}{\ifmmode^{a}\else\textsuperscript{\ensuremath{a}}\fi}
\newunicodechar{ᵉ}{\ifmmode^{e}\else\textsuperscript{\ensuremath{e}}\fi}
\newunicodechar{ᵏ}{\ifmmode^{k}\else\textsuperscript{\ensuremath{k}}\fi}
\newunicodechar{ᵖ}{\ifmmode^{p}\else\textsuperscript{\ensuremath{p}}\fi}
\newunicodechar{ᴺ}{\ifmmode^{N}\else\textsuperscript{\ensuremath{N}}\fi}
\newunicodechar{ᶜ}{\ifmmode^{c}\else\textsuperscript{\ensuremath{c}}\fi}
\newunicodechar{ʰ}{\ifmmode^{h}\else\textsuperscript{\ensuremath{h}}\fi}
\newunicodechar{ᵠ}{\ifmmode^{q}\else\textsuperscript{\ensuremath{q}}\fi}
\newunicodechar{ₙ}{\ifmmode_{n}\else\textsubscript{\ensuremath{n}}\fi}
\newunicodechar{ₐ}{\ifmmode_{a}\else\textsubscript{\ensuremath{a}}\fi}
\newunicodechar{ᵢ}{\ifmmode_{i}\else\textsubscript{\ensuremath{i}}\fi}
\newunicodechar{ᵣ}{\ifmmode_{r}\else\textsubscript{\ensuremath{r}}\fi}
\newunicodechar{ₖ}{\ifmmode_{k}\else\textsubscript{\ensuremath{k}}\fi}
\newunicodechar{ᵦ}{\ifmmode_{\beta}\else\textsubscript{\ensuremath{\beta}}\fi}
\newunicodechar{ₓ}{\ifmmode_{x}\else\textsubscript{\ensuremath{x}}\fi}
\newfontfamily\popdisplay{ArchivoBlack-Regular.ttf}[Path=../../../fonts/]
\newfontfamily\poplogo{SpaceGrotesk-Bold.ttf}[Path=../../../fonts/]
\newfontfamily\popsemi{PlusJakartaSans-SemiBold.ttf}[Path=../../../fonts/]
\newfontfamily\popxbold{PlusJakartaSans-ExtraBold.ttf}[Path=../../../fonts/]
\newfontfamily\popxboldls{PlusJakartaSans-ExtraBold.ttf}[Path=../../../fonts/,LetterSpace=3.0]

\setlist[enumerate]{leftmargin=1.7em,itemsep=5pt,topsep=4pt}
\setlist[itemize]{leftmargin=1.7em,itemsep=4pt,topsep=4pt,label={\color{poporange}\textbullet}}
\setlength{\parindent}{0pt}
\setlength{\parskip}{5pt}
\renewcommand{\arraystretch}{1.25}

% pgfplots : style Pop par défaut
\pgfplotsset{
  every axis/.append style={axis line style={popink,line width=1pt},tick style={popink},
    grid style={popline,line width=0.5pt},label style={font=\small},tick label style={font=\footnotesize}},
  popcurve/.style={poporange,line width=1.8pt,no markers,smooth},
}
% Même style disponible dans un \draw TikZ simple (hors pgfplots)
\tikzset{popcurve/.style={poporange,line width=1.8pt}}
\tikzset{popgrid/.style={popline,very thin}}
\colorlet{popcurve}{poporange} % le nom du style est parfois utilisé comme couleur

% ============================================================
% Pill / squiggle
% ============================================================
\newcommand{\poppill}[3]{%
  \tikz[baseline=-0.55ex]\node[draw=popink,line width=1.2pt,rounded corners=2.6pt,%
    fill=#2,inner xsep=6.5pt,inner ysep=3pt]{%
    \popxboldls\fontsize{7}{7}\selectfont\color{#3}\MakeUppercase{#1}};%
}
\newcommand{\popsquiggle}[1]{%
  \tikz[baseline=-0.4ex]{
    \draw[#1,line width=2.6pt,line cap=round]
      plot [smooth] coordinates {(0,0.09) (0.34,0.01) (0.68,0.21) (1.02,0.01) (1.36,0.10)};
  }%
}
% Mot-clé surligné (marqueur orange sous le texte)
\newcommand{\cle}[1]{{\popxbold\color{popdark}#1}}

% ============================================================
% En-tête / pied
% ============================================================
\pagestyle{fancy}
\fancyhf{}
\renewcommand{\headrulewidth}{0pt}
\renewcommand{\footrulewidth}{0pt}
\lhead{\raisebox{-0.15ex}{\poplogo\fontsize{13}{13}\selectfont\color{popink}OnBuch\color{poporange}+}}
\rhead{\poppill{\DOCMATIERE\ \textbullet\ \DOCNIVEAU\ \textbullet\ Leçon \DOCLECON}{white}{popink}}
\lfoot{\popxbold\fontsize{7.6}{7.6}\selectfont\color{popmuted}\MakeUppercase{Cours signé OnBuch+}\ \textbullet\ {\popsemi\DOCTITRE}}
\rfoot{\tikz[baseline=-0.6ex]\node[rounded corners=8.5pt,fill=popink,minimum height=17pt,minimum width=17pt,inner xsep=5pt,inner sep=0]{\popxbold\fontsize{8}{8}\selectfont\color{popcream}\thepage\,/\,\pageref*{LastPage}};}

% ============================================================
% Titres
% ============================================================
\newcommand{\chaptitre}[2]{%
  \begin{tcolorbox}[enhanced,colback=poporange,colframe=popink,boxrule=2.6pt,arc=11pt,
      drop shadow=popink,left=15pt,right=15pt,top=12pt,bottom=11pt,before skip=3pt,after skip=18pt]
    \poppill{#2}{popink}{popcream}\par\vspace{9pt}
    {\raggedright\hyphenpenalty=10000\exhyphenpenalty=10000\popdisplay\fontsize{22}{24}\selectfont\color{popcream}\MakeUppercase{#1}\par}
  \end{tcolorbox}
}
% Section numérotée : pastille orange avec le numéro + titre Archivo
\newcounter{popsec}
\newcommand{\coursec}[1]{%
  \stepcounter{popsec}\setcounter{popsub}{0}%
  \par\vspace{16pt}\noindent
  \begin{minipage}{\linewidth}
  \tikz[baseline=-0.7ex]\node[draw=popink,line width=1.6pt,fill=poporange,rounded corners=4pt,minimum size=22pt,inner sep=0]{\popdisplay\fontsize{12}{12}\selectfont\color{popcream}\thepopsec};%
  \hspace{8pt}{\popdisplay\fontsize{15}{17}\selectfont\color{popink}\MakeUppercase{#1}}\par
  \vspace{4pt}\noindent{\color{poporange}\rule{36pt}{3.6pt}}
  \end{minipage}\par\nopagebreak\vspace{8pt}
}
\newcounter{popsub}
\newcommand{\courssub}[1]{%
  \stepcounter{popsub}%
  \par\vspace{9pt}\noindent{\popxbold\fontsize{11.5}{13}\selectfont\color{popink}{\color{poporange}\thepopsec.\thepopsub}\hspace{6pt}#1}\par\nopagebreak\vspace{4pt}
}

% ============================================================
% Boîtes pédagogiques — même recette : fond blanc, bordure épaisse colorée,
% ombre dure encre, badge plein. Titre optionnel [#1].
% ============================================================
\tcbset{popbase/.style={breakable,enhanced,colback=white,boxrule=2.3pt,arc=9pt,
  shadow={3pt}{-3pt}{0pt}{popink},left=14pt,right=14pt,top=11pt,bottom=11pt,before skip=11pt,after skip=15pt,
  fontupper=\popsemi\fontsize{10.6}{14.5}\selectfont\color{popink},
  fontlower=\popsemi\fontsize{10.6}{14.5}\selectfont\color{popink}}}
% Coche dessinée (le glyphe ✓ n'existe pas dans les polices du document)
\newcommand{\popcheck}[1]{\tikz[baseline=-0.5ex]\draw[#1,line width=1.6pt,line cap=round,line join=round](-0.13,0.0)--(-0.04,-0.09)--(0.13,0.1);}
\providecommand{\coche}{\popcheck{popgreen}}
\providecommand{\resultat}[1]{\par\smallskip\noindent\fcolorbox{popgreen}{popgreenL}{\parbox{\dimexpr\linewidth-2\fboxsep-2\fboxrule\relax}{\textbf{Résultat.} #1}}\par\smallskip}
\newcommand{\popboxhead}[3]{\poppill{#1}{#2}{white}\ifx\relax#3\relax\else\hspace{6pt}{\popxbold\fontsize{10.6}{12}\selectfont\color{popink}#3}\fi\par\vspace{7pt}}

\newtcolorbox{aretenir}[1][]{popbase,colframe=popgreen,before upper={\popboxhead{À retenir}{popgreen}{#1}}}
% Texte en chinois (document de lecture, dialogue, extrait) : cadre
% d'étude. Un titre optionnel (source, auteur) s'affiche dans l'en-tête.
\newtcolorbox{textebox}[1][]{popbase,colframe=popblue,before upper={\popboxhead{文本 · Texte}{popblue}{#1}},after upper={}}
% Marques ✓ / ✗ (forme correcte / fautive) souvent utilisées par les modèles
\providecommand{\cross}{\textcolor{poppink}{\ensuremath{\times}}}
\providecommand{\xmark}{\cross}\providecommand{\crossmark}{\cross}\providecommand{\cmark}{\ensuremath{\checkmark}}
% Ligne de réponse / justification à compléter (inventée par les modèles)
\providecommand{\justification}[1]{\par\noindent\textit{Justification~:} \dotfill\par}
% \textipa (package tipa non chargé) : la police ne couvre pas l'API -> simple passage
\providecommand{\textipa}[1]{#1}
% Soulignement (ulem non chargé) : \uline, \ul
\providecommand{\uline}[1]{\underline{#1}}\providecommand{\ul}[1]{\underline{#1}}
% \glossaire{\item ...} inventée par les modèles : liste de glossaire
\providecommand{\glossaire}[1]{\begin{itemize}#1\end{itemize}}
% \alert (beamer) inventée par les modèles : texte mis en évidence
\providecommand{\alert}[1]{\textbf{\textcolor{poppink}{#1}}}
% variantes ASCII d'ordinaux écrites en macro
\providecommand{\iere}{\textsuperscript{re}}\providecommand{\ieme}{\textsuperscript{e}}\providecommand{\ere}{\textsuperscript{re}}\providecommand{\eme}{\textsuperscript{e}}\providecommand{\er}{\textsuperscript{er}}\providecommand{\ie}{\textsuperscript{e}}
% Ordinaux français écrits en macro par les modèles : 1\ière, 3\ième
\newcommand{\ière}{\textsuperscript{re}}\newcommand{\ième}{\textsuperscript{e}}
% \exemple (inventée par les modèles) : étiquette « Exemple : »
\providecommand{\exemple}{\textbf{Exemple~:}\ }
% \square utilisable aussi hors mode math (cases à cocher)
\AtBeginDocument{\let\squareMath\square\renewcommand{\square}{\ensuremath{\squareMath}}}
% Mot ou phrase en chinois à l'intérieur d'un paragraphe français
\newcommand{\zh}[1]{\ifmmode\text{#1}\else{#1}\fi}
\let\eng\zh \let\esp\zh \let\all\zh % alias tolérés
% flèches d'intonation inventées par les modèles
\providecommand{\textsearrow}{}\renewcommand{\textsearrow}{\ensuremath{\searrow}}\providecommand{\textnearrow}{}\renewcommand{\textnearrow}{\ensuremath{\nearrow}}
% \fr{..} : mot français (inventé par les modèles)
\providecommand{\fr}[1]{\textit{#1}}
% zhbox : encadré de texte chinois inventé par les modèles
\newenvironment{zhbox}{\begin{quote}}{\end{quote}}
% \vs : « versus » inventé par les modèles
\providecommand{\vs}{\textit{vs}\ }
% \blank : case à compléter inventée par les modèles
\providecommand{\blank}[1][]{\underline{\hspace{1.6em}}}
\newtcolorbox{exemplebox}[1][]{popbase,colframe=poporange,before upper={\popboxhead{Exemple}{poporange}{#1}}}
\newtcolorbox{definition}[1][]{popbase,colframe=popblue,before upper={\popboxhead{Définition}{popblue}{#1}}}
\newtcolorbox{propriete}[1][]{popbase,colframe=poppurple,before upper={\popboxhead{Propriété}{poppurple}{#1}}}
\newtcolorbox{methode}[1][]{popbase,colframe=popgold,before upper={\popboxhead{Méthode}{popgold}{#1}}}
\newtcolorbox{attention}[1][]{popbase,colframe=poppink,before upper={\popboxhead{Attention}{poppink}{#1}}}
\newtcolorbox{experience}[1][]{popbase,colframe=popblue,colback=popblueL,before upper={\popboxhead{Expérience}{popblue}{#1}}}
% « Le savais-tu ? » : carte violette pleine (contexte, culture, Cameroun)
\newtcolorbox{savaistu}[1][]{popbase,colback=poppurple,colframe=popink,
  fontupper=\popsemi\fontsize{10.4}{14.2}\selectfont\color{white},
  before upper={\poppill{Le savais-tu\,?}{white}{poppurple}\ifx\relax#1\relax\else\hspace{6pt}{\popxbold\color{white}#1}\fi\par\vspace{7pt}}}
% Exercice résolu : énoncé en haut, \tcblower puis la solution sur fond crème
\newcounter{popexo}\newcounter{popexr}
\newtcolorbox{exoresolu}[1][]{popbase,colframe=poporange,colbacklower=popcream,
  segmentation style={popink,line width=1.2pt,dashed},
  before upper={\stepcounter{popexr}\popboxhead{Exercice résolu \thepopexr}{poporange}{#1}},
  before lower={\poppill{Solution}{popink}{popcream}\par\vspace{6pt}}}
% Exercice (non résolu en place) — la correction est dans la partie Corrigés
\newtcolorbox{exercice}[1][]{popbase,colframe=popink,
  before upper={\stepcounter{popexo}\popboxhead{Exercice \thepopexo}{popink}{#1}}}
% Corrigé
\newtcolorbox{corrige}[1][]{popbase,colframe=popgreen,colback=popgreenL,
  before upper={\popboxhead{Corrigé}{popgreen}{#1}}}
% Objectifs / prérequis / bilan
\newtcolorbox{objectifs}{popbase,colframe=popink,colback=poporangeL,
  before upper={\popboxhead{Objectifs de la leçon}{popink}{}}}
\newtcolorbox{prerequis}{popbase,colframe=popink,colback=popblueL,
  before upper={\popboxhead{Prérequis}{popblue}{}}}
\newtcolorbox{bilan}{popbase,colframe=popink,colback=white,boxrule=2.8pt,
  before upper={\popboxhead{Fiche bilan}{poporange}{L'essentiel à connaître pour l'examen}}}
% Figure encadrée avec légende
\newtcolorbox{popfigure}[1][]{enhanced,breakable=false,colback=white,colframe=popline,boxrule=1.4pt,arc=8pt,
  left=10pt,right=10pt,top=10pt,bottom=8pt,before skip=10pt,after skip=12pt,halign=center,
  lower separated=false,halign lower=center,
  fontlower=\popsemi\fontsize{9}{11.5}\selectfont\color{popmuted},
  #1}
\newcommand{\legende}[1]{\par\vspace{6pt}{\popsemi\fontsize{9}{11.5}\selectfont\color{popmuted}#1}}

% ============================================================
% Signature OnBuch+
% ============================================================
\newcommand{\poplogoinline}[1]{{\poplogo\fontsize{#1}{#1}\selectfont\color{popink}OnBuch\color{poporange}+}}
\newcommand{\signatureonbuch}{%
  \begin{tcolorbox}[enhanced,colback=popink,colframe=popink,boxrule=2.2pt,arc=10pt,
      drop shadow=poporange,left=14pt,right=14pt,top=10pt,bottom=10pt,before skip=0pt,after skip=0pt]
    \tikz[baseline=-0.7ex]\node[circle,fill=popgreen,draw=popcream,line width=1.3pt,minimum size=22pt,inner sep=0]{\popcheck{white}};%
    \hspace{9pt}\begin{minipage}[c]{0.62\linewidth}
      {\popxbold\fontsize{12}{13}\selectfont\color{popcream}Signé\ \textbullet\ {\poplogo OnBuch\color{poporange}+}}\par\vspace{2pt}
      {\raggedright\popsemi\fontsize{8}{10.5}\selectfont\color{popline}\MakeUppercase{Équipe pédagogique OnBuch+}\\\MakeUppercase{Conforme au programme officiel MINESEC}\par}
    \end{minipage}\hfill
    {\poplogo\fontsize{22}{22}\selectfont\color{popcream}OnBuch\color{poporange}+}
  \end{tcolorbox}%
}

% ============================================================
% Page de garde
% #1 titre · #2 matière · #3 niveau · #4 module · #5 n° leçon · #6 durée ·
% #7 description / objectifs (texte ou itemize)
% ============================================================
\newcommand{\pagegarde}[7]{%
  \thispagestyle{empty}
  \newgeometry{a4paper,margin=1.7cm,top=1.6cm,bottom=1.4cm}%
  \begin{tikzpicture}[remember picture,overlay]
    \fill[poporange!18] ([xshift=-3.2cm,yshift=-3.6cm]current page.north east) circle (4.6cm);
    \fill[poppurple!14] ([xshift=2.4cm,yshift=7.6cm]current page.south west) circle (3.8cm);
  \end{tikzpicture}%
  {\poplogo\fontsize{30}{30}\selectfont\color{popink}OnBuch\color{poporange}+}\hfill
  \raisebox{0.6ex}{\poppill{Cours signé \textbullet\ Premium}{popink}{popcream}}\par
  \vspace{1pt}\popsquiggle{poppurple}\par
  \vspace{8pt}
  \begin{tcolorbox}[enhanced,colback=poporange,colframe=popink,boxrule=2.8pt,arc=13pt,
      drop shadow=popink,left=18pt,right=18pt,top=12pt,bottom=13pt,after skip=10pt]
    \poppill{#2 \textbullet\ #3}{popink}{popcream}\hspace{5pt}\poppill{#4}{white}{popink}\par\vspace{8pt}
    {\popdisplay\fontsize{14}{14}\selectfont\color{popink}LEÇON #5}\par\vspace{4pt}
    {\raggedright\hyphenpenalty=10000\exhyphenpenalty=10000\popdisplay\fontsize{25}{27}\selectfont\color{popcream}\MakeUppercase{#1}\par}
    \vspace{8pt}
    {\popxbold\fontsize{9.5}{9.5}\selectfont\color{popink}\MakeUppercase{Durée conseillée : #6}}
  \end{tcolorbox}
  \begin{tcolorbox}[enhanced,colback=white,colframe=popink,boxrule=2.2pt,arc=10pt,
      drop shadow=popink,left=16pt,right=16pt,top=10pt,bottom=10pt,after skip=8pt]
    \poppill{Dans cette leçon}{poppurple}{white}\par\vspace{5pt}
    \popsemi\fontsize{9.3}{12.6}\selectfont\color{popink}#7
  \end{tcolorbox}
  \noindent\begin{minipage}[t]{0.487\linewidth}
    \begin{tcolorbox}[enhanced,colback=popgreen,colframe=popink,boxrule=2.2pt,arc=10pt,
        drop shadow=popink,left=11pt,right=11pt,top=10pt,bottom=10pt,height=3.1cm,valign=center]
      \tcbox[colback=white,colframe=popink,boxrule=1.3pt,arc=5pt,boxsep=0pt,nobeforeafter,
        left=3pt,right=3pt,top=3pt,bottom=3pt,valign=center]{\qrcode[height=1.9cm]{https://play.google.com/store/apps/details?id=com.luvvix.onbuch&hl=fr}}%
      \hspace{8pt}\begin{minipage}[c]{0.5\linewidth}
        {\popdisplay\fontsize{11}{12.5}\selectfont\color{white}\MakeUppercase{Télécharge OnBuch}}\par\vspace{4pt}
        {\raggedright\popsemi\fontsize{8.4}{11}\selectfont\color{white}Gratuit \textbullet\ Google Play\\Tuteur IA, annales, concours, résultats\par}
      \end{minipage}
    \end{tcolorbox}
  \end{minipage}\hfill
  \begin{minipage}[t]{0.487\linewidth}
    \begin{tcolorbox}[enhanced,colback=poppurple,colframe=popink,boxrule=2.2pt,arc=10pt,
        drop shadow=popink,left=11pt,right=11pt,top=10pt,bottom=10pt,height=3.1cm,valign=center]
      \tikz[baseline=-0.7ex]\node[circle,fill=white,draw=popink,line width=1.3pt,minimum size=1.6cm,inner sep=0]{\popdisplay\fontsize{24}{24}\selectfont\color{poppurple}+};%
      \hspace{8pt}\begin{minipage}[c]{0.6\linewidth}
        {\popdisplay\fontsize{11}{12.5}\selectfont\color{white}\MakeUppercase{Passe à OnBuch+}}\par\vspace{4pt}
        {\raggedright\popsemi\fontsize{8.4}{11}\selectfont\color{white}Tous les cours signés, Tuteur IA illimité, fiches PDF — onglet OnBuch+ de l'app\par}
      \end{minipage}
    \end{tcolorbox}
  \end{minipage}
  \vspace{8pt}
  \popcontenu
  \vfill
  \signatureonbuch
  \restoregeometry
  \newpage
}

% Rangée « Ce cours contient » (page de garde)
\newcommand{\poptile}[3]{% #1 couleur #2 gros texte #3 légende
  \begin{tcolorbox}[enhanced,colback=#1,colframe=popink,boxrule=2pt,arc=9pt,shadow={2.5pt}{-2.5pt}{0pt}{popink},
      left=6pt,right=6pt,top=9pt,bottom=9pt,halign=center,valign=center,height=1.75cm,width=\linewidth,nobeforeafter]
    {\popdisplay\fontsize{12.5}{13}\selectfont\color{white}\MakeUppercase{#2}}\par\vspace{3pt}
    {\centering\popsemi\fontsize{7.8}{9.5}\selectfont\color{white}#3\par}
  \end{tcolorbox}}
\newcommand{\popcontenu}{%
  \par\noindent{\popxboldls\fontsize{7.5}{7.5}\selectfont\color{popmuted}CE COURS SIGNÉ CONTIENT}\par\vspace{7pt}
  \noindent\begin{minipage}[t]{0.235\linewidth}\poptile{popblue}{Cours}{complet, illustré, pas à pas}\end{minipage}\hfill
  \begin{minipage}[t]{0.235\linewidth}\poptile{poporange}{Méthodes}{et exercices résolus}\end{minipage}\hfill
  \begin{minipage}[t]{0.235\linewidth}\poptile{poppink}{Exercices}{type examen + corrigés}\end{minipage}\hfill
  \begin{minipage}[t]{0.235\linewidth}\poptile{popgreen}{Fiche bilan}{l'essentiel à retenir}\end{minipage}\par}

% Sommaire
\newtcolorbox{sommaire}{enhanced,colback=white,colframe=popink,boxrule=2.2pt,arc=10pt,
  drop shadow=popink,left=18pt,right=18pt,top=13pt,bottom=13pt,before skip=6pt,after skip=20pt,
  before upper={\poppill{Sommaire}{poppurple}{white}\par\vspace{9pt}\popsemi\fontsize{10.6}{14.5}\selectfont\color{popink}}}

% Signature de fin de cours — poussée tout en bas de la dernière page
\newcommand{\findecours}{%
  \par\vfill\nopagebreak
  \begin{center}\popsquiggle{poporange}\end{center}\vspace{4pt}
  \signatureonbuch
}
\AtBeginDocument{\def\llbracket{\mathopen{[\mkern-3.5mu[}}\def\rrbracket{\mathclose{]\mkern-3.5mu]}}} % intervalles d'entiers (glyphe ⟦ absent de la police)
% Glyphes absents de Fira Math : redéfinis à partir de glyphes disponibles
\AtBeginDocument{%
  \def\setminus{\mathbin{\backslash}}\let\smallsetminus\setminus
  \def\vdots{\mathord{\vbox{\baselineskip3.2pt\lineskiplimit0pt\kern4pt\hbox{.}\hbox{.}\hbox{.}}}}%
  \def\ddots{\mathinner{\mkern1mu\raise6.5pt\hbox{.}\mkern2mu\raise3.6pt\hbox{.}\mkern2mu\raise0.7pt\hbox{.}\mkern1mu}}%
  \def\triangle{\mathord{\scalebox{0.9}{$\symup{\Delta}$}}}%
  \def\nabla{\mathord{\raisebox{\depth}{\scalebox{1}[-1]{$\symup{\Delta}$}}}}%
  \def\checkmark{\popcheck{popdark}}%
  \def\star{\mathbin{\ast}}\def\bigstar{\mathord{\ast}}%
  \def\diamond{\mathbin{\scalebox{0.6}{$\Diamond$}}}%
  \def\bigcup{\mathop{\mathchoice{\raisebox{-0.3em}{\scalebox{1.7}{$\cup$}}}{\scalebox{1.25}{$\cup$}}{\cup}{\cup}}\displaylimits}%
  \def\bigcap{\mathop{\mathchoice{\raisebox{-0.3em}{\scalebox{1.7}{$\cap$}}}{\scalebox{1.25}{$\cap$}}{\cap}{\cap}}\displaylimits}%
  \def\lesseqqgtr{\mathrel{\lessgtr}}\def\bigoplus{\mathop{\oplus}\displaylimits}%
}
\providecommand{\ssi}{si et seulement si} % abréviation fréquente
\AtBeginDocument{\providecommand{\wideparen}{\overparen}} % arc de cercle

%% FIGFIX-BEGIN v3 — correctifs globaux des figures (docs/onbuch-generation/tools/figfix)
\usepackage{adjustbox}\usepackage{etoolbox}
% toute figure TikZ/pgfplots plus large que la ligne est réduite à la largeur disponible
\BeforeBeginEnvironment{tikzpicture}{\begin{adjustbox}{max width=\linewidth}}
\AfterEndEnvironment{tikzpicture}{\end{adjustbox}}
% légendes pgfplots lisibles par-dessus les courbes
\pgfplotsset{every axis/.append style={legend style={fill=white,fill opacity=0.92,text opacity=1,draw=black!15,inner sep=3pt}}}
% étiquettes d'arêtes (syntaxe edge["..."]) avec fond blanc
\tikzset{every edge quotes/.append style={fill=white,inner sep=1.2pt}}
% bulles de cartes mentales : texte à la ligne dans la bulle, bulle un peu plus grande
\tikzset{every concept/.append style={text width=2.0cm,align=center,minimum size=2.3cm,inner sep=2pt}}
%% FIGFIX-END
\begin{document}

\pagegarde{Leçon 18 — Expression écrite : protection de l'environnement (你应该怎么保护学校的环境)}{Chinois}{1ère A}{Module 2 : Environnement et santé}{ZHA18}{2 h}{Cette leçon guide les élèves de 1ère A dans la rédaction d'un texte simple en chinois sur la protection de l'environnement scolaire, en intégrant l'écriture correcte des caractères, les connecteurs logiques et la traduction thème/version. Elle prépare aux épreuves écrites du baccalauréat tout en développant une conscience écologique.
\vspace{4pt}
\begin{multicols}{2}\small
\begin{itemize}
\item À la fin de cette leçon, je sais rédiger un paragraphe cohérent en chinois sur la protection de l'environnement scolaire.
\item Je sais identifier et utiliser les principaux connecteurs logiques (d'abord, ensuite, enfin, donc, parce que).
\item Je sais écrire correctement les caractères clés du thème (环, 保, 护, 学, 校, 环境) en respectant l'ordre des traits et les radicaux.
\item Je sais traduire du français vers le chinois (thème) et du chinois vers le français (version) des phrases simples sur le sujet.
\item Je sais évaluer mon propre texte à l'aide d'une grille de critères (contenu, structure, caractères, langue).
\item Je sais présenter oralement mon affiche en chinois avec une prononciation correcte des tons.
\end{itemize}\end{multicols}}

\begin{sommaire}
\begin{multicols}{2}
\begin{enumerate}
\item Modèle de texte commenté : « Comment protéger l'environnement de l'école »
\item Structure et connecteurs logiques
\item Écriture correcte des caractères du thème
\item Exercices de thème (français → chinois) et version (chinois → français)
\item Grille d'évaluation et auto‑correction
\item Activité d'intégration
\item Méthodes et exercices résolus
\item Exercices
\item Corrigés des exercices
\item Fiche bilan
\end{enumerate}
\end{multicols}
\end{sommaire}

\begin{objectifs}
\begin{itemize}
\item À la fin de cette leçon, je sais rédiger un paragraphe cohérent en chinois sur la protection de l'environnement scolaire.
\item Je sais identifier et utiliser les principaux connecteurs logiques (d'abord, ensuite, enfin, donc, parce que).
\item Je sais écrire correctement les caractères clés du thème (环, 保, 护, 学, 校, 环境) en respectant l'ordre des traits et les radicaux.
\item Je sais traduire du français vers le chinois (thème) et du chinois vers le français (version) des phrases simples sur le sujet.
\item Je sais évaluer mon propre texte à l'aide d'une grille de critères (contenu, structure, caractères, langue).
\item Je sais présenter oralement mon affiche en chinois avec une prononciation correcte des tons.
\item Je sais comparer les pratiques écologiques scolaires au Cameroun et en Chine.
\end{itemize}
\end{objectifs}

\begin{prerequis}
\begin{itemize}
\item Maîtrise des radicaux de base (氵, 土, 木, 手) et de l'ordre des traits élémentaires (Seconde).
\item Connaissance des structures simples SVO et des particules 的, 地, 得 (début d'année).
\item Vocabulaire de base sur l'école (学校, 老师, 同学) et l'environnement (环境, 垃圾, 回收).
\item Capacité à lire le pinyin avec les quatre tons.
\item Notions de connecteurs de base en français (d'abord, puis, enfin).
\end{itemize}
\end{prerequis}

\newpage

\coursec{Modèle de texte commenté : « Comment protéger l'environnement de l'école »}

\courssub{Situation d'accroche et problématique}

Il est dix heures au lycée de Yaoundé. La sonnerie vient de retentir : c'est la fin de la grande pause. Les élèves regagnent leurs salles de classe, mais la cour de récréation reste jonchée de sachets plastiques, de papiers et de bouteilles vides. Quelques élèves de Première A s'arrêtent, déçus : « Notre école était si belle ce matin ! » En discutant avec leurs camarades correspondants du Nigeria et du Ghana, ils réalisent que ce problème est commun à de nombreux pays d'Afrique de l'Ouest, anglophones comme francophones : partout, les déchets plastiques envahissent les cours d'école.

Leur professeur de chinois saisit l'occasion : « Et si vous écriviez, en chinois, une courte affiche pour sensibiliser toute l'école à la protection de l'environnement ? » Écrire en chinois sur ce sujet exige un vocabulaire précis (\zh{环境} environnement, \zh{保护} protéger, \zh{垃圾} déchets…), des connecteurs logiques (\zh{首先}, \zh{其次}, \zh{然后}, \zh{最后}, \zh{因此}) et une structure claire.

\textbf{Problématique :} comment rédiger en chinois un texte simple, cohérent et convaincant pour expliquer comment protéger l'environnement de son école ? Pour y répondre, nous allons d'abord lire et observer un texte modèle, puis en analyser la structure, le vocabulaire et les connecteurs, avant de nous entraîner à l'écriture des caractères et à la traduction.

\courssub{Le texte modèle : lecture et observation}

Lisons d'abord le texte modèle. Il compte environ 90 caractères : c'est la longueur idéale pour une production écrite de niveau Première (HSK 2-3).

\begin{textebox}[我们应该怎么保护学校的环境？]
我们应该保护学校的环境。首先，不要乱扔垃圾。其次，做好垃圾分类和回收。然后，多种树、多种花。最后，向同学们宣传环保知识。因此，学校会变得更美丽。\\
\emph{Wǒmen yīnggāi bǎohù xuéxiào de huánjìng. Shǒuxiān, bú yào luàn rēng lājī. Qícì, zuò hǎo lājī fēnlèi hé huíshōu. Ránhòu, duō zhòng shù, duō zhòng huā. Zuìhòu, xiàng tóngxuémen xuānchuán huánbǎo zhīshi. Yīncǐ, xuéxiào huì biàn de gèng měilì.}\\
« Nous devons protéger l'environnement de l'école. D'abord, il ne faut pas jeter les déchets n'importe où. Ensuite, il faut bien faire le tri et le recyclage des déchets. Puis, plantons plus d'arbres et plus de fleurs. Enfin, faisons la promotion des connaissances écologiques auprès des camarades. Ainsi, l'école deviendra plus belle. »
\end{textebox}

Avant toute explication, observons. Relisez le texte et répondez mentalement à trois questions :

\begin{enumerate}
\item Combien y a-t-il de phrases ? (Réponse : six.)
\item Quels mots reviennent au début des phrases ? (Réponse : \zh{首先}, \zh{其次}, \zh{然后}, \zh{最后}, \zh{因此}.)
\item La première phrase annonce-t-elle une action précise ou une idée générale ? (Réponse : une idée générale, une \emph{thèse} : il faut protéger l'environnement.)
\end{enumerate}

Ces observations nous conduisent à la notion centrale de cette leçon : le \cle{connecteur logique}.

\begin{definition}[Connecteur logique]
Un \cle{connecteur logique} est un mot ou une expression qui relie les idées entre elles et montre la progression du raisonnement. En chinois comme en français, il se place généralement \textbf{en début de phrase}, suivi d'une virgule chinoise \zh{，}. Exemples : \zh{首先} (d'abord), \zh{其次} (ensuite, de plus), \zh{然后} (puis), \zh{最后} (enfin), \zh{因此} (c'est pourquoi, donc).
\end{definition}

\courssub{Analyse phrase par phrase}

Notre texte modèle, bien que court, possède une architecture complète en trois temps : \textbf{introduction (thèse)}, \textbf{actions concrètes (arguments)}, \textbf{conclusion}.

\textbf{1. L'introduction : la thèse.}\par
\zh{我们应该保护学校的环境。} \emph{Wǒmen yīnggāi bǎohù xuéxiào de huánjìng.} « Nous devons protéger l'environnement de l'école. »

Cette phrase unique joue le rôle d'introduction. Elle pose le sujet (\zh{我们} nous), la modalité (\zh{应该} devoir, il faut) et le thème (\zh{保护学校的环境} protéger l'environnement de l'école). Remarquez l'ordre des mots, identique au français : Sujet + modal + verbe + complément. Le petit mot \zh{的} relie \zh{学校} (école) à \zh{环境} (environnement) : « l'environnement \emph{de} l'école ».

\textbf{2. Le développement : quatre actions concrètes.}\par
Chaque action est introduite par un connecteur d'énumération :

\begin{itemize}
\item \zh{首先，不要乱扔垃圾。} « D'abord, ne pas jeter les déchets n'importe où. » — \zh{不要} + verbe exprime l'interdiction (« ne pas… ») ; \zh{乱扔} = jeter (\zh{扔}) de façon désordonnée (\zh{乱}).
\item \zh{其次，做好垃圾分类和回收。} « Ensuite, bien faire le tri et le recyclage. » — \zh{做好} signifie « bien faire, faire correctement » (\zh{好} est ici un complément de résultat : faire \emph{jusqu'au bon résultat}) ; \zh{和} coordonne \zh{分类} (tri) et \zh{回收} (recyclage).
\item \zh{然后，多种树、多种花。} « Puis, planter plus d'arbres et de fleurs. » — \zh{多} placé avant le verbe \zh{种} signifie « davantage, plus » ; la virgule énumérative chinoise \zh{、} sépare les deux actions parallèles.
\item \zh{最后，向同学们宣传环保知识。} « Enfin, sensibiliser les camarades. » — la préposition \zh{向} (vers, auprès de) introduit le destinataire : \zh{向同学们} « auprès des camarades » ; ce complément se place \textbf{avant le verbe} \zh{宣传}.
\end{itemize}

\textbf{3. La conclusion : la conséquence.}\par
\zh{因此，学校会变得更美丽。} \emph{Yīncǐ, xuéxiào huì biàn de gèng měilì.} « Ainsi, l'école deviendra plus belle. »

Le connecteur \zh{因此} (donc, par conséquent) introduit le résultat logique des actions précédentes. Le verbe modal \zh{会} exprime ici la probabilité future (« deviendra ») ; \zh{变得} = devenir ; \zh{更} = encore plus.

\begin{propriete}[Structure type : thèse – arguments – conclusion]
Un texte d'expression écrite simple en chinois suit la structure \cle{thèse – arguments – conclusion}, qui garantit la cohérence du propos :
\begin{enumerate}
\item \textbf{Thèse} (une phrase) : on annonce l'idée générale, souvent avec un modal (\zh{应该}, \zh{要}).
\item \textbf{Arguments / actions} (trois à quatre phrases) : chaque phrase est introduite par un connecteur d'énumération (\zh{首先} → \zh{其次} → \zh{然后} → \zh{最后}).
\item \textbf{Conclusion} (une phrase) : on tire la conséquence avec \zh{因此} ou \zh{所以}, souvent avec \zh{会} + verbe pour exprimer le résultat futur.
\end{enumerate}
Exception : dans un texte très court (affiche, slogan), on peut supprimer la conclusion et finir sur \zh{最后}.
\end{propriete}

\begin{popfigure}
\begin{tikzpicture}[node distance=0.55cm]
\node[draw=popblue, fill=popblueL, rounded corners, minimum width=4.6cm, minimum height=1.1cm, align=center] (intro) {\textbf{Introduction (thèse)}\\ 我们应该保护学校的环境。};
\node[draw=popgreen, fill=popgreenL, rounded corners, minimum width=4.6cm, minimum height=2.6cm, align=center, below=of intro] (args) {\textbf{Arguments (actions)}\\ 首先，不要乱扔垃圾。\\ 其次，做好垃圾分类和回收。\\ 然后，多种树、多种花。\\ 最后，宣传环保知识。};
\node[draw=poporange, fill=poporangeL, rounded corners, minimum width=4.6cm, minimum height=1.1cm, align=center, below=of args] (concl) {\textbf{Conclusion (conséquence)}\\ 因此，学校会变得更美丽。};
\draw[-{Stealth}, thick, popdark] (intro) -- (args);
\draw[-{Stealth}, thick, popdark] (args) -- (concl);
\end{tikzpicture}
\legende{Schéma de la structure du texte modèle : thèse → arguments → conclusion.}
\end{popfigure}

\courssub{Le vocabulaire clé du thème}

\begin{center}
\begin{tabularx}{\linewidth}{|X|X|X|X|}
\hline
\rowcolor{popblueL} Caractères & Pinyin & Nature & Traduction \\
\hline
环境 & huánjìng & nom & environnement \\
\hline
保护 & bǎohù & verbe & protéger \\
\hline
垃圾 & lājī & nom & déchets, ordures \\
\hline
分类 & fēnlèi & verbe / nom & trier ; le tri \\
\hline
回收 & huíshōu & verbe / nom & recycler ; le recyclage \\
\hline
种树 & zhòng shù & verbe + compl. & planter des arbres \\
\hline
宣传 & xuānchuán & verbe & faire la promotion, sensibiliser \\
\hline
环保 & huánbǎo & nom (abrév.) & protection de l'environnement, écologie \\
\hline
知识 & zhīshi & nom & connaissances \\
\hline
美丽 & měilì & adjectif & beau, magnifique \\
\hline
干净 & gānjìng & adjectif & propre \\
\hline
垃圾桶 & lājītǒng & nom & poubelle \\
\hline
\end{tabularx}
\end{center}

Remarquez que \zh{环保} est l'abréviation de \zh{环境保护} (\emph{huánjìng bǎohù}, protection de l'environnement) : le chinois adore ces mots-valises de deux caractères. De même, \zh{垃圾分类} (le tri des déchets) est un mot composé très fréquent dans les médias chinois actuels.

\courssub{Les connecteurs logiques en détail}

\begin{center}
\begin{tabularx}{\linewidth}{|X|X|X|}
\hline
\rowcolor{popblueL} Connecteur & Fonction & Équivalent français \\
\hline
首先 shǒuxiān & premier point & d'abord, premièrement \\
\hline
其次 qícì & point suivant & ensuite, de plus \\
\hline
然后 ránhòu & étape chronologique & puis, après cela \\
\hline
最后 zuìhòu & dernier point & enfin, pour finir \\
\hline
因此 yīncǐ & conséquence logique & donc, c'est pourquoi \\
\hline
\end{tabularx}
\end{center}

\begin{methode}[Repérer les connecteurs dans un texte]
\begin{enumerate}
\item Lisez le texte une première fois pour le sens global.
\item Surlignez ou soulignez chaque connecteur : \zh{首先}, \zh{其次}, \zh{然后}, \zh{最后}, \zh{因此}.
\item Numérotez les idées : chaque connecteur d'énumération introduit une nouvelle action.
\item Vérifiez la logique : les connecteurs \zh{首先} → \zh{最后} ordonnent les actions ; \zh{因此} annonce la conclusion. Si \zh{因此} apparaît au milieu, c'est que la phrase exprime une conséquence partielle.
\item Servez-vous de ce squelette pour rédiger votre propre texte en remplaçant les actions.
\end{enumerate}
\end{methode}

\begin{exemplebox}[Connecteurs en contexte]
\zh{首先，我们要捡垃圾。} \emph{Shǒuxiān, wǒmen yào jiǎn lājī.} « D'abord, nous devons ramasser les déchets. »\\
\zh{然后，我们把垃圾放进垃圾桶里。} \emph{Ránhòu, wǒmen bǎ lājī fàng jìn lājītǒng lǐ.} « Ensuite, nous mettons les déchets dans la poubelle. »\\
\zh{因此，环境变好了。} \emph{Yīncǐ, huánjìng biàn hǎo le.} « Donc, l'environnement s'est amélioré. »
\end{exemplebox}

\begin{attention}[Ne confondez pas 然后 et 因此]
Erreur fréquente des francophones : confondre \zh{然后} (\emph{ránhòu}, « ensuite ») et \zh{因此} (\emph{yīncǐ}, « donc »). Le premier indique la \textbf{suite chronologique} des actions ; le second exprime la \textbf{conséquence logique}. Comparez :\\
\zh{我们先扫地，然后擦桌子。} \emph{Wǒmen xiān sǎo dì, ránhòu cā zhuōzi.} « Nous balayons d'abord, \emph{puis} nous essuyons les tables. » (chronologie)\\
\zh{他每天回收垃圾，因此学校很干净。} \emph{Tā měi tiān huíshōu lājī, yīncǐ xuéxiào hěn gānjìng.} « Il recycle chaque jour, \emph{donc} l'école est propre. » (conséquence)\\
Un autre piège : ne commencez jamais un texte par \zh{因此} si aucune cause n'a été exprimée avant.
\end{attention}

\courssub{Temps et modalité dans le texte modèle}

Le chinois ne conjugue pas les verbes : ce sont des \textbf{mots-outils} qui indiquent le temps et la modalité.

\begin{itemize}
\item \textbf{Marques de temps.} Le texte modèle peut être enrichi avec \zh{现在} (\emph{xiànzài}, maintenant) et \zh{每天} (\emph{měi tiān}, chaque jour), placés \textbf{après le sujet et avant le verbe} : \zh{现在我们每天保护学校的环境。} \emph{Xiànzài wǒmen měi tiān bǎohù xuéxiào de huánjìng.} « Maintenant, nous protégeons chaque jour l'environnement de l'école. » Attention : contrairement au français, on ne met jamais le complément de temps à la fin de la phrase.
\item \textbf{Modalité.} \zh{应该} (\emph{yīnggāi}) exprime le devoir moral, le conseil (« il faudrait ») ; \zh{要} (\emph{yào}) exprime la nécessité ou la volonté ferme (« il faut »). \zh{应该} est plus doux et plus écrit ; c'est le modal idéal pour une affiche de sensibilisation. La négation se fait avec \zh{不应该} ou \zh{不要} + verbe : \zh{不要乱扔垃圾} « ne jetez pas les déchets n'importe où ».
\end{itemize}

\courssub{Écriture correcte des caractères du thème}

Pour rédiger sans faute, il faut maîtriser les traits et les clés (radicaux) des caractères du thème. Rappel des règles générales d'écriture : de haut en bas, de gauche à droite, l'horizontal avant le vertical, l'extérieur avant l'intérieur, et on ferme les cadres en dernier.

\begin{center}
\begin{tabularx}{\linewidth}{|X|X|X|X|}
\hline
\rowcolor{popblueL} Caractère & Clé (radical) & Ordre des traits (conseil) & Piège à éviter \\
\hline
环 huán & 王 (jade, à gauche) & 王 d'abord (4 traits), puis 不 à droite & ne pas écrire la clé 王 comme 玉 (pas de point) \\
\hline
境 jìng & 土 (terre, à gauche) & 土 (3 traits), puis 竟 à droite de haut en bas & ne pas confondre avec 镜 jìng (miroir, clé 钅métal) \\
\hline
保 bǎo & 亻(homme, à gauche) & 亻(2 traits), puis 呆 à droite & 亻se trace en 2 traits, pas en 1 \\
\hline
护 hù & 扌(main, à gauche) & 扌(3 traits), puis 户 à droite & ne pas confondre 户 avec 尸 \\
\hline
垃 lā & 土 (terre) & 土 puis 立 & 垃 et 圾 ont tous deux la clé 土 \\
\hline
圾 jī & 土 (terre) & 土 puis 及 (3 traits) & 及 : le 撇 part en premier \\
\hline
种 zhòng & 禾 (céréale, à gauche) & 禾 (5 traits), puis 中 & ne pas confondre avec 钟 zhōng (horloge, clé 钅) \\
\hline
树 shù & 木 (arbre, à gauche) & 木, puis 又, puis 寸 (trois parties) & caractère en trois morceaux : rester compact \\
\hline
\end{tabularx}
\end{center}

\begin{attention}[Homophones et caractères proches du thème]
\begin{itemize}
\item \zh{境} jìng (environnement, clé 土 terre) ≠ \zh{镜} jìng (miroir, clé 钅métal) : même son, sens différents ; la clé donne le sens.
\item \zh{种} zhòng (planter, verbe) se prononce aussi \emph{zhǒng} dans \zh{种子} (graine) : même caractère, deux lectures.
\item \zh{回收} huíshōu : \zh{收} (recevoir) ≠ \zh{放} fàng (poser) : ne pas inverser dans « mettre à la poubelle » (\zh{放进垃圾桶}).
\item \zh{美丽} měilì : \zh{丽} s'écrit avec un trait horizontal unique en haut, puis les deux parties symétriques ; ne pas écrire deux traits horizontaux séparés.
\end{itemize}
\end{attention}

\courssub{Comparaison avec un modèle français équivalent}

Un texte français équivalent serait : « Nous devons protéger l'environnement de notre école. D'abord, ne jetons pas les déchets par terre. Ensuite, trions et recyclons. Puis, plantons des arbres et des fleurs. Enfin, sensibilisons nos camarades. Ainsi, notre école deviendra plus belle. »

\begin{center}
\begin{tabularx}{\linewidth}{|X|X|X|}
\hline
\rowcolor{popblueL} Point de comparaison & Français & Chinois \\
\hline
Connecteurs & d'abord, ensuite, puis, enfin, ainsi & 首先、其次、然后、最后、因此 \\
\hline
Modalité & « il faut », impératif & 应该 / 要 + verbe, ou 不要 + verbe \\
\hline
Place du temps & souvent en fin de phrase & après le sujet, avant le verbe \\
\hline
Verbe & conjugué & invariable, particule 了 pour l'accompli \\
\hline
Structure & thèse – arguments – conclusion & identique : 总 – 分 – 总 \\
\hline
\end{tabularx}
\end{center}

Bonne nouvelle : la structure argumentative est la même dans les deux langues. Les Chinois appellent cette architecture \zh{总—分—总} (\emph{zǒng – fēn – zǒng}) : « général – détaillé – général ». Votre habitude française de rédiger une introduction, un développement et une conclusion est donc directement transférable.

\begin{exoresolu}[Repérage des connecteurs et de la structure]
Énoncé : dans le texte suivant, soulignez les connecteurs, puis indiquez la fonction de chaque phrase (thèse, argument, conclusion).\\
\zh{我们应该爱护教室。首先，每天扫地。其次，不要在墙上乱画。最后，保持桌子干净。因此，教室会很舒服。}
\tcblower
Solution détaillée :
\begin{enumerate}
\item Connecteurs à souligner : \zh{首先}, \zh{其次}, \zh{最后}, \zh{因此}.
\item \zh{我们应该爱护教室。} \emph{Wǒmen yīnggāi àihù jiàoshì.} → \textbf{thèse} (modal \zh{应该} + idée générale : « Nous devons prendre soin de la salle de classe »).
\item \zh{首先，每天扫地。} \emph{Shǒuxiān, měi tiān sǎo dì.} → \textbf{argument 1} (« D'abord, balayer chaque jour »).
\item \zh{其次，不要在墙上乱画。} \emph{Qícì, bú yào zài qiáng shang luàn huà.} → \textbf{argument 2} (« Ensuite, ne pas gribouiller sur les murs » ; \zh{不要} + verbe = interdiction).
\item \zh{最后，保持桌子干净。} \emph{Zuìhòu, bǎochí zhuōzi gānjìng.} → \textbf{argument 3} (« Enfin, garder les tables propres »).
\item \zh{因此，教室会很舒服。} \emph{Yīncǐ, jiàoshì huì hěn shūfu.} → \textbf{conclusion} (« Donc, la salle sera agréable » ; \zh{会} exprime le résultat futur).
\end{enumerate}
La structure est bien : thèse → trois arguments ordonnés → conséquence. Texte cohérent !
\end{exoresolu}

\begin{exoresolu}[Thème : traduction du français vers le chinois]
Énoncé : traduisez en chinois : « D'abord, nous ne devons pas jeter les déchets par terre ; ensuite, nous devons planter plus d'arbres. »
\tcblower
Solution détaillée :\\
\zh{首先，我们不应该乱扔垃圾；其次，我们应该多种树。}\\
\emph{Shǒuxiān, wǒmen bù yīnggāi luàn rēng lājī ; qícì, wǒmen yīnggāi duō zhòng shù.}\\
Étapes : (1) « d'abord » → \zh{首先} en tête de phrase, suivi de \zh{，} ; (2) « ne devons pas » → négation du modal : \zh{不应该} ; (3) « jeter les déchets n'importe où » → \zh{乱扔垃圾} ; (4) « ensuite » → \zh{其次} ; (5) « planter plus d'arbres » → \zh{多种树}, avec \zh{多} avant le verbe. Le point-virgule chinois \zh{；} relie les deux propositions.
\end{exoresolu}

\begin{exercice}[Version : traduction du chinois vers le français]
Traduisez en français :\\
\zh{我们应该保护学校的环境。首先，不要乱扔垃圾。其次，做好垃圾分类。最后，多种树。因此，学校会更美丽。}
\end{exercice}

\begin{corrige}[Version : traduction du chinois vers le français]
« Nous devons protéger l'environnement de l'école. D'abord, il ne faut pas jeter les déchets n'importe où. Ensuite, il faut bien faire le tri des déchets. Enfin, plantons plus d'arbres. Ainsi, l'école sera plus belle. »\\
Points de vigilance : \zh{做好} = « bien faire » (complément de résultat) ; \zh{更} = « encore plus » ; \zh{会} marque ici le futur probable.
\end{corrige}

\begin{exercice}[Production écrite guidée]
Rédigez en chinois un court texte (5 à 6 phrases) intitulé \zh{我们应该怎么保护教室的环境？} en suivant le squelette : thèse avec \zh{应该} → trois actions avec \zh{首先}、\zh{其次}、\zh{最后} → conclusion avec \zh{因此} et \zh{会}. Utilisez au moins quatre mots du vocabulaire de la leçon.
\end{exercice}

\begin{corrige}[Production écrite guidée (modèle de réponse)]
\zh{我们应该保护教室的环境。首先，不要在地上乱扔垃圾。其次，每天扫地、擦黑板。最后，把垃圾放进垃圾桶里。因此，教室会变得又干净又美丽。}\\
\emph{Wǒmen yīnggāi bǎohù jiàoshì de huánjìng. Shǒuxiān, bú yào zài dìshang luàn rēng lājī. Qícì, měi tiān sǎo dì, cā hēibǎn. Zuìhòu, bǎ lājī fàng jìn lājītǒng lǐ. Yīncǐ, jiàoshì huì biàn de yòu gānjìng yòu měilì.}\\
« Nous devons protéger l'environnement de la salle de classe. D'abord, ne pas jeter les déchets par terre. Ensuite, balayer et nettoyer le tableau chaque jour. Enfin, mettre les déchets dans la poubelle. Ainsi, la salle deviendra à la fois propre et belle. »
\end{corrige}

\begin{savaistu}[La semaine de l'environnement scolaire en Chine]
En Chine, de nombreuses écoles organisent chaque année au printemps une « Semaine de l'environnement scolaire » (\zh{学校环保周}, \emph{xuéxiào huánbǎo zhōu}) : affiches, concours de slogans, plantations d'arbres et tri des déchets. Depuis 2019, la ville de Shanghai a rendu le tri des déchets (\zh{垃圾分类}) obligatoire pour tous les habitants, et les écoliers apprennent très tôt à distinguer déchets recyclables, déchets de cuisine et déchets dangereux. Au Cameroun, des initiatives similaires existent, comme les journées de nettoyage « keep your school clean » dans plusieurs lycées de Yaoundé et de Douala : un beau terrain de rencontre entre les deux cultures !
\end{savaistu}

\begin{aretenir}[L'essentiel de la section]
\begin{itemize}
\item Un texte simple chinois suit la structure \cle{thèse – arguments – conclusion} (\zh{总—分—总}).
\item Les connecteurs \zh{首先}、\zh{其次}、\zh{然后}、\zh{最后} ordonnent les actions ; \zh{因此} introduit la conséquence.
\item \zh{应该} + verbe exprime le devoir-conseil ; \zh{不要} + verbe exprime l'interdiction ; la négation du devoir est \zh{不应该}.
\item Le complément de temps (\zh{现在}, \zh{每天}) se place après le sujet, avant le verbe ; le complément introduit par \zh{向} se place avant le verbe.
\item Vocabulaire clé : \zh{环境}、\zh{保护}、\zh{垃圾}、\zh{分类}、\zh{回收}、\zh{种树}、\zh{宣传}、\zh{干净}.
\item Écriture : attention aux clés (\zh{境} clé 土 ≠ \zh{镜} clé 钅) et à la double lecture de \zh{种} (zhòng / zhǒng).
\end{itemize}
\end{aretenir}

\coursec{Structure et connecteurs logiques}

Dans la section précédente, nous avons lu et commenté un texte modèle : \zh{你应该怎么保护学校的环境？} Nous avons remarqué que ce texte n'est pas une simple suite de phrases juxtaposées : les idées s'enchaînent grâce à de petits mots-outils comme \zh{首先}, \zh{其次}, \zh{最后}, \zh{因为……所以……}. Ces mots, appelés \cle{connecteurs logiques}, sont la charpente invisible de tout texte bien écrit. Sans eux, votre production écrite ressemblera à une liste télégraphique ; avec eux, elle deviendra un paragraphe fluide et argumenté, exactement ce que l'examinateur attend. Cette section vous apprend à les choisir, à les placer correctement dans la phrase et à les combiner.

\courssub{Observation : repérer les connecteurs dans un texte}

Lisons d'abord un court paragraphe et soulignons mentalement les mots qui relient les phrases entre elles.

\begin{textebox}[保护学校的环境]
我们应该保护学校的环境。首先，我们不要乱扔垃圾，要把垃圾放进垃圾桶里。其次，我们要爱护花草树木，不能摘花，也不能踩草坪。因为干净的校园让老师和同学都很舒服，所以每个人都应该参加大扫除。最后，我们可以一起种树，让学校更美丽。总之，保护环境就是保护我们自己。\\
Wǒmen yīnggāi bǎohù xuéxiào de huánjìng. Shǒuxiān, wǒmen bú yào luàn rēng lājī, yào bǎ lājī fàng jìn lājītǒng lǐ. Qícì, wǒmen yào àihù huācǎo shùmù, bù néng zhāi huā, yě bù néng cǎi cǎopíng. Yīnwèi gānjìng de xiàoyuán ràng lǎoshī hé tóngxué dōu hěn shūfu, suǒyǐ měi ge rén dōu yīnggāi cānjiā dàsǎochú. Zuìhòu, wǒmen kěyǐ yìqǐ zhòng shù, ràng xuéxiào gèng měilì. Zǒngzhī, bǎohù huánjìng jiùshì bǎohù wǒmen zìjǐ.\\
« Nous devons protéger l'environnement de l'école. D'abord, nous ne devons pas jeter les déchets n'importe où, il faut les mettre dans la poubelle. Ensuite, nous devons prendre soin des fleurs, des plantes et des arbres : on ne peut ni cueillir les fleurs, ni marcher sur la pelouse. Parce qu'un campus propre met les professeurs et les élèves à l'aise, chacun doit participer au grand nettoyage. Enfin, nous pouvons planter des arbres ensemble pour rendre l'école plus belle. En somme, protéger l'environnement, c'est nous protéger nous-mêmes. »
\end{textebox}

\textbf{Questions d'observation :}
\begin{enumerate}
\item Quels mots introduisent le premier, le deuxième et le dernier argument ?
\item Quel couple de mots exprime la cause et la conséquence ?
\item Quel mot annonce la conclusion générale ?
\end{enumerate}

Réponses : \zh{首先} (d'abord), \zh{其次} (ensuite), \zh{最后} (enfin) organisent les arguments ; le couple \zh{因为……所以……} exprime cause et conséquence ; \zh{总之} introduit la conclusion.

\courssub{Les connecteurs de séquence : ordonner les idées}

\begin{definition}[Connecteur de séquence]
Un \cle{connecteur de séquence} est un mot-outil qui organise les idées dans l'ordre chronologique ou logique d'un texte. En chinois, la série écrite la plus fréquente est : \zh{首先} (d'abord) → \zh{其次} (ensuite, deuxièmement) → \zh{再次} (puis, une fois de plus) → \zh{然后} (puis, après) → \zh{最后} (enfin).
\end{definition}

\begin{propriete}[Position des connecteurs de séquence]
Le connecteur de séquence se place \textbf{en début de phrase ou de proposition}, suivi d'une virgule chinoise \zh{，}, puis viennent le sujet et le verbe :\\
\cle{Connecteur ，+ Sujet + Verbe + Complément}.\\
Exemple : \zh{首先，我们不要乱扔垃圾。} \emph{Shǒuxiān, wǒmen bú yào luàn rēng lājī.} « D'abord, nous ne devons pas jeter les déchets n'importe où. »
\end{propriete}

\begin{exemplebox}[La série de séquence en action]
\zh{首先，我们打扫教室。} \emph{Shǒuxiān, wǒmen dǎsǎo jiàoshì.} « D'abord, nous nettoyons la salle de classe. »\\
\zh{其次，我们分类垃圾。} \emph{Qícì, wǒmen fēnlèi lājī.} « Ensuite, nous trions les déchets. »\\
\zh{然后，我们给花浇水。} \emph{Ránhòu, wǒmen gěi huā jiāo shuǐ.} « Puis, nous arrosons les fleurs. »\\
\zh{最后，大家一起种树。} \emph{Zuìhòu, dàjiā yìqǐ zhòng shù.} « Enfin, tout le monde plante des arbres ensemble. »
\end{exemplebox}

\begin{attention}[« puis » ne se traduit pas toujours par 然后]
Erreur fréquente des francophones : traduire systématiquement « puis » par \zh{然后}. Or \zh{然后} marque surtout la succession \emph{chronologique} d'actions (souvent à l'oral) : \zh{我先吃饭，然后做作业。} \emph{Wǒ xiān chīfàn, ránhòu zuò zuòyè.} « Je mange d'abord, puis je fais mes devoirs. » Dans un texte argumentatif écrit, pour énumérer des arguments, on préfère \zh{其次} (deuxièmement), plus formel. De même, n'inversez pas \zh{然后} et \zh{最后} : \zh{最后} clôt obligatoirement la série ; rien ne vient après lui.
\end{attention}

\courssub{Addition, opposition, cause et conclusion}

\begin{definition}[Autres familles de connecteurs]
Outre la séquence, quatre familles sont indispensables : l'\cle{addition} (\zh{而且} « de plus », \zh{也} « aussi »), l'\cle{opposition} (\zh{但是} « mais »), la \cle{cause-conséquence} (\zh{因为……所以……} « parce que… donc… », \zh{因此} « c'est pourquoi ») et la \cle{conclusion} (\zh{总之} « en somme »).
\end{definition}

\begin{propriete}[因为……所以…… : un couple inséparable]
\zh{因为……所以……} forme un \textbf{couple inséparable} à l'écrit : \cle{因为 + cause ，所以 + conséquence}. L'omission de \zh{所以} rend la phrase incomplète ou maladroite en chinois standard écrit.\\
Exemple : \zh{因为我们爱护环境，所以我们分类垃圾。} \emph{Yīnwèi wǒmen àihù huánjìng, suǒyǐ wǒmen fēnlèi lājī.} « Parce que nous protégeons l'environnement, nous trions les déchets. »\\
Remarque : contrairement au français (« parce que » et « donc » ne se combinent pas dans la même phrase), le chinois \emph{exige} souvent les deux. À l'oral rapide, on peut n'en garder qu'un, mais à l'examen, écrivez le couple complet. On peut aussi remplacer \zh{所以} par \zh{因此}, placé en tête de la proposition de conséquence : \zh{垃圾很多，因此我们必须分类。} \emph{Lājī hěn duō, yīncǐ wǒmen bìxū fēnlèi.} « Il y a beaucoup de déchets, c'est pourquoi nous devons les trier. »
\end{propriete}

\begin{exemplebox}[Addition, opposition, conclusion]
\zh{种树可以美化校园，而且可以净化空气。} \emph{Zhòng shù kěyǐ měihuà xiàoyuán, érqiě kěyǐ jìnghuà kōngqì.} « Planter des arbres embellit le campus et, de plus, purifie l'air. »\\
\zh{保护环境很难，但是我们必须努力。} \emph{Bǎohù huánjìng hěn nán, dànshì wǒmen bìxū nǔlì.} « Protéger l'environnement est difficile, mais nous devons faire des efforts. »\\
\zh{总之，保护环境是每个人的责任。} \emph{Zǒngzhī, bǎohù huánjìng shì měi ge rén de zérèn.} « En somme, protéger l'environnement est la responsabilité de chacun. »
\end{exemplebox}

\begin{savaistu}[而且 : plutôt à l'oral]
Le connecteur \zh{而且} (de plus) s'emploie surtout à l'oral ou dans un style courant. Dans une rédaction soignée, les Chinois lui préfèrent \zh{此外} (cǐwài, « en outre ») ou \zh{另外} (lìngwài, « par ailleurs »). Exemple : \zh{此外，我们还要节约用水。} \emph{Cǐwài, wǒmen hái yào jiéyuē yòngshuǐ.} « En outre, nous devons aussi économiser l'eau. » Pensez-y pour enrichir votre style à l'écrit.
\end{savaistu}

\courssub{Tableau récapitulatif des connecteurs}

\begin{center}
\begin{tabularx}{\linewidth}{|X|X|X|X|}\hline
\rowcolor{popblueL} Caractères & Pinyin & Fonction & Traduction \\ \hline
首先 & shǒuxiān & séquence & d'abord, premièrement \\ \hline
其次 & qícì & séquence (écrit) & ensuite, deuxièmement \\ \hline
再次 & zàicì & séquence & puis, une fois de plus \\ \hline
然后 & ránhòu & séquence (oral) & puis, après \\ \hline
最后 & zuìhòu & séquence & enfin, pour finir \\ \hline
而且 & érqiě & addition & de plus, et aussi \\ \hline
但是 & dànshì & opposition & mais, cependant \\ \hline
因为……所以…… & yīnwèi… suǒyǐ… & cause-conséquence & parce que… donc… \\ \hline
因此 & yīncǐ & conséquence & c'est pourquoi, donc \\ \hline
总之 & zǒngzhī & conclusion & en somme, bref \\ \hline
\end{tabularx}
\end{center}

\begin{popfigure}
\begin{tikzpicture}[
  node distance=0.55cm and 0.9cm,
  root/.style={rectangle, rounded corners, draw=popdark, fill=popblue!30, font=\bfseries, align=center, minimum height=0.9cm},
  fam/.style={rectangle, rounded corners, draw=popdark, align=center, minimum height=0.8cm, font=\small},
  ex/.style={rectangle, draw=popline, fill=popcream, rounded corners, align=center, font=\small, minimum height=0.8cm}
]
\node[root, align=center] (r){Connecteurs logiques\\连接词};
\node[fam, fill=popgreenL, below left=0.8cm and 2.6cm of r] (f1) {Séquence};
\node[fam, fill=poporangeL, below=0.8cm of r, align=center] (f2){Addition /\\Opposition};
\node[fam, fill=poppurpleL, below right=0.8cm and 2.6cm of r, align=center] (f3){Cause /\\Conclusion};
\draw[-{Stealth}, thick, popdark] (r) -- (f1);
\draw[-{Stealth}, thick, popdark] (r) -- (f2);
\draw[-{Stealth}, thick, popdark] (r) -- (f3);
\node[ex, below=0.4cm of f1, align=center] (e1){首先 其次\\再次 然后 最后};
\node[ex, below=0.4cm of f2, align=center] (e2){而且 也\\但是};
\node[ex, below=0.4cm of f3, align=center] (e3){因为…所以…\\因此 总之};
\draw[-, popline] (f1) -- (e1);
\draw[-, popline] (f2) -- (e2);
\draw[-, popline] (f3) -- (e3);
\end{tikzpicture}
\legende{Organigramme des connecteurs classés par fonction : séquence, addition/opposition, cause-conclusion.}
\end{popfigure}

\courssub{Méthode : construire un paragraphe en quatre étapes}

\begin{methode}[Du plan au paragraphe réussi]
Pour rédiger un paragraphe argumentatif sur la protection de l'environnement, suivez quatre étapes :
\begin{enumerate}
\item \textbf{Thèse} : annoncez le sujet avec \zh{我们应该……} « Nous devons… ». Exemple : \zh{我们应该保护学校的环境。}
\item \textbf{Argument 1} : introduisez-le par \zh{首先}. Exemple : \zh{首先，我们不要乱扔垃圾。}
\item \textbf{Argument 2} : introduisez-le par \zh{其次} ; renforcez avec \zh{因为……所以……} si besoin. Exemple : \zh{其次，因为树木能净化空气，所以我们要爱护树木。}
\item \textbf{Conclusion} : terminez par \zh{总之}. Exemple : \zh{总之，保护环境就是保护我们自己。}
\end{enumerate}
Vérifiez ensuite : chaque connecteur est-il en tête de proposition ? Le couple \zh{因为……所以……} est-il complet ?
\end{methode}

\begin{exoresolu}[Remettre un paragraphe dans l'ordre]
Énoncé : remettez ces quatre phrases dans l'ordre logique pour reconstruire un paragraphe cohérent.\\
(a) \zh{总之，干净的学校让大家都高兴。} (b) \zh{其次，我们要节约用水用电。} (c) \zh{我们应该保护学校的环境。} (d) \zh{首先，我们不要乱扔垃圾。}
\tcblower
Solution détaillée : on cherche d'abord la thèse, qui ne contient aucun connecteur de séquence : c'est (c). Ensuite, les connecteurs indiquent l'ordre : \zh{首先} « d'abord » = (d), puis \zh{其次} « ensuite » = (b). Enfin, \zh{总之} « en somme » marque la conclusion : (a). Ordre correct : \textbf{c – d – b – a}.\\
Paragraphe reconstruit : \zh{我们应该保护学校的环境。首先，我们不要乱扔垃圾。其次，我们要节约用水用电。总之，干净的学校让大家都高兴。} « Nous devons protéger l'environnement de l'école. D'abord, nous ne devons pas jeter les déchets n'importe où. Ensuite, nous devons économiser l'eau et l'électricité. En somme, une école propre rend tout le monde heureux. »
\end{exoresolu}

\begin{exoresolu}[Corriger les erreurs de connecteurs]
Énoncé : les deux phrases suivantes contiennent chacune une erreur typique de francophone. Corrigez-les.\\
(1) \zh{因为校园很脏，我们要大扫除。} (2) \zh{最后，我们扫地。然后，我们种树。}
\tcblower
Solution détaillée :\\
(1) Le couple \zh{因为……所以……} est incomplet : il manque \zh{所以}. Correction : \zh{因为校园很脏，所以我们要大扫除。} \emph{Yīnwèi xiàoyuán hěn zāng, suǒyǐ wǒmen yào dàsǎochú.} « Parce que le campus est sale, nous devons faire le grand nettoyage. »\\
(2) Inversion de \zh{然后} et \zh{最后} : \zh{最后} clôt la série, il doit venir en dernier. Correction : \zh{首先，我们扫地。然后，我们种树。} \emph{Shǒuxiān, wǒmen sǎo dì. Ránhòu, wǒmen zhòng shù.} ou mieux : \zh{然后，我们扫地。最后，我们种树。} \emph{Ránhòu, wǒmen sǎo dì. Zuìhòu, wǒmen zhòng shù.} « Puis nous balayons le sol ; enfin, nous plantons des arbres. »
\end{exoresolu}

\begin{attention}[Les trois pièges de l'examen]
\begin{itemize}
\item \textbf{Omission de \zh{所以}} : après une proposition en \zh{因为}, vérifiez toujours que \zh{所以} (ou \zh{因此}) introduit la conséquence.
\item \textbf{Inversion \zh{然后} / \zh{最后}} : \zh{然后} peut se répéter plusieurs fois ; \zh{最后} n'apparaît qu'une seule fois, à la fin.
\item \textbf{Traduction mot à mot du français} : en français on dit « parce que… » sans « donc » ; en chinois écrit, gardez les deux moitiés du couple.
\end{itemize}
\end{attention}

\courssub{À vous de produire}

\begin{exercice}[Production personnelle guidée]
Rédigez \textbf{cinq phrases} sur la protection de l'environnement de votre école (à Yaoundé, Douala, Buea…) en utilisant \textbf{au moins trois connecteurs différents} parmi : \zh{首先}, \zh{其次}, \zh{然后}, \zh{最后}, \zh{而且}, \zh{但是}, \zh{因为……所以……}, \zh{因此}, \zh{总之}. Soulignez chaque connecteur utilisé. Modèle de départ : \zh{首先，我们……}
\end{exercice}

\begin{exercice}[Thème : traduisez en chinois]
(1) « Parce que nous aimons notre école, nous devons la garder propre. » (2) « D'abord nous trions les déchets, enfin nous plantons des arbres ensemble. »
\end{exercice}

\begin{corrige}[Exercice de thème]
(1) \zh{因为我们爱我们的学校，所以我们要保持学校干净。} \emph{Yīnwèi wǒmen ài wǒmen de xuéxiào, suǒyǐ wǒmen yào bǎochí xuéxiào gānjìng.} — Couple \zh{因为……所以……} complet.\\
(2) \zh{首先，我们分类垃圾；最后，我们一起种树。} \emph{Shǒuxiān, wǒmen fēnlèi lājī; zuìhòu, wǒmen yìqǐ zhòng shù.} — \zh{首先} ouvre la série, \zh{最后} la clôt.
\end{corrige}

\begin{aretenir}[L'essentiel de la section]
\begin{itemize}
\item Les connecteurs de séquence (\zh{首先}, \zh{其次}, \zh{再次}, \zh{然后}, \zh{最后}) se placent en tête de proposition, suivis de \zh{，}.
\item \zh{因为……所以……} est un couple inséparable à l'écrit ; \zh{因此} peut remplacer \zh{所以}.
\item \zh{但是} marque l'opposition, \zh{而且} l'addition (plutôt oral ; à l'écrit, préférez \zh{此外}), \zh{总之} la conclusion.
\item Un bon paragraphe = thèse + argument 1 (\zh{首先}) + argument 2 (\zh{其次}) + conclusion (\zh{总之}).
\end{itemize}
\end{aretenir}

\coursec{Écriture correcte des caractères du thème}

Dans la section précédente, nous avons appris à organiser notre texte sur la protection de l'environnement grâce aux connecteurs logiques : \zh{首先} (d'abord), \zh{然后} (ensuite), \zh{最后} (enfin). Mais un texte bien structuré ne suffit pas : il faut aussi \cle{écrire correctement les caractères}. Un caractère mal tracé peut devenir un autre mot, et votre beau texte perd tout son sens. Cette section vous donne les clés pour écrire sans faute les huit caractères essentiels du thème \zh{保护环境} (protéger l'environnement).

\courssub{Les huit caractères cibles du thème}

\begin{definition}[Le radical (部首 bùshǒu)]
Le \cle{radical} (ou clé) est l'élément classificatoire d'un caractère chinois. Il indique souvent le \textbf{domaine de sens} du caractère. Par exemple, les caractères contenant \zh{氵} (l'eau) parlent souvent de liquides : \zh{河} (rivière), \zh{海} (mer). Le reste du caractère, appelé \cle{composant phonétique}, donne souvent une indication sur la \textbf{prononciation}. Connaître le radical permet de deviner le sens et de chercher le caractère dans un dictionnaire.
\end{definition}

\begin{center}
\begin{tabularx}{\linewidth}{|X|X|X|X|}
\hline
\rowcolor{popblueL} Caractère & Pinyin & Radical (sens) & Composant phonétique / structurel \\
\hline
环 & huán & 王 (jade) & 睘 (jiǒng/quán) \rightarrow simplifié en 不 \\
\hline
境 & jìng & 土 (terre) & 竟 jìng (donne le son) \\
\hline
保 & bǎo & 亻 (personne) & 呆 dāi (partie droite, indicateur phonétique ancien) \\
\hline
护 & hù & 扌 (main) & 户 hù (donne exactement le son) \\
\hline
学 & xué & 子 (enfant) & (trait brisé) (partie supérieure : 宀 + 爻/㐆) \\
\hline
校 & xiào & 木 (bois) & 交 jiāo (donne le son) \\
\hline
垃 & lā & 土 (terre) & 立 lì (proche du son) \\
\hline
圾 & jī & 土 (terre) & 及 jí (donne presque le son) \\
\hline
\end{tabularx}
\end{center}

\begin{exemplebox}[Des mots construits par association de sens]
\zh{环境} \emph{huánjìng} « environnement » = \zh{环} (entourer) + \zh{境} (frontière, territoire) : littéralement « ce qui nous entoure ».\\
\zh{保护} \emph{bǎohù} « protéger » = \zh{保} (garder) + \zh{护} (protéger) : deux synonymes qui se renforcent.\\
\zh{学校} \emph{xuéxiào} « école » = \zh{学} (étudier) + \zh{校} (établissement scolaire).\\
\zh{垃圾} \emph{lājī} « déchets, ordures » : les deux caractères portent le radical \zh{土} (terre).
\end{exemplebox}

\begin{savaistu}[Pourquoi 垃圾 a-t-il le radical de la terre ?]
Le caractère \zh{垃} (déchet) partage le radical \zh{土} (terre) avec \zh{圾}. Pourquoi ? Parce que, traditionnellement, les déchets retournent à la terre : on les brûlait ou on les enterrait. Au Cameroun comme en Chine, la question des ordures dans les cours d'école et les marchés est un vrai sujet de société : c'est exactement le thème de votre production écrite \zh{你应该怎么保护学校的环境} (Comment dois-tu protéger l'environnement de l'école ?).
\end{savaistu}

\courssub{Les règles générales de l'ordre des traits}

\begin{propriete}[L'ordre des traits (笔顺 bǐshùn)]
L'écriture d'un caractère suit des règles strictes :
\begin{enumerate}
\item \textbf{Horizontal avant vertical} : dans \zh{十} (dix), on trace d'abord 一, puis 丨.
\item \textbf{Gauche avant droite} : dans \zh{保}, on écrit d'abord le radical 亻 à gauche, puis la partie droite.
\item \textbf{Haut avant bas} : dans \zh{学}, on écrit d'abord le haut, puis 子 en bas.
\item \textbf{Centre avant les ailes} : dans \zh{小}, le trait central 亅 précède les deux points latéraux.
\item \textbf{Extérieur avant intérieur} (et on ferme en dernier) : dans \zh{国}, le cadre se ferme par le trait du bas à la fin.
\end{enumerate}
Respecter l'ordre des traits n'est pas un caprice : il garantit la \textbf{proportion}, la \textbf{vitesse} et la \textbf{lisibilité} de l'écriture.
\end{propriete}

\begin{attention}[Erreur fréquente des francophones]
Beaucoup d'élèves écrivent le trait vertical central de \zh{保} avant le trait horizontal supérieur de la partie droite, comme on le ferait en dessinant. Résultat : le caractère se déforme et devient illisible. Respecter l'ordre 1 à 9 (voir ci-dessous) évite systématiquement cette déformation. Autre piège : en français, nos lettres s'écrivent de gauche à droite en une ligne ; en chinois, chaque caractère occupe un \textbf{carré imaginaire} et se construit en couches, du haut vers le bas.
\end{attention}

\courssub{Ordre des traits détaillé : 环, 保, 护}

\textbf{环 huán (8 traits)}\par
Radical 王 à gauche (4 traits), partie droite 不 (4 traits). Étapes : 1) premier horizontal de 王 ; 2) deuxième horizontal ; 3) vertical ; 4) troisième horizontal (plus long, légèrement montant) ; 5) horizontal de 不 ; 6) trait tombant à gauche (撇) ; 7) vertical ; 8) point à droite (丶). Le radical 王 est \textbf{étroit} (environ un tiers de la largeur), la partie droite est \textbf{large}.

\begin{popfigure}
\begin{tikzpicture}[scale=0.9]
\draw[step=1, popline, very thin] (0,0) grid (3,3);
\draw[popmuted, dashed] (1.5,0) -- (1.5,3);
\draw[popmuted, dashed] (0,1.5) -- (3,1.5);
\draw[very thick, popink, -{Stealth}] (0.25,2.35) -- (1.15,2.35) node[midway, above, font=\small, popblue] {1};
\draw[very thick, popink, -{Stealth}] (0.25,1.75) -- (1.05,1.75) node[midway, above, font=\small, popblue] {2};
\draw[very thick, popink, -{Stealth}] (0.7,2.35) -- (0.7,0.95) node[midway, left, font=\small, popblue] {3};
\draw[very thick, popink, -{Stealth}] (0.2,0.95) -- (1.2,0.75) node[midway, below, font=\small, popblue] {4};
\draw[very thick, poporange, -{Stealth}] (1.45,2.5) -- (2.8,2.5) node[midway, above, font=\small, poporange] {5};
\draw[very thick, poporange, -{Stealth}] (2.15,2.35) -- (1.55,1.35) node[midway, left, font=\small, poporange] {6};
\draw[very thick, poporange, -{Stealth}] (2.15,2.3) -- (2.15,0.7) node[midway, right, font=\small, poporange] {7};
\draw[very thick, poporange, -{Stealth}] (2.35,1.3) -- (2.75,0.85) node[midway, right, font=\small, poporange] {8};
\node[font=\small, popdark] at (1.5,-0.5) {环 : traits 1--4 = 王 (gauche, étroit), traits 5--8 = 不 (droite, large)};
\end{tikzpicture}
\legende{Ordre des traits de 环 (huán) : d'abord le radical 王 à gauche, puis 不 à droite.}
\end{popfigure}

\textbf{保 bǎo (9 traits)}\par
Radical 亻 à gauche (2 traits), partie droite 呆 (7 traits). Étapes : 1) trait tombant à gauche (撇) ; 2) vertical ; puis à droite : 3) vertical du 口 ; 4) horizontal-vertical du 口 (un seul trait) ; 5) horizontal de fermeture du 口 ; 6) horizontal du 木 ; 7) vertical du 木 ; 8) trait tombant à gauche ; 9) trait tombant à droite (捺). Règle appliquée : gauche avant droite, haut (口) avant bas (木), horizontal avant vertical.

\begin{popfigure}
\begin{tikzpicture}[scale=0.9]
\draw[step=1, popline, very thin] (0,0) grid (3,3);
\draw[popmuted, dashed] (1.5,0) -- (1.5,3);
\draw[popmuted, dashed] (0,1.5) -- (3,1.5);
\draw[very thick, popink, -{Stealth}] (0.85,2.7) -- (0.35,1.7) node[midway, left, font=\small, popblue] {1};
\draw[very thick, popink, -{Stealth}] (0.6,2.0) -- (0.6,0.5) node[midway, left, font=\small, popblue] {2};
\draw[very thick, poporange, -{Stealth}] (1.6,2.75) -- (1.6,2.15) node[midway, left, font=\small, poporange] {3};
\draw[very thick, poporange, -{Stealth}] (1.6,2.75) -- (2.55,2.75) -- (2.55,2.15) node[pos=0.3, above, font=\small, poporange] {4};
\draw[very thick, poporange, -{Stealth}] (1.6,2.15) -- (2.55,2.15) node[midway, below, font=\small, poporange] {5};
\draw[very thick, poporange, -{Stealth}] (1.35,1.55) -- (2.8,1.55) node[midway, above, font=\small, poporange] {6};
\draw[very thick, poporange, -{Stealth}] (2.05,1.9) -- (2.05,0.45) node[midway, right, font=\small, poporange] {7};
\draw[very thick, poporange, -{Stealth}] (2.0,1.2) -- (1.35,0.5) node[midway, left, font=\small, poporange] {8};
\draw[very thick, poporange, -{Stealth}] (2.1,1.2) -- (2.8,0.5) node[midway, right, font=\small, poporange] {9};
\node[font=\small, popdark] at (1.5,-0.5) {保 : traits 1--2 = 亻, traits 3--5 = 口, traits 6--9 = 木};
\end{tikzpicture}
\legende{Ordre des traits de 保 (bǎo) : la personne 亻 protège, à droite 呆 se construit de haut en bas.}
\end{popfigure}

\textbf{护 hù (7 traits)}\par
Radical 扌 (main) à gauche (3 traits), 户 à droite (4 traits). Étapes : 1) horizontal ; 2) crochet vertical (亅) ; 3) trait montant (提) ; puis à droite : 4) point ; 5) horizontal-cassé (横折) ; 6) horizontal ; 7) trait tombant à gauche (撇). Le radical 扌 est étroit ; 户 occupe les deux tiers droits du carré.

\begin{popfigure}
\begin{tikzpicture}[scale=0.9]
\draw[step=1, popline, very thin] (0,0) grid (3,3);
\draw[popmuted, dashed] (1.5,0) -- (1.5,3);
\draw[popmuted, dashed] (0,1.5) -- (3,1.5);
\draw[very thick, popink, -{Stealth}] (0.2,2.2) -- (1.15,2.2) node[midway, above, font=\small, popblue] {1};
\draw[very thick, popink, -{Stealth}] (0.7,2.7) -- (0.7,0.6) -- (0.55,0.85) node[pos=0.5, right, font=\small, popblue] {2};
\draw[very thick, popink, -{Stealth}] (0.2,1.15) -- (1.15,1.55) node[midway, below, font=\small, popblue] {3};
\draw[very thick, popgreen, -{Stealth}] (2.35,2.75) -- (2.1,2.5) node[midway, right, font=\small, popgreen] {4};
\draw[very thick, popgreen, -{Stealth}] (1.7,2.3) -- (2.7,2.3) -- (2.7,1.75) node[pos=0.3, above, font=\small, popgreen] {5};
\draw[very thick, popgreen, -{Stealth}] (1.7,1.75) -- (2.7,1.75) node[midway, below, font=\small, popgreen] {6};
\draw[very thick, popgreen, -{Stealth}] (1.75,1.7) -- (1.45,0.5) node[midway, left, font=\small, popgreen] {7};
\node[font=\small, popdark] at (1.5,-0.5) {护 : traits 1--3 = 扌 (main), traits 4--7 = 户 (porte)};
\end{tikzpicture}
\legende{Ordre des traits de 护 (hù) : la main 扌 à gauche, la porte 户 à droite — protéger, c'est tenir la porte.}
\end{popfigure}

\begin{methode}[Mémoriser 环 par décomposition]
Pour ne jamais oublier \zh{环}, décomposez-le en trois gestes :
\begin{enumerate}
\item Écrivez \zh{王} (jade, roi) à gauche, bien étroit : 4 traits.
\item À droite, commencez par le trait horizontal supérieur de \zh{不}.
\item Terminez par le crochet et le point : le « jade entouré » évoque ce qui nous \textbf{entoure} → l'\textbf{environnement}.
\end{enumerate}
Créez une petite histoire : « Le roi (王) dit non (不) à la pollution : il protège son environnement. » Les histoires mnémotechniques sont la méthode la plus efficace pour retenir les caractères composés.
\end{methode}

\courssub{Attention aux caractères proches visuellement}

\begin{center}
\begin{tabularx}{\linewidth}{|X|X|X|}
\hline
\rowcolor{popblueL} Paire à ne pas confondre & Différence visuelle & Traduction \\
\hline
环 huán / 还 hái & 环 = 王 + 不 ; 还 = 不 + 辶 (marche) & entourer / encore, retourner \\
\hline
护 hù / 把 bǎ & 护 = 扌+ 户 ; 把 = 扌+ 巴 & protéger / prendre (préposition) \\
\hline
校 xiào / 核 hé & 校 = 木 + 交 ; 核 = 木 + 亥 & école / noyau, vérifier \\
\hline
\end{tabularx}
\end{center}

\begin{attention}[Le piège du radical identique]
\zh{护} et \zh{把} partagent le radical \zh{扌} (main) : seule la partie droite change. De même, \zh{校} et \zh{核} partagent \zh{木} (bois). Quand deux caractères ont le même radical, c'est le \textbf{composant phonétique} qui fait toute la différence de sens et de son. Dans votre copie, un seul trait mal placé transforme « école » (\zh{学校}) en mot inexistant : relisez chaque caractère en vérifiant les deux parties.
\end{attention}

\courssub{Les règles de proportion dans le carré d'écriture}

\begin{propriete}[Proportion gauche-droite]
Dans les caractères à structure gauche-droite comme \zh{环}, \zh{境}, \zh{保}, \zh{护}, \zh{校}, \zh{垃}, \zh{圾} :
\begin{itemize}
\item la partie gauche (le radical) occupe environ \textbf{un tiers} de la largeur et est légèrement \textbf{compressée} ;
\item la partie droite occupe environ \textbf{deux tiers} et porte l'essentiel de l'information ;
\item les deux parties s'alignent sur le \textbf{même axe vertical imaginaire} : ni la gauche ni la droite ne doit « tomber ».
\end{itemize}
Exception notable : \zh{学} est un caractère \textbf{haut-bas} : la partie supérieure est compressée, \zh{子} en bas est large et stable, comme un socle.
\end{propriete}

\courssub{Utiliser le dictionnaire de radicaux}

\begin{methode}[Retrouver un caractère inconnu]
Vous rencontrez un caractère inconnu, par exemple \zh{境} ? Voici la démarche :
\begin{enumerate}
\item \textbf{Identifiez le radical} : ici \zh{土} (terre), à gauche.
\item \textbf{Comptez les traits du radical} : 土 a 3 traits → cherchez dans la table des radicaux à 3 traits.
\item \textbf{Comptez les traits restants} : 竟 a 11 traits → dans la section 土, cherchez les caractères à +11 traits.
\item \textbf{Lisez le numéro de page} indiqué, puis vérifiez le pinyin (jìng) et le sens.
\end{enumerate}
Cette compétence est précieuse à l'examen : elle vous rend autonome face à tout texte nouveau.
\end{methode}

\courssub{Exercices d'écriture guidée}

\begin{exoresolu}[Dictée et tracé sur quadrillage]
Énoncé : (1) Tracez \zh{护} en respectant l'ordre des 7 traits, dans un carré de 2 cm de côté. (2) Recopiez cinq fois le mot \zh{保护环境}. (3) Dictée : votre camarade dicte \zh{学校}, \zh{垃圾}, \zh{环境} ; écrivez-les puis vérifiez.
\tcblower
Solution détaillée : (1) Ordre de \zh{护} : horizontal → crochet vertical → trait montant (radical 扌, étroit à gauche), puis point → horizontal-cassé → horizontal → trait tombant (户, large à droite). (2) \zh{保护环境} \emph{bǎohù huánjìng} « protéger l'environnement » : vérifiez que 保 a bien 9 traits et que le 口 de sa partie droite est fermé avant d'écrire 木. (3) Corrigé de la dictée : \zh{学校} \emph{xuéxiào} « école », \zh{垃圾} \emph{lājī} « déchets », \zh{环境} \emph{huánjìng} « environnement ». Erreur à repérer : si votre \zh{垃} ressemble à \zh{拉} (tirer), c'est que le radical 土 a été tracé comme 扌 — comparez le nombre de traits : 土 = 3 traits (le 3e horizontal), 扌 = 3 traits (le 3e montant/tí).
\end{exoresolu}

\begin{exercice}[Entraînement personnel]
\begin{enumerate}
\item Écrivez dix fois chacun des huit caractères cibles en respectant l'ordre des traits, puis couvrez votre cahier et réécrivez-les de mémoire.
\item Classez ces caractères selon leur radical : 河, 境, 保, 海, 垃, 护, 校, 圾, 你, 林.
\item Écrivez la phrase \zh{我们应该保护学校的环境。} (Wǒmen yīnggāi bǎohù xuéxiào de huánjìng. « Nous devons protéger l'environnement de l'école. ») en soignant les proportions de chaque caractère.
\end{enumerate}
\end{exercice}

\begin{corrige}[Exercice d'entraînement]
1) Auto-vérification : comparez chaque caractère au modèle du tableau, comptez les traits (环 8, 保 9, 护 7). 2) Radical 氵(eau) : 河, 海 ; radical 土 (terre) : 境, 垃, 圾 ; radical 亻 (personne) : 保, 你 ; radical 扌 (main) : 护 ; radical 木 (bois) : 校, 林. 3) Phrase modèle : \zh{我们应该保护学校的环境。} — attention à \zh{们} (radical 亻 étroit) et à la place du point chinois 。 en bas à gauche de la case finale.
\end{corrige}

\begin{aretenir}[L'essentiel de la section]
\begin{itemize}
\item Huit caractères cibles : \zh{环境} (environnement), \zh{保护} (protéger), \zh{学校} (école), \zh{垃圾} (déchets).
\item Le radical donne le sens, le composant phonétique donne le son : c'est la clé de la mémorisation et de la recherche au dictionnaire.
\item Ordre des traits : horizontal avant vertical, gauche avant droite, haut avant bas, centre avant ailes.
\item Proportions : radical étroit à gauche (un tiers), partie droite large (deux tiers).
\item Méfiez-vous des sosies : 环/还, 护/把, 校/核.
\end{itemize}
\end{aretenir}

\coursec{Exercices de thème (français → chinois) et version (chinois → français)}

Dans la section précédente, nous avons appris à écrire correctement les caractères du thème de l'environnement (\zh{环}, \zh{境}, \zh{垃}, \zh{圾}, \zh{保}, \zh{护}…) en respectant les radicaux et l'ordre des traits. Il est temps maintenant de mettre ces caractères en action : nous allons les utiliser dans des exercices de \cle{traduction}. Traduire, c'est l'épreuve reine de l'écrit en chinois : elle vérifie à la fois votre vocabulaire, votre grammaire, vos connecteurs et la lisibilité de vos caractères.

\courssub{Thème et version : deux directions de la traduction}

\begin{definition}[Thème et version]
Le \cle{thème} est la traduction du \textbf{français vers le chinois}. La \cle{version} est la traduction du \textbf{chinois vers le français}. Dans les deux cas, il ne s'agit pas de remplacer mot à mot, mais de transporter le \textbf{sens} dans une phrase correcte et naturelle dans la langue d'arrivée.
\end{definition}

Observez d'abord ces deux exemples avant d'énoncer la méthode.

\begin{exemplebox}[Exemples traduits et commentés]
\textbf{Thème :} « Enfin, nous plantons des arbres. » → \zh{最后，我们种树。} (\emph{zuìhòu, wǒmen zhòng shù.})\\[2pt]
Le connecteur « enfin » devient \zh{最后} ; le verbe « planter (des arbres) » se dit simplement \zh{种树} ; pas d'article, pas de pluriel en chinois.\\[4pt]
\textbf{Version :} \zh{因为我们爱护环境，所以我们分类垃圾。} (\emph{yīnwèi wǒmen àihù huánjìng, suǒyǐ wǒmen fēnlèi lājī.}) → « Parce que nous protégeons l'environnement, nous trions les déchets. »\\[2pt]
En chinois, \zh{因为…所以…} s'emploient \textbf{ensemble} ; en français, on ne garde qu'un seul des deux (« parce que… » ou « …donc… »), jamais les deux.
\end{exemplebox}

\begin{propriete}[Règle d'or de la version]
La version doit restituer le \textbf{sens global}, pas le mot à mot. Il faut \textbf{adapter les connecteurs} au français naturel : \zh{因此} → « donc / c'est pourquoi », \zh{首先} → « d'abord », \zh{其次} → « ensuite », \zh{最后} → « enfin ». De même, en thème, « donc » se rendra par \zh{所以} ou \zh{因此} (plus formel, idéal à l'écrit).
\end{propriete}

\begin{methode}[Traduire par étapes]
\begin{enumerate}
\item \textbf{Identifier le verbe principal} de la phrase française et trouver son équivalent chinois (ex. « jeter » → \zh{扔}, « trier » → \zh{分类}, « protéger » → \zh{保护}).
\item \textbf{Choisir le connecteur approprié} (\zh{首先}, \zh{其次}, \zh{然后}, \zh{最后}, \zh{因为…所以…}, \zh{应该}) et le placer correctement, souvent en tête de phrase suivi d'une virgule chinoise \zh{，}.
\item \textbf{Construire la phrase} selon l'ordre chinois : Sujet + (complément de temps/lieu) + Verbe + Complément.
\item \textbf{Écrire les caractères en respectant l'ordre des traits} (haut → bas, gauche → droite) et vérifier la ponctuation pleine chasse \zh{。} \zh{，}.
\item \textbf{Relire} : sens complet ? caractères lisibles ? tons du pinyin corrects si demandé ?
\end{enumerate}
\end{methode}

\begin{attention}[Piège : « ne pas » = 不 ou 别 ?]
« Ne pas » se rend par \zh{不} ou \zh{别} selon le registre. \zh{别} est une \textbf{interdiction orale} directe (\zh{别扔垃圾！} « Ne jette pas de déchets ! »). À l'\textbf{écrit formel}, dans un texte argumentatif, on utilise \zh{不} : \zh{我们不应该乱扔垃圾。} (\emph{wǒmen bù yīnggāi luàn rēng lājī.}) « Nous ne devons pas jeter les déchets n'importe où. » Attention au changement de ton (sandhi) : devant un 4e ton (\zh{乱} \emph{luàn}), \zh{不} se prononce \emph{bú} (\zh{我们不乱扔垃圾} \emph{wǒmen \textbf{bú} luàn rēng lājī}). Retenez aussi : \zh{乱扔} = « jeter n'importe où / en désordre ».
\end{attention}

\begin{savaistu}[Le thème-version à l'examen]
Au baccalauréat, l'épreuve de thème/version représente environ 20\,\% de la note d'écrit de chinois, et la \textbf{lisibilité des caractères} est explicitement notée : un caractère mal tracé (trait manquant, radical confondu) coûte des points même si le sens est bon. D'où l'importance de la section précédente sur l'écriture des caractères !
\end{savaistu}

\courssub{Texte de référence : le modèle à imiter}

Avant de traduire, relisons un court texte modèle qui servira de « banque de phrases » pour vos traductions.

\begin{textebox}[我们应该怎么保护学校的环境？]
我们的学校很美丽，但是有时候地上有垃圾。我们应该怎么保护学校的环境呢？\\
首先，我们不应该乱扔垃圾，要把垃圾放进垃圾桶里。\\
其次，我们要做好垃圾分类：纸、塑料和旧书要分开放。\\
然后，我们应该节约用水用电。离开教室的时候，要关灯。\\
最后，我们要爱护花草树木，还可以种树。\\
因为学校是我们的第二个家，所以我们要一起保护它。保护环境，从我做起！\\[4pt]
\emph{Wǒmen de xuéxiào hěn měilì, dànshì yǒu shíhou dìshang yǒu lājī. Wǒmen yīnggāi zěnme bǎohù xuéxiào de huánjìng ne? Shǒuxiān, wǒmen bù yīnggāi luàn rēng lājī, yào bǎ lājī fàng jìn lājītǒng lǐ. Qícì, wǒmen yào zuò hǎo lājī fēnlèi: zhǐ, sùliào hé jiù shū yào fēnkāi fàng. Ránhòu, wǒmen yīnggāi jiéyuē yòng shuǐ yòng diàn. Líkāi jiàoshì de shíhou, yào guān dēng. Zuìhòu, wǒmen yào àihù huācǎo shùmù, hái kěyǐ zhòng shù. Yīnwèi xuéxiào shì wǒmen de dì èr gè jiā, suǒyǐ wǒmen yào yìqǐ bǎohù tā. Bǎohù huánjìng, cóng wǒ zuòqǐ!}\\[4pt]
« Notre école est très belle, mais parfois il y a des déchets par terre. Comment devons-nous protéger l'environnement de l'école ? D'abord, nous ne devons pas jeter les déchets n'importe où, nous devons les mettre dans la poubelle. Ensuite, nous devons bien trier les déchets : le papier, le plastique et les vieux livres doivent être mis séparément. Puis, nous devons économiser l'eau et l'électricité. En quittant la salle de classe, il faut éteindre la lumière. Enfin, nous devons prendre soin des fleurs, des plantes et des arbres, et nous pouvons aussi planter des arbres. Parce que l'école est notre deuxième maison, nous devons la protéger ensemble. Protéger l'environnement commence par moi ! »
\end{textebox}

\textbf{Glossaire rapide :} \zh{垃圾桶} (\emph{lājītǒng}) poubelle ; \zh{分类} (\emph{fēnlèi}) trier ; \zh{节约} (\emph{jiéyuē}) économiser ; \zh{关灯} (\emph{guān dēng}) éteindre la lumière ; \zh{爱护} (\emph{àihù}) prendre soin de ; \zh{从我做起} (\emph{cóng wǒ zuòqǐ}) « commencer par soi-même » (slogan célèbre en Chine).

\courssub{Tableau de correspondance français – chinois}

\begin{popfigure}
\begin{center}
\begin{tabularx}{\linewidth}{|X|X|X|X|}
\hline
\rowcolor{popblueL} Phrase française & Pinyin & Caractères & Notes \\
\hline
D'abord, nous ne jetons pas les déchets n'importe où. & shǒuxiān, wǒmen \textbf{bú} luàn rēng lājī. & 首先，我们不乱扔垃圾。 & \zh{首先} + virgule ; \zh{不乱扔} (ton \emph{bú} devant 4e ton) \\
\hline
Ensuite, nous devons bien trier les déchets. & qícì, wǒmen yào zuò hǎo lājī fēnlèi. & 其次，我们要做好垃圾分类。 & \zh{做好} = bien faire \\
\hline
Puis, nous économisons l'eau et l'électricité. & ránhòu, wǒmen jiéyuē yòng shuǐ yòng diàn. & 然后，我们节约用水用电。 & \zh{用水用电} sans répéter « économiser » \\
\hline
Enfin, nous plantons des arbres. & zuìhòu, wǒmen zhòng shù. & 最后，我们种树。 & pas de « des » en chinois \\
\hline
Parce que nous aimons notre école, nous la protégeons. & yīnwèi wǒmen ài wǒmen de xuéxiào, suǒyǐ wǒmen bǎohù tā. & 因为我们爱我们的学校，所以我们保护它。 & \zh{因为…所以…} ensemble \\
\hline
\end{tabularx}
\end{center}
\legende{Correspondances français–chinois des phrases-clés du thème de l'environnement.}
\end{popfigure}

\begin{popfigure}
\begin{center}
\begin{tikzpicture}[mindmap, grow cyclic, every node/.style={concept, align=center, font=\small},
root concept/.append style={concept color=popgreen, font=\bfseries},
level 1/.append style={level distance=3.2cm, sibling angle=90},
level 1 concept/.append style={concept color=popgreenL}]
\node[root concept, align=center]{保护环境\\ \emph{bǎohù huánjìng}}
  child { node[align=center]{Connecteurs\\ 首先 其次\\ 然后 最后} }
  child { node[align=center]{Actions\\ 扔垃圾 分类\\ 种树 关灯} }
  child { node[align=center]{Négation\\ 不应该\\ 不乱扔} }
  child { node[align=center]{Cause\\ 因为…所以…\\ 因此} };
\end{tikzpicture}
\end{center}
\legende{Carte mentale du lexique utile pour le thème « protection de l'environnement ».}
\end{popfigure}

\courssub{Série A — Thème : du français vers le chinois}

\begin{exercice}[Thème : traduisez en chinois (6 phrases)]
Consignes : respectez la structure thèse–arguments–conclusion si vous enchaînez les phrases ; utilisez \textbf{au moins 3 connecteurs} différents ; écrivez les caractères cibles (\zh{环境}, \zh{垃圾}, \zh{分类}, \zh{保护}, \zh{种}) avec soin ; ponctuation chinoise obligatoire.
\begin{enumerate}
\item D'abord, nous ne jetons pas les déchets n'importe où.
\item Ensuite, nous devons bien trier les déchets.
\item Puis, nous économisons l'eau et l'électricité.
\item En quittant la salle de classe, il faut éteindre la lumière.
\item Enfin, nous plantons des arbres.
\item Parce que l'école est notre deuxième maison, nous devons la protéger ensemble.
\end{enumerate}
\end{exercice}

\begin{exoresolu}[Thème guidé : phrase 1]
Énoncé : traduisez « D'abord, nous ne jetons pas les déchets n'importe où. »
\tcblower
Solution détaillée :
\begin{enumerate}
\item Verbe principal : « jeter » → \zh{扔} (\emph{rēng}) ; « n'importe où » → l'adverbe \zh{乱} (\emph{luàn}) placé \textbf{avant} le verbe : \zh{乱扔}.
\item Connecteur : « d'abord » → \zh{首先}, placé en tête suivi de \zh{，}.
\item Négation : registre écrit → \zh{不} (et non \zh{别}) : \zh{我们不乱扔垃圾}. Prononciation : \emph{wǒmen \textbf{bú} luàn rēng lājī} (sandhi sur \zh{不} devant \zh{乱}).
\item Phrase finale : \zh{首先，我们不乱扔垃圾。} (\emph{shǒuxiān, wǒmen bú luàn rēng lājī.})
\end{enumerate}
Vérification : sujet \zh{我们} + négation \zh{不} + adverbe \zh{乱} + verbe \zh{扔} + complément \zh{垃圾} : l'ordre est correct.
\end{exoresolu}

\begin{corrige}[Série A — Thème]
\begin{enumerate}
\item \zh{首先，我们不乱扔垃圾。} (\emph{shǒuxiān, wǒmen bú luàn rēng lājī.})
\item \zh{其次，我们要做好垃圾分类。} (\emph{qícì, wǒmen yào zuò hǎo lājī fēnlèi.})
\item \zh{然后，我们节约用水用电。} (\emph{ránhòu, wǒmen jiéyuē yòng shuǐ yòng diàn.})
\item \zh{离开教室的时候，要关灯。} (\emph{líkāi jiàoshì de shíhou, yào guān dēng.})
\item \zh{最后，我们种树。} (\emph{zuìhòu, wǒmen zhòng shù.})
\item \zh{因为学校是我们的第二个家，所以我们要一起保护它。} (\emph{yīnwèi xuéxiào shì wǒmen de dì èr gè jiā, suǒyǐ wǒmen yào yìqǐ bǎohù tā.})
\end{enumerate}
\end{corrige}

\courssub{Série B — Version : du chinois vers le français}

\begin{exercice}[Version : traduisez en français (5 phrases)]
Consignes : restituez le sens global en français naturel ; adaptez les connecteurs ; n'oubliez pas les articles et les marques du pluriel français.
\begin{enumerate}
\item \zh{其次，我们要做好垃圾分类。}
\item \zh{离开教室的时候，我们应该关灯。}
\item \zh{因为学校是我们的第二个家，所以我们要爱护它。}
\item \zh{保护环境，从我做起。}
\item \zh{我们不应该乱扔垃圾，要把垃圾放进垃圾桶里。}
\end{enumerate}
\end{exercice}

\begin{corrige}[Série B — Version]
\begin{enumerate}
\item Ensuite, nous devons bien trier les déchets. (litt. « faire bien le tri des déchets »)
\item En quittant la salle de classe, nous devons éteindre la lumière.
\item Parce que l'école est notre deuxième maison, nous devons en prendre soin. (un seul connecteur en français !)
\item Protéger l'environnement commence par moi. (slogan : « à chacun d'agir »)
\item Nous ne devons pas jeter les déchets n'importe où ; il faut les mettre dans la poubelle.
\end{enumerate}
\end{corrige}

\courssub{Correction guidée, auto-évaluation et remédiation}

\begin{methode}[Grille de correction (sur 20 points)]
Attribuez les points phrase par phrase :
\begin{itemize}
\item \textbf{Vocabulaire} (5 pts) : le bon mot est choisi (\zh{扔} pour « jeter », \zh{分类} pour « trier », pas de confusion \zh{垃圾}/\zh{拉圾}).
\item \textbf{Grammaire} (5 pts) : ordre des mots, place de \zh{不} et de \zh{乱} avant le verbe, \zh{的时候} correctement construit.
\item \textbf{Connecteurs} (4 pts) : au moins 3 connecteurs variés, bien placés, suivis de \zh{，}.
\item \textbf{Caractères} (4 pts) : traits complets, radicaux corrects, écriture lisible.
\item \textbf{Ponctuation} (2 pts) : \zh{，} et \zh{。} pleine chasse, pas de point français.
\end{itemize}
\end{methode}

\begin{exercice}[Auto-évaluation : cochez ce que vous avez réussi]
\begin{itemize}
\item[$\square$] J'ai utilisé au moins 3 connecteurs différents (\zh{首先}, \zh{其次}, \zh{然后}, \zh{最后}…).
\item[$\square$] J'ai écrit \zh{不} (et non \zh{别}) pour la négation à l'écrit.
\item[$\square$] J'ai placé \zh{因为…所以…} ensemble en chinois, un seul connecteur en français.
\item[$\square$] Mes caractères \zh{环境}, \zh{垃圾}, \zh{分类} sont complets et lisibles.
\item[$\square$] Ma ponctuation est chinoise (\zh{。} \zh{，}).
\end{itemize}
Notez ici vos deux principales difficultés : \rule{3cm}{0.4pt} et \rule{3cm}{0.4pt}.
\end{exercice}

\begin{exercice}[Remédiation : réécrivez ces phrases incorrectes à l'aide du modèle]
\begin{enumerate}
\item \zh{我们别乱扔垃圾。} (registre écrit exigé) → attendu : \zh{我们不应该乱扔垃圾。}
\item \zh{因为我们爱学校，我们保护它。} (il manque un élément) → attendu : \zh{因为我们爱学校，所以我们保护它。}
\item \zh{最后，我们种很多树们。} (erreur de pluriel) → attendu : \zh{最后，我们种很多树。} (\zh{们} ne s'emploie qu'avec les personnes).
\end{enumerate}
\end{exercice}

\begin{exercice}[Extension : production libre]
Rédigez un court paragraphe d'environ \textbf{50 caractères} sur le thème « Comment je protège l'environnement de mon école ». Utilisez au moins 3 connecteurs, une structure \zh{因为…所以…} et terminez par le slogan \zh{保护环境，从我做起！}. Vous pouvez mentionner votre ville (Yaoundé, Douala, Buea…) : \zh{我们在雅温得的学校很美丽…}
\end{exercice}

\begin{aretenir}[L'essentiel de la section]
\begin{itemize}
\item Thème = français → chinois ; version = chinois → français ; dans les deux cas, on traduit le \textbf{sens}, pas les mots.
\item Méthode en 5 étapes : verbe principal → connecteur → ordre de la phrase → caractères soignés → relecture.
\item Négation à l'écrit : \zh{不} ; \zh{别} est réservé à l'oral et aux interdictions directes. Attention au sandhi : \zh{不} → \emph{bú} devant un 4e ton.
\item \zh{因为…所以…} vont ensemble en chinois, jamais en français.
\item La lisibilité des caractères et la ponctuation chinoise comptent dans la note.
\end{itemize}
\end{aretenir}

\coursec{Grille d'évaluation et auto‑correction}

Vous savez maintenant rédiger un texte sur la protection de l'environnement de l'école et traduire des phrases dans les deux sens (thème et version). Il reste une dernière étape, souvent négligée mais décisive : \textbf{évaluer et corriger soi-même sa production}. Un bon élève n'est pas celui qui ne fait jamais d'erreurs, mais celui qui sait les repérer et les corriger avant de rendre sa copie. Cette section vous donne les outils du correcteur : la grille officielle d'évaluation, deux exemples de copies notées, une checklist d'auto-correction et des conseils concrets pour gagner des points.

\courssub{Présentation de la grille d'évaluation}

\begin{definition}[Grille d'évaluation]
Une \cle{grille d'évaluation} est un outil critérié qui permet une notation \textbf{objective et transparente} d'une production écrite. Au lieu de donner une note globale « au feeling », le correcteur évalue séparément plusieurs \cle{critères} (contenu, structure, vocabulaire, grammaire, présentation), chacun selon plusieurs \cle{niveaux} de réussite (Excellent, Bien, Passable, Insuffisant). L'élève connaît ainsi exactement ce qu'on attend de lui et où il doit progresser.
\end{definition}

Pour l'expression écrite \zh{你应该怎么保护学校的环境} (\emph{Nǐ yīnggāi zěnme bǎohù xuéxiào de huánjìng} — « Comment dois-tu protéger l'environnement de l'école »), la grille comporte \textbf{5 critères}, chacun noté sur 4 niveaux :

\begin{center}
\begin{tabularx}{\linewidth}{|X|X|X|X|X|}
\hline
\rowcolor{popblueL} \textbf{Critère} & \textbf{Excellent (4)} & \textbf{Bien (3)} & \textbf{Passable (2)} & \textbf{Insuffisant (1)} \\
\hline
1. Contenu et idées & Thèse claire, 3 arguments ou plus, exemples concrets (école, Cameroun) & Thèse présente, 2 arguments développés & Thèse vague, 1 seul argument, idées répétitives & Hors sujet ou aucune idée développée \\
\hline
2. Structure et connecteurs & Introduction, développement, conclusion ; 4 connecteurs variés (首先，然后，所以，但是…) & Plan visible ; 2 à 3 connecteurs corrects & Plan confus ; 1 seul connecteur ou connecteurs mal employés & Aucun plan, aucun connecteur, phrases juxtaposées \\
\hline
3. Vocabulaire et caractères & Lexique du thème riche et précis ; les 8 caractères cibles tous corrects & Bon lexique ; 1 à 2 fautes de caractères & Lexique pauvre ou répétitif ; 3 à 4 fautes de caractères & Vocabulaire hors sujet ; plus de 5 fautes de caractères \\
\hline
4. Grammaire et syntaxe & Ordre des mots toujours correct ; 因为…所以…, 应该， 得 maîtrisés ; zéro faute & 1 à 2 fautes légères qui ne gênent pas la compréhension & 3 à 4 fautes (place de 应该， double 因为…所以…) & Fautes systématiques empêchant la compréhension \\
\hline
5. Présentation et écriture & Traits dans le bon ordre, caractères équilibrés, copie propre et lisible & Écriture lisible, quelques traits hésitants & Écriture difficile, proportions irrégulières, ratures & Écriture illisible, caractères déformés, copie négligée \\
\hline
\end{tabularx}
\end{center}

\begin{propriete}[Calcul de la note finale]
Chaque critère pèse \cle{20 \%} de la note finale. On convertit le niveau obtenu en points : Excellent = 4, Bien = 3, Passable = 2, Insuffisant = 1. La somme des 5 critères donne la note sur 20.
\begin{itemize}
\item Exemple : Contenu 3 + Structure 4 + Vocabulaire 3 + Grammaire 4 + Présentation 3 = \textbf{17/20}.
\item Exemple : Contenu 2 + Structure 2 + Vocabulaire 3 + Grammaire 2 + Présentation 3 = \textbf{12/20}.
\end{itemize}
Aucun critère n'est « secondaire » : une belle écriture ne sauve pas un contenu vide, et de bonnes idées ne compensent pas une grammaire fautive.
\end{propriete}

\begin{attention}[Le piège de l'écriture négligée]
Erreur fréquente des francophones : oublier de vérifier l'\textbf{ordre des traits} sur les caractères nouveaux du thème (环，境，保，护，垃，圾…). Un caractère tracé dans le désordre est souvent déformé ou méconnaissable, et cela coûte des points sur le critère 5 (présentation) \emph{et} parfois sur le critère 3 (caractères). Règle d'or : \textbf{haut avant bas, gauche avant droite, horizontal avant vertical, l'extérieur avant l'intérieur, et on ferme la boîte en dernier} (comme dans 国 : on écrit la dernière horizontale de fermeture à la fin).
\end{attention}

\courssub{Exemples de textes notés avec commentaires du correcteur}

Observons deux copies d'élèves (anonymisées), notées avec la grille. Lisez-les avant de regarder les notes : essayez de deviner laquelle est la meilleure, et pourquoi.

\begin{textebox}[Copie A — texte noté « Bien » (16/20)]
我们应该保护学校的环境。\\
Wǒmen yīnggāi bǎohù xuéxiào de huánjìng.\\
« Nous devons protéger l'environnement de l'école. »\\[0.3em]
首先，我们不要在地上扔垃圾。然后，我们应该每天打扫教室。\\
Shǒuxiān, wǒmen bú yào zài dìshang rēng lājī. Ránhòu, wǒmen yīnggāi měitiān dǎsǎo jiàoshì.\\
« D'abord, nous ne devons pas jeter de déchets par terre. Ensuite, nous devons nettoyer la salle de classe tous les jours. »\\[0.3em]
干净的环境对我们的健康很好。\\
Gānjìng de huánjìng duì wǒmen de jiànkāng hěn hǎo.\\
« Un environnement propre est très bon pour notre santé. »
\end{textebox}

\textbf{Commentaires du correcteur (Copie A) :}
\begin{itemize}
\item \textbf{Contenu (3/4)} : thèse présente dès la première phrase, deux arguments concrets (déchets, nettoyage). Un troisième argument (arbres, eau, affiches) aurait enrichi le texte.
\item \textbf{Structure (3/4)} : plan visible, mais seulement \textbf{2 connecteurs} (首先，然后). Il manque une vraie conclusion introduite par 所以 ou 总之.
\item \textbf{Vocabulaire (3/4)} : lexique correct (垃圾，打扫，健康), mais \textbf{2 fautes de caractères} : 境 écrit sans le trait final du composant 竟， et 扔 confondu avec 仍.
\item \textbf{Grammaire (4/4)} : ordre des mots correct, 应该 bien placé avant le verbe, complément de lieu 在地上 bien placé avant le verbe 扔， zéro faute de syntaxe.
\item \textbf{Présentation (3/4)} : copie lisible, quelques traits hésitants. \textbf{Total : 3 + 3 + 3 + 4 + 3 = 16/20.}
\end{itemize}

\begin{textebox}[Copie B — texte noté « Excellent » (19/20)]
学校是我们每天学习的地方，所以我们应该保护学校的环境。\\
Xuéxiào shì wǒmen měitiān xuéxí de dìfang, suǒyǐ wǒmen yīnggāi bǎohù xuéxiào de huánjìng.\\
« L'école est le lieu où nous étudions chaque jour, c'est pourquoi nous devons protéger son environnement. »\\[0.3em]
首先，我们不应该在地上扔垃圾，应该把垃圾放进垃圾桶里。\\
Shǒuxiān, wǒmen bù yīnggāi zài dìshang rēng lājī, yīnggāi bǎ lājī fàng jìn lājītǒng lǐ.\\
« D'abord, nous ne devons pas jeter de déchets par terre, nous devons les mettre dans la poubelle. »\\[0.3em]
然后，我们要每天打扫教室和校园。另外，我们也应该多种树，因为树可以让空气更干净。\\
Ránhòu, wǒmen yào měitiān dǎsǎo jiàoshì hé xiàoyuán. Lìngwài, wǒmen yě yīnggāi duō zhòng shù, yīnwèi shù kěyǐ ràng kōngqì gèng gānjìng.\\
« Ensuite, nous devons nettoyer chaque jour la classe et la cour. De plus, nous devrions aussi planter plus d'arbres, car les arbres rendent l'air plus propre. »\\[0.3em]
但是，保护环境不只是老师的事，也是每个学生的事。\\
Dànshì, bǎohù huánjìng bù zhǐ shì lǎoshī de shì, yě shì měi ge xuésheng de shì.\\
« Mais protéger l'environnement n'est pas seulement l'affaire des professeurs, c'est aussi l'affaire de chaque élève. »\\[0.3em]
总之，如果每个人都努力，我们的学校就会更美丽、更干净。\\
Zǒngzhī, rúguǒ měi ge rén dōu nǔlì, wǒmen de xuéxiào jiù huì gèng měilì, gèng gānjìng.\\
« En résumé, si chacun fait des efforts, notre école sera plus belle et plus propre. »
\end{textebox}

\begin{exemplebox}[Pourquoi la copie B obtient « Excellent »]
\begin{itemize}
\item \textbf{Contenu (4/4)} : thèse claire + 3 arguments développés (déchets, nettoyage, arbres) + une idée de responsabilité partagée.
\item \textbf{Structure (4/4)} : introduction, développement, conclusion ; \textbf{4 connecteurs variés} : 首先，然后，另外，总之， plus 所以 et 但是 pour la logique interne.
\item \textbf{Vocabulaire (4/4)} : lexique riche et précis (垃圾桶，校园，种树，空气), tous les caractères cibles corrects.
\item \textbf{Grammaire (4/4)} : 因为 et 所以 jamais employés ensemble dans la même phrase, 把 correctement employée (\zh{把垃圾放进垃圾桶里}), structure \zh{如果…就…} bien construite.
\item \textbf{Présentation (3/4)} : écriture soignée, un seul caractère un peu déséquilibré (境). \textbf{Total : 4 + 4 + 4 + 4 + 3 = 19/20.}
\end{itemize}
\end{exemplebox}

\begin{center}
\begin{tabularx}{\linewidth}{|X|X|X|X|}
\hline
\rowcolor{popblueL} \textbf{Critère} & \textbf{Copie A} & \textbf{Copie B} & \textbf{Écart expliqué} \\
\hline
Contenu et idées & 3 & 4 & B développe 3 arguments au lieu de 2 \\
\hline
Structure et connecteurs & 3 & 4 & B utilise 4 connecteurs variés + conclusion \\
\hline
Vocabulaire et caractères & 3 & 4 & A a 2 fautes de caractères (境，扔) \\
\hline
Grammaire et syntaxe & 4 & 4 & Les deux copies sont syntaxiquement correctes \\
\hline
Présentation et écriture & 3 & 3 & Lisibilité comparable \\
\hline
\textbf{Total} & \textbf{16/20} & \textbf{19/20} & La régularité sur tous les critères paie \\
\hline
\end{tabularx}
\end{center}

\courssub{Guide d'auto‑correction : la checklist avant de rendre la copie}

\begin{methode}[Relire sa copie en 5 points de contrôle]
Après la rédaction, ne rendez jamais votre texte immédiatement. Prenez trois minutes pour relire en cochant mentalement chaque case :
\begin{enumerate}
\item \textbf{Thèse claire} : ma première phrase annonce-t-elle le sujet ? (\zh{我们应该保护学校的环境。})
\item \textbf{3 connecteurs minimum} : ai-je au moins 首先，然后， et une conclusion avec 所以 ou 总之 ?
\item \textbf{8 caractères cibles corrects} : 环，境，保，护，垃，圾，打，扫 — les ai-je vérifiés un par un, trait par trait ?
\item \textbf{Pas d'erreur sur 因为…所以…} : ai-je évité de mettre les deux dans la même phrase (faute typique du francophone qui traduit « parce que… donc… ») ?
\item \textbf{Écriture lisible} : mes caractères sont-ils équilibrés, dans des cases imaginaires de même taille, sans ratures ?
\end{enumerate}
Si une case n'est pas cochée, corrigez \emph{avant} de rendre. Cette relecture vaut souvent 2 à 3 points.
\end{methode}

\begin{attention}[Les 3 fautes qui coûtent le plus cher]
\begin{itemize}
\item \textbf{Double connecteur} : \zh{因为学校很脏，所以我们要打扫。} se rencontre parfois à l'oral, mais à l'écrit, au lycée, on préfère n'en garder qu'un seul : \zh{学校很脏，所以我们要打扫。} (\emph{Xuéxiào hěn zàng, suǒyǐ wǒmen yào dǎsǎo.} — « L'école est sale, donc nous devons nettoyer. »)
\item \textbf{Place de 应该} : il se place \emph{avant} le verbe, jamais après : \zh{我们应该保护环境。} (\emph{Wǒmen yīnggāi bǎohù huánjìng.}) et non « 我们保护环境应该 ».
\item \textbf{Caractères proches confondus} : 扔 (rēng, jeter — clé de la main 扌) ≠ 仍 (réng, encore) ; 境 (jìng, environnement — clé de la terre 土) ≠ 镜 (jìng, miroir — clé du métal 钅).
\end{itemize}
\end{attention}

\courssub{Conseils pour améliorer sa note}

\begin{itemize}
\item \textbf{Enrichir le vocabulaire} : remplacez les mots trop simples par le lexique précis du thème. Au lieu de répéter 干净 (propre), utilisez 美丽 (beau), 健康 (santé), 空气 (air), 校园 (cour de l'école), 垃圾桶 (poubelle), 种树 (planter des arbres), 节约用水 (économiser l'eau).
\item \textbf{Varier les connecteurs} : ne vous contentez pas de 然后 partout. Alternez : 首先 (d'abord), 然后 (ensuite), 另外 (de plus), 但是 (mais), 因为 (parce que), 所以 (donc), 如果…就… (si… alors…), 总之 (en résumé).
\item \textbf{Soigner l'écriture} : entraînez-vous à tracer les 8 caractères cibles en respectant l'ordre des traits ; un caractère bien proportionné (clé à gauche plus étroite, composant à droite plus large, comme dans 环，保，护，垃，圾，打，扫) donne immédiatement une impression de maîtrise.
\item \textbf{Ajouter un exemple local} : une phrase comme \zh{在雅温得的学校里，我们每个星期五打扫校园。} (\emph{Zài Yǎwēndé de xuéxiào lǐ, wǒmen měi ge xīngqīwǔ dǎsǎo xiàoyuán.} — « Dans les écoles de Yaoundé, nous nettoyons la cour chaque vendredi. ») montre au correcteur que vous savez réutiliser le vocabulaire dans un contexte réel.
\end{itemize}

\begin{exoresolu}[Auto-correction d'une phrase fautive]
\textbf{Énoncé.} Un élève a écrit : \zh{因为我们应该保护学校的环境，所以我们要每天打扫教室，不要扔拉圾在地上。} Relevez toutes les erreurs, corrigez la phrase, et indiquez quels critères de la grille sont concernés.
\tcblower
\textbf{Solution détaillée.}
\begin{enumerate}
\item \textbf{Faute de caractère} : 拉圾 → \textbf{垃圾} (lājī, déchets). Le premier caractère est 垃 (clé de la terre 土), et non 拉 (lā, tirer — clé de la main 扌). Critère 3 (vocabulaire et caractères).
\item \textbf{Ordre des mots} : en chinois, le complément de lieu 在地上 se place \emph{avant} le verbe, contrairement au français (« jeter des déchets \emph{par terre} »). Il faut donc écrire \zh{不要在地上扔垃圾} et non « 不要扔垃圾在地上 ». Critère 4 (grammaire).
\item \textbf{因为…所以…} : les deux ensemble sont à éviter à l'écrit ; on supprime 因为 : la phrase reste logique et plus élégante. Critère 4 (grammaire).
\end{enumerate}
\textbf{Phrase corrigée} : \zh{我们应该保护学校的环境，所以要每天打扫教室，不要在地上扔垃圾。} (\emph{Wǒmen yīnggāi bǎohù xuéxiào de huánjìng, suǒyǐ yào měitiān dǎsǎo jiàoshì, bú yào zài dìshang rēng lājī.} — « Nous devons protéger l'environnement de l'école, donc nous devons nettoyer la classe chaque jour et ne pas jeter de déchets par terre. »)
\end{exoresolu}

\courssub{Le lien avec l'oral : lire sa copie à voix haute}

Une dernière astuce de correcteur : \textbf{relisez votre texte à voix haute} (ou à mi-voix) avant de le rendre. La lecture orale révèle ce que l'œil ne voit plus :

\begin{itemize}
\item \textbf{Fluidité} : si vous « butez » sur une phrase, c'est souvent qu'elle est mal construite (ordre des mots, connecteur manquant). Une phrase correcte se lit d'une traite.
\item \textbf{Tons} : prononcer les tons vous oblige à vérifier les mots : 环境 (huánjìng, 2\textsuperscript{e}–4\textsuperscript{e} tons), 保护 (bǎohù, 3\textsuperscript{e}–4\textsuperscript{e}), 垃圾 (lājī, 1\textsuperscript{er}–1\textsuperscript{er}). Si vous hésitez sur un ton, vérifiez le caractère : l'hésitation signale souvent un caractère mal appris.
\item \textbf{Pauses logiques} : chaque connecteur (首先，然后，总之) marque une petite pause ; si votre texte n'a aucun endroit où respirer, il manque de connecteurs — critère 2 de la grille.
\end{itemize}

\begin{experience}[Correction croisée en binôme]
Échangez votre production avec celle d'un camarade. Notez sa copie avec la grille des 5 critères, puis lisez-lui vos commentaires en commençant par un point fort avant chaque point faible. Comparez ensuite vos deux grilles : les écarts de notation sont d'excellents sujets de discussion. Terminez en lisant votre propre texte corrigé à voix haute devant la classe.
\end{experience}

\begin{savaistu}[Ce que regardent vraiment les correcteurs]
Aux examens officiels, au Cameroun comme en Chine, les correcteurs de production écrite accordent une attention particulière à deux choses : la \textbf{cohérence logique} (les connecteurs, qui montrent que l'élève sait organiser sa pensée) et la \textbf{propreté de l'écriture} (des traits nets, dans le bon ordre, témoignent d'un apprentissage sérieux des caractères). En Chine, on dit qu'une belle écriture est « le visage de l'élève » : \zh{字如其人} (\emph{zì rú qí rén} — « les caractères ressemblent à leur auteur »). Soigner sa copie, c'est déjà gagner des points avant même d'être lu.
\end{savaistu}

\begin{aretenir}[L'essentiel de la section]
\begin{itemize}
\item La grille d'évaluation comporte \textbf{5 critères} (contenu, structure, vocabulaire, grammaire, présentation), chacun pesant \textbf{20 \%} de la note sur 20.
\item Avant de rendre : checklist en 5 points — thèse claire, 3 connecteurs minimum, 8 caractères cibles vérifiés, pas de double 因为…所以…, écriture lisible.
\item Les fautes les plus coûteuses : caractères proches confondus (扔/仍，境/镜，垃/拉), complément de lieu mal placé (在地上 \emph{avant} le verbe), connecteurs absents.
\item Pour progresser : enrichir le lexique du thème, varier les connecteurs, s'entraîner à l'ordre des traits, et \textbf{relire à voix haute} pour tester fluidité et tons.
\end{itemize}
\end{aretenir}

\newpage

\coursec{Activité d'intégration}

\coursec{Leçon 18 — Expression écrite : Protection de l'environnement (你应该怎么保护学校的环境)}

\courssub{Objectifs de la leçon}
À l'issue de cette leçon, l'élève sera capable de :
\begin{itemize}
\item \textbf{Comprendre} un texte authentique (annonce de concours) et en extraire le vocabulaire clé de l'environnement ;
\item \textbf{Rédiger} un texte court, structuré et correct en chinois (titre, 4--5 actions, conclusion) sur le thème « Protéger l'environnement de l'école » ;
\item \textbf{Maîtriser} les structures grammaticales : \zh{应该} + V (devoir/conseil), \zh{不要} + V (interdiction), \zh{为了} + V/Nom, … (but), connecteurs \zh{首先/然后/还有/最后} ;
\item \textbf{Écrire} correctement les caractères du thème (radicaux, ordre des traits, distinctions de caractères proches) ;
\item \textbf{Traduire} fidèlement du chinois vers le français (version) et du français vers le chinois (thème) des phrases simples ;
\item \textbf{Présenter} oralement sa production en chinois (1 min) puis en français, avec tons et prononciation soignés ;
\item \textbf{S'auto-évaluer} à l'aide d'une grille critériée (contenu, grammaire, caractères, traduction, oral).
\end{itemize}

\courssub{Situation déclenchante : La Semaine de l'environnement au Lycée bilingue de Yaoundé}

\begin{textebox}[Annonce officielle du concours — 通知]
通知：环境周海报比赛\\
Tōngzhī: Huánjìng zhōu hǎibào bǐsài.\\
« Avis : Concours d'affiches de la Semaine de l'environnement. »\\[4pt]
为了保护我们学校的环境，请每个小组做一张海报。\\
Wèile bǎohù wǒmen xuéxiào de huánjìng, qǐng měi ge xiǎozǔ zuò yì zhāng hǎibào.\\
« Pour protéger l'environnement de notre école, chaque groupe est invité à réaliser une affiche. »\\[4pt]
海报上要有：一个中文题目、四到五个行动（用中文写）、法语翻译和图画。\\
Hǎibào shang yào yǒu: yí ge Zhōngwén tímù, sì dào wǔ ge xíngdòng (yòng Zhōngwén xiě), Fǎyǔ fānyì hé túhuà.\\
« L'affiche doit comporter : un titre en chinois, quatre à cinq actions (rédigées en chinois), la traduction française et des dessins. »\\[4pt]
每个小组用中文介绍一分钟，然后用法语介绍。\\
Měi ge xiǎozǔ yòng Zhōngwén jièshào yì fēnzhōng, ránhòu yòng Fǎyǔ jièshào.\\
« Chaque groupe présente son affiche en chinois pendant une minute, puis en français. »
\end{textebox}

\begin{experience}[Remue-méninges guidé — 10 min]
\textbf{Consigne :} En groupe de 3, lisez l'annonce. Soulignez les mots connus du champ lexical « environnement ». Listez 8 actions concrètes pour verdir l'école (ex. : ne pas jeter de papiers, trier les déchets, économiser l'eau/électricité, planter des arbres, protéger les fleurs, ramasser les déchets, utiliser les deux faces du papier, apporter sa gourde). Partagez au tableau.
\end{experience}

\courssub{Structures grammaticales et connecteurs logiques}

\begin{propriete}[Structure 1 : Devoir / Conseil — \zh{应该} (yīnggāi)]
\textbf{Schéma :} Sujet + \zh{应该} + Verbe (+ Complément).\\
\textbf{Emploi :} Exprime une obligation morale, un conseil, une recommandation. Négation : \zh{不应该} (ne pas devoir).\\
\textbf{Exemples :}
\begin{itemize}
\item \zh{我们应该保护环境。} (Wǒmen yīnggāi bǎohù huánjìng.) — Nous devons protéger l'environnement.
\item \zh{同学们应该节约用水。} (Tóngxuémen yīnggāi jiéyuē yòng shuǐ.) — Les camarades doivent économiser l'eau.
\item \zh{你不应该乱扔垃圾。} (Nǐ bù yīnggāi luàn rēng lājī.) — Tu ne devrais pas jeter les déchets n'importe où.
\end{itemize}
\textbf{Attention :} \zh{应该} se place \textbf{avant} le verbe. *\zh{保护环境很应该} est faux.
\end{propriete}

\begin{propriete}[Structure 2 : Interdiction / Défense — \zh{不要} (bú yào)]
\textbf{Schéma :} Sujet + \zh{不要} + Verbe (+ Complément).\\
\textbf{Emploi :} Ordre négatif, défense, conseil de ne pas faire. Plus fort que \zh{不} + V seul.\\
\textbf{Exemples :}
\begin{itemize}
\item \zh{我们不要在校园里扔垃圾。} (Wǒmen bú yào zài xiàoyuán lǐ rēng lājī.) — Ne jetons pas de déchets dans l'école.
\item \zh{请大家不要摘花。} (Qǐng dàjiā bú yào zhāi huā.) — Merci à tous de ne pas cueillir les fleurs.
\end{itemize}
\textbf{Note phonétique :} \zh{不} (bù) 4\textsuperscript{e} ton $\to$ \zh{不} (bú) 2\textsuperscript{e} ton devant \zh{要} (yào) 4\textsuperscript{e} ton (règle du tone sandhi).
\end{propriete}

\begin{propriete}[Structure 3 : But / Finalité — \zh{为了} (wèile)]
\textbf{Schéma :} \zh{为了} + Verbe / Nom, Sujet + \zh{应该/要/会} + Verbe \dots\\
\textbf{Emploi :} Introduit le but de l'action principale. La proposition en \zh{为了} précède le sujet.\\
\textbf{Exemples :}
\begin{itemize}
\item \zh{为了保护环境，我们应该种树。} (Wèile bǎohù huánjìng, wǒmen yīnggāi zhòng shù.) — Pour protéger l'environnement, nous devons planter des arbres.
\item \zh{为了节约电，离开教室请关灯。} (Wèile jiéyuē diàn, líkāi jiàoshì qǐng guān dēng.) — Pour économiser l'électricité, merci d'éteindre en quittant la classe.
\end{itemize}
\end{propriete}

\begin{propriete}[Connecteurs logiques pour structurer un texte]
\begin{center}
\begin{tabularx}{\linewidth}{|X|X|X|}
\hline
\rowcolor{popblueL} Connecteur & Pinyin & Emploi / Traduction \\
\hline
首先 & shǒuxiān & D'abord, en premier lieu (initie la liste) \\
\hline
然后 & ránhòu & Ensuite, puis (chronologie / suite logique) \\
\hline
还有 & háiyǒu & De plus, en outre (ajout d'un argument) \\
\hline
最后 & zuìhòu & Enfin, pour terminer (clôture la liste) \\
\hline
\end{tabularx}
\end{center}
\textbf{Modèle de paragraphe :} \zh{首先，\dots 然后，\dots 还有，\dots 最后，\dots}
\end{propriete}

\courssub{Vocabulaire thématique : Environnement et actions éco-citoyennes}

\begin{center}
\begin{tabularx}{\linewidth}{|X|X|X|X|}
\hline
\rowcolor{popblueL} Caractères & Pinyin & Nature & Traduction \\
\hline
环境 & huánjìng & n. & environnement \\
\hline
保护 & bǎohù & v. & protéger \\
\hline
垃圾 & lājī & n. & déchets, ordures \\
\hline
扔 & rēng & v. & jeter (lancer) \\
\hline
节约 & jiéyuē & v. & économiser \\
\hline
用水 & yòng shuǐ & loc. v. & utiliser de l'eau / consommer de l'eau \\
\hline
用电 & yòng diàn & loc. v. & utiliser de l'électricité / consommer de l'électricité \\
\hline
关灯 & guān dēng & loc. v. & éteindre la lumière/lampe \\
\hline
开灯 & kāi dēng & loc. v. & allumer la lumière \\
\hline
种树 & zhòng shù & loc. v. & planter des arbres (\zh{zhòng} = planter) \\
\hline
校园 & xiàoyuán & n. & campus, cour de l'école \\
\hline
爱护 & àihù & v. & prendre soin de, chérir \\
\hline
花草 & huācǎo & n. & fleurs et plantes, verdure \\
\hline
摘 & zhāi & v. & cueillir, ramasser (fruit, fleur) \\
\hline
绿色 & lǜsè & n./adj. & vert, écologique \\
\hline
海报 & hǎibào & n. & affiche \\
\hline
题目 & tímù & n. & titre, sujet \\
\hline
行动 & xíngdòng & n. & action \\
\hline
美丽 & měilì & adj. & beau, joli \\
\hline
让 & ràng & v. & faire que, laisser, rendre (causatif) \\
\hline
从\dots 做起 & cóng\dots zuòqǐ & loc. & commencer par\dots, à partir de\dots \\
\hline
\end{tabularx}
\end{center}

\courssub{Écriture correcte des caractères du thème : Radicaux, traits et pièges}

\begin{aretenir}[Analyse graphique des caractères clés]
\begin{itemize}
\item \zh{环} (huán) : Radical \zh{王} (jade/roi) + \zh{奂}. Ordre : haut \zh{奂} (point, horizontales, verticale), bas \zh{王}. \textbf{Piège :} ne pas confondre avec \zh{还} (huán/hái, radical \zh{辶}).
\item \zh{境} (jìng) : Radical \zh{土} (terre) + \zh{镸} (phonétique). \textbf{Distinction :} \zh{境} (environnement, frontière) vs \zh{镜} (jìng, miroir, radical \zh{钅} métal).
\item \zh{保} (bǎo) : Radical \zh{亻} (homme) + \zh{呆}. Homme qui veille sur quelque chose de statique.
\item \zh{护} (hù) : Radical \zh{扌} (main) + \zh{护} (phonétique \zh{户} + \zh{言}). \textbf{Distinction :} \zh{护} (protéger, main) vs \zh{户} (hù, porte, foyer, pas de main).
\item \zh{垃} (lā) \& \zh{圾} (jī) : Tous deux radical \zh{土} (terre/déchet). \zh{垃} : \zh{土} + \zh{拉} ; \zh{圾} : \zh{土} + \zh{及}.
\item \zh{节} (jié) : Radical \zh{竹} (bambou, haut) + \zh{即}. \textbf{Distinction :} \zh{节} (économiser, nœud, fête) vs \zh{苗} (miáo, plantule, radical \zh{艹} herbe).
\item \zh{约} (yuē) : Radical \zh{纟} (soie) + \zh{勺}. Lier avec du fil $\to$ accord, économie.
\item \zh{树} (shù) : Radical \zh{木} (bois/arbre) + \zh{对}. \textbf{Ordre :} horizontale, verticale, gauche, droite, horizontale (bois), puis \zh{对}.
\item \zh{种} (zhòng/zhǒng) : Radical \zh{禾} (riz/herbe) + \zh{中}. \zh{zhòng} (verbe, planter) vs \zh{zhǒng} (nom, espèce/graine).
\item \zh{洁} (jié) : Radical \zh{氵} (eau) + \zh{齐}. Propre $\to$ eau qui rend net. (Apparaît dans \zh{干净} gānjìng).
\end{itemize}
\end{aretenir}

\begin{methode}[Ordre des traits : Règles d'or pour l'examen]
\begin{enumerate}
\item De haut en haut (上 → 下) : \zh{三}、\zh{草} (partie haute).
\item De gauche à droite (左 → 右) : \zh{好}、\zh{树} (radical bois à gauche).
\item Horizontales avant verticales (横 → 竖) : \zh{十}、\zh{王}.
\item Traits centraux avant latéraux symétriques : \zh{小}、\zh{水}.
\item Fermer en dernier (fermer le cadre) : \zh{口}、\zh{国}、\zh{园}.
\item Le trait qui traverse s'écrit à la fin : \zh{中}、\zh{王}、\zh{董}.
\end{enumerate}
\textbf{Exercice immédiat :} Sur votre cahier d'écriture (格子纸), recopiez 5 fois chaque caractère ci-dessus en respectant l'ordre. Vérifiez par un pair.
\end{methode}

\courssub{Entraînement : Thème (Français $\to$ Chinois) et Version (Chinois $\to$ Français)}

\begin{exercice}[Thème : Traduisez en chinois (caractères + pinyin)]
\begin{enumerate}
\item Pour protéger l'environnement, nous devons économiser l'eau.
\item Ne jetez pas les déchets dans la cour de l'école.
\item Merci à tous de prendre soin des fleurs et des plantes.
\item D'abord, nous devons éteindre les lumières ; ensuite, nous devons planter des arbres.
\item Enfin, commençons par nous-mêmes pour rendre l'école plus belle.
\end{enumerate}
\end{exercice}

\begin{corrige}[Exercice Thème]
\begin{enumerate}
\item \zh{为了保护环境，我们应该节约用水。} (Wèile bǎohù huánjìng, wǒmen yīnggāi jiéyuē yòng shuǐ.)
\item \zh{不要在校园里扔垃圾。} / \zh{请大家不要在校园里扔垃圾。} (Bú yào zài xiàoyuán lǐ rēng lājī. / Qǐng dàjiā bú yào zài xiàoyuán lǐ rēng lājī.)
\item \zh{请大家爱护花草。} (Qǐng dàjiā àihù huācǎo.)
\item \zh{首先，我们应该关灯；然后，我们应该种树。} (Shǒuxiān, wǒmen yīnggāi guān dēng; ránhòu, wǒmen yīnggāi zhòng shù.)
\item \zh{最后，从我做起，让学校更美丽。} (Zuìhòu, cóng wǒ zuòqǐ, ràng xuéxiào gèng měilì.)
\end{enumerate}
\end{corrige}

\begin{exercice}[Version : Traduisez en français]
\begin{enumerate}
\item 我们应该每天捡垃圾。
\item 为了节约用电，离开教室的时候要关灯。
\item 请大家不要踩草坪。
\item 首先，我们要分类垃圾；然后，我们要回收纸张；最后，我们要种树。
\item 绿色校园，从我做起，让我们的学校更美丽！
\end{enumerate}
\end{exercice}

\begin{corrige}[Exercice Version]
\begin{enumerate}
\item Nous devons ramasser les déchets chaque jour. (\zh{捡} jiǎn = ramasser).
\item Pour économiser l'électricité, il faut éteindre la lumière en quittant la salle de classe.
\item Merci à tous de ne pas marcher sur le gazon. (\zh{踩} cǎi = marcher sur/piétiner ; \zh{草坪} cǎopíng = pelouse/gazon).
\item D'abord, nous devons trier les déchets ; ensuite, nous devons recycler le papier ; enfin, nous devons planter des arbres. (\zh{分类} fēnlèi = trier/classer ; \zh{回收} huíshōu = recycler ; \zh{纸张} zhǐzhāng = papier).
\item Un campus vert commence par moi, rendons notre école plus belle !
\end{enumerate}
\end{corrige}

\courssub{Tâche finale intégrée : Création de l'affiche bilingue « École verte »}

\begin{methode}[Consignes détaillées — Durée totale : 2 h]
\begin{enumerate}
\item \textbf{Préparation \& Rédaction (1 h 15).}
\begin{itemize}
\item \textit{Étape 1 (10 min) :} Lecture de l'annonce, brainstorming vocabulaire (voir \textit{Expérience} p.1).
\item \textit{Étape 2 (25 min) :} Choix de 5 actions. Rédaction en chinois \textbf{obligatoire} : 1 phrase avec \zh{我们应该}, 1 avec \zh{我们不要}, 1 avec \zh{请大家}, 2 libres. Utilisation des 4 connecteurs \zh{首先/然后/还有/最后}. Titre imposé ou au choix (ex. \zh{绿色校园，从我做起}).
\item \textit{Étape 3 (15 min) :} \textbf{Calligraphie / Copie au propre.} Recopie sur feuille A4 en gros caractères. Vérification croisée (binôme) : radicaux, traits, classificateurs (\zh{一棵树}, \zh{一张纸}, \zh{一个瓶子}). \textbf{Interdit :} correcteur automatique, dictionnaire pour l'écriture des caractères (mémoire only).
\item \textit{Étape 4 (15 min) :} Traduction française en regard (version). Vérification fidélité (sujet, temps, négation).
\item \textit{Étape 5 (10 min) :} Mise en page : Titre centré, phrases numérotées, traduction alignée, pictogrammes (poubelle \zh{}, robinet \zh{}, ampoule \zh{}, arbre \zh{}, main/feuille \zh{}).
\end{itemize}
\item \textbf{Présentation orale \& Évaluation (45 min).}
\begin{itemize}
\item Chaque groupe passe au tableau (1 min chinois + 30 s français). Les 3 membres parlent en chinois.
\item Auditoire : Grille d'écoute (Titre ? 5 actions ? Tons ?).
\item Auto-évaluation groupe par la \textbf{Grille officielle} (ci-dessous). Remise à l'enseignant pour validation.
\end{itemize}
\end{enumerate}
\end{methode}

\begin{aretenir}[Grille d'évaluation — Total 20 points]
\begin{center}
\begin{tabularx}{\linewidth}{|X|c|X|X|}
\hline
\rowcolor{popblueL} \textbf{Critère} & \textbf{Pts} & \textbf{Indicateurs de réussite} \\
\hline
Contenu / Respect consignes & 4 & Titre chinois + 5 actions + traduction + dessins. Thème respecté. \\
\hline
Grammaire & Structures & 4 & \zh{应该/不要/为了/请大家} utilisés correctement. 4 connecteurs présents. Ordre mots correct (Suj + Adv + V). \\
\hline
Caractères & Écriture & 4 & Caractères lisibles, radicaux justes, pas de confusion (境/镜, 护/户, 节/苗), classificateurs justes (\zh{棵}, \zh{张}). \\
\hline
Traduction (Version) & 4 & Fidèle, naturel en français, pas de contresens, pas d'invention. \\
\hline
Expression orale & 4 & Prononciation claire, tons respectés, fluidité, participation équitable des 3 membres. \\
\hline
\textbf{TOTAL} & \textbf{20} &  \\
\hline
\end{tabularx}
\end{center}
\end{aretenir}

\courssub{Modèle commenté : Affiche « École verte » (Production attendue)}

\begin{exoresolu}[Modèle d'affiche complète — Corrigé type]
Énoncé : Produire le texte final de l'affiche bilingue pour le concours.
\tcblower
\textbf{Titre :}\\
\zh{绿色校园，从我做起}\\
\emph{Lǜsè xiàoyuán, cóng wǒ zuòqǐ.}\\
« Un campus vert, commençons par moi. »\\[6pt]
\textbf{Corps de l'affiche :}\\
\zh{为了保护我们学校的环境，我们应该这样做：}\\
\emph{Wèile bǎohù wǒmen xuéxiào de huánjìng, wǒmen yīnggāi zhèyàng zuò:}\\
« Pour protéger l'environnement de notre école, nous devons agir ainsi : »\\[4pt]
\zh{首先，我们不要在校园里乱扔垃圾。}\\
\emph{Shǒuxiān, wǒmen bú yào zài xiàoyuán lǐ luàn rēng lājī.}\\
« D'abord, ne jetons pas les déchets n'importe où dans l'école. »\\[4pt]
\zh{然后，我们应该节约用水，也要节约用电。}\\
\emph{Ránhòu, wǒmen yīnggāi jiéyuē yòng shuǐ, yě yào jiéyuē yòng diàn.}\\
« Ensuite, nous devons économiser l'eau et aussi économiser l'électricité. »\\[4pt]
\zh{还有，请大家爱护花草树木，不要踩草坪、不要摘花。}\\
\emph{Háiyǒu, qǐng dàjiā àihù huācǎo shùmù, bú yào cǎi cǎopíng, bú yào zhāi huā.}\\
« De plus, merci à tous de prendre soin des fleurs et arbres, ne marchez pas sur le gazon, ne cueillez pas les fleurs. »\\[4pt]
\zh{最后，我们应该每个学期种一棵树。}\\
\emph{Zuìhòu, wǒmen yīnggāi měi ge xuéqī zhòng yì kē shù.}\\
« Enfin, nous devrions planter un arbre chaque semestre. »\\[6pt]
\textbf{Conclusion / Slogan :}\\
\zh{行动起来，让校园更美丽！}\\
\emph{Xíngdòng qǐlái, ràng xiàoyuán gèng měilì!}\\
« Passons à l'action, rendons le campus plus beau ! »\\[6pt]
\textbf{Analyse linguistique du modèle :}
\begin{itemize}
\item \textbf{Cohésion :} Les 4 connecteurs structurent visuellement et logiquement le discours.
\item \textbf{Variété syntaxique :} Alternance \zh{不要} (interdiction n°1), \zh{应该} (obligation n°2, n°5), \zh{请大家} + \zh{不要} (appel poli + interdiction n°3). Utilisation de \zh{也要} (aussi devoir) pour lier eau/électricité.
\item \textbf{Précision lexicale :} \zh{乱扔} (jeter n'importe comment), \zh{踩草坪} (piétiner pelouse), \zh{花草树木} (terme composé standard « flore »), \zh{每个学期} (chaque semestre - réaliste), \zh{一棵树} (classificateur \zh{棵} pour arbres).
\item \textbf{Conclusion :} \zh{行动起来} (se mettre en action / agir maintenant), \zh{让} + \zh{更} + adj (structure causative comparative).
\end{itemize}
\end{exoresolu}

\courssub{Approfondissement socioculturel : Cameroun — Chine}

\begin{savaistu}[Gestion de l'environnement : Regards croisés Cameroun / Chine]
\begin{itemize}
\item \textbf{Cameroun :} Yaoundé et Douala font face aux défis de la gestion des déchets plastiques (bouteilles, sachets \zh{塑料袋} sùliào dài). Initiatives : « Opération ville propre » (nettoyage communautaire le samedi), tri sélectif naissant (ONG, jeunes), reboisement « Green Sahel » au Nord. Le \zh{ndolé} (feuilles de bitterleaf) et le bois de chauffe posent la question de la déforestation locale.
\item \textbf{Chine :} Politique « \zh{垃圾分类} (lājī fēnlèi) » obligatoire à Shanghai (2019), Pékin, Shenzhen : 4 poubelles (recyclable \zh{可回收}, dangereux \zh{有害}, humide/cuisine \zh{厨余}, autres \zh{其他}). Amendes pour non-tri. Projet « \zh{三北防护林} » (Grande Muraille Verte) : plantation massive d'arbres au Nord-Ouest contre la désertification (similaire à la Grande Muraille Verte africaine). Objectif « \zh{碳达峰} (tándáfēng) pic carbone 2030 / \zh{碳中和} (tánzhōnghé) neutralité 2060 ».
\item \textbf{Point commun :} L'école est le levier principal. En Chine, « \zh{绿色学校} » (écoles vertes) est un label officiel. Au Cameroun, les clubs « Environnement » et « Jardinage » se multiplient dans les lycées bilingues.
\item \textbf{Partenariat :} Échanges scolaires Yaoundé–Kunming (Yunnan, ville « Printemps éternel », biodiversité) : projets communs pépinières, énergie solaire, affiches bilingues (comme ce concours !).
\end{itemize}
\end{savaistu}

\courssub{Bilan et perspectives}

Cette leçon a permis de \textbf{mobiliser en situation authentique} l'ensemble des compétences du Module 2 : compréhension de l'écrit (annonce), expression écrite (affiche structurée), grammaire (\zh{应该/不要/为了} + connecteurs), lexique (environnement), écriture (caractères à radicaux \zh{土}, \zh{木}, \zh{氵}, \zh{扌}, \zh{竹}), traduction (thème/version), expression orale (présentation), culture (comparaison CM/CN).

\begin{attention}[Check-list finale avant l'évaluation certificative (Épreuve écrite 1h)]
\begin{itemize}
\item Je sais écrire \textbf{de mémoire} les 15 caractères du tableau vocabulaire (radicaux, traits, pas de confusion).
\item Je sais construire une phrase avec \zh{为了\dots{}，应该\dots} et une avec \zh{不要\dots}.
\item J'utilise \textbf{systématiquement} \zh{首先/然后/还有/最后} pour structurer un paragraphe de 4--5 phrases.
\item Je connais les classificateurs : \zh{棵} (arbre), \zh{张} (affiche, papier, table), \zh{个} (général), \zh{只} (animal), \zh{本} (livre).
\item Je relis mon texte en me posant 3 questions : \textbf{Sens ?} (Traduction OK) \textbf{Grammaire ?} (Ordres, outils) \textbf{Caractères ?} (Traits, radicaux).
\end{itemize}
\textbf{Devoirs :} Finalisez votre affiche personnelle (version individuelle du travail de groupe) pour la note de devoir. Révisez les caractères pour le contrôle de la semaine prochaine (dictée 10 mots + écriture 5 mots).
\end{attention}


\newpage

\coursec{Méthodes et exercices résolus}

\coursec{Leçon 18 — Expression écrite : protection de l'environnement (你应该怎么保护学校的环境)}

\begin{center}
\begin{tabularx}{\linewidth}{|X|X|X|X|}
\hline
\rowcolor{popblueL} Caractères & Pinyin & Nature & Traduction \\
\hline
环境 & huánjìng & nom & environnement \\
\hline
保护 & bǎohù & verbe & protéger \\
\hline
垃圾 & lājī & nom & déchets, ordures \\
\hline
垃圾桶 & lājītǒng & nom & poubelle \\
\hline
扔 & rēng & verbe & jeter \\
\hline
地上 & dìshang & nom de lieu & par terre, au sol \\
\hline
干净 & gānjìng & adj. & propre \\
\hline
脏 & zāng & adj. & sale \\
\hline
节约 & jiéyuē & verbe & économiser \\
\hline
用水 & yòngshuǐ & verbe + obj. & utiliser de l'eau / consommer de l'eau \\
\hline
种树 & zhòng shù & verbe + obj. & planter des arbres \\
\hline
空气 & kōngqì & nom & air \\
\hline
应该 & yīnggāi & verbe modal & devoir (obligation morale) \\
\hline
不要 & bú yào & verbe modal & ne pas (interdiction) \\
\hline
首先 & shǒuxiān & adv. & d'abord, en premier \\
\hline
然后 & ránhòu & adv. & ensuite, puis \\
\hline
因为……所以…… & yīnwèi… suǒyǐ… & conj. & parce que… donc… \\
\hline
但是 & dànshì & conj. & mais \\
\hline
还 & hái & adv. & aussi, en plus \\
\hline
乱扔 & luàn rēng & verbe & jeter n'importe où / salir \\
\hline
爱护 & àihù & verbe & prendre soin de, chérir \\
\hline
公园 & gōngyuán & nom & parc \\
\hline
河 & hé & nom & fleuve, rivière \\
\hline
\end{tabularx}
\end{center}

\courssub{Modèle de texte commenté : « Comment protéger l'environnement de l'école »}

\begin{textebox}[Modèle : 我们应该保护学校的环境]
\zh{我们的学校很大，也很漂亮。但是有时候地上有很多垃圾。我们应该保护学校的环境。首先，我们不要把垃圾扔在地上，要把垃圾放进垃圾桶里。然后，我们应该节约用水。我们还应该多种树，因为树可以让空气更干净。学校是我们的家，所以我们要爱护它。} \\

(Wǒmen de xuéxiào hěn dà, yě hěn piàoliang. Dànshì yǒu shíhou dìshang yǒu hěn duō lājī. Wǒmen yīnggāi bǎohù xuéxiào de huánjìng. Shǒuxiān, wǒmen bú yào bǎ lājī rēng zài dìshang, yào bǎ lājī fàng jìn lājītǒng lǐ. Ránhòu, wǒmen yīnggāi jiéyuē yòngshuǐ. Wǒmen hái yīnggāi duō zhòng shù, yīnwèi shù kěyǐ ràng kōngqì gèng gānjìng. Xuéxiào shì wǒmen de jiā, suǒyǐ wǒmen yào àihù tā.)

« Notre école est grande et très belle. Mais parfois il y a beaucoup de déchets par terre. Nous devons protéger l'environnement de l'école. D'abord, nous ne devons pas jeter les déchets par terre, nous devons les mettre dans la poubelle. Ensuite, nous devons économiser l'eau. Nous devons aussi planter plus d'arbres, parce que les arbres rendent l'air plus propre. L'école est notre maison, donc nous devons en prendre soin. »
\end{textebox}

\begin{methode}[Analyse du modèle et plan de rédaction]
Pour écrire un texte du type \zh{你应该怎么保护学校的环境} (Nǐ yīnggāi zěnme bǎohù xuéxiào de huánjìng ?), suis ces étapes :

\begin{enumerate}
\item \textbf{Comprendre le sujet} : repère le thème (\zh{保护} protéger), le lieu (\zh{学校的环境} l'environnement de l'école) et le mot interrogatif (\zh{怎么} comment → on attend des \emph{moyens d'action} concrets).
\item \textbf{Structurer en 3 parties} :
  \begin{itemize}
  \item \textbf{Introduction} (1-2 phrases) : présenter le constat (ex. \zh{我们的学校很漂亮，但是有时候不太干净}).
  \item \textbf{Développement} (3-4 phrases) : lister 2 ou 3 actions concrètes avec \zh{我们应该……} (nous devons…) ou \zh{我们不要……} (nous ne devons pas…). Utiliser \zh{首先}, \zh{然后}, \zh{还}.
  \item \textbf{Conclusion} (1 phrase) : appel à l'action ou souhait (ex. \zh{学校是我们的家，所以我们要爱护它}).
  \end{itemize}
\item \textbf{Mobiliser les structures clés} :
  \begin{itemize}
  \item \zh{我们应该 + verbe} (obligation) : \zh{我们应该保护环境}.
  \item \zh{不要 + verbe} (interdiction) : \zh{不要乱扔垃圾}.
  \item \zh{把 + objet + verbe + compl.} (manipulation d'objet) : \zh{把垃圾放进垃圾桶里}.
  \item \zh{因为……所以……} (cause/conséquence) : \zh{因为树可以让空气更干净，所以我们应该多种树}.
  \end{itemize}
\item \textbf{Relire} : vérifier l'ordre \textbf{Sujet + Temps + Lieu (\zh{在}...) + Verbe + Objet}, les tons, la ponctuation pleine chasse (。 ，) et les classificateurs (\zh{棵} pour les arbres).
\end{enumerate}
\end{methode}

\begin{exoresolu}[Rédaction guidée : produire un texte de 6 à 8 phrases]
\textbf{Énoncé.} Rédige en chinois un court texte intitulé \zh{我们应该保护学校的环境} en utilisant au moins deux connecteurs (\zh{首先}, \zh{然后}, \zh{因为}, \zh{所以}) et la structure \zh{应该}.
\tcblower
\textbf{Solution détaillée.}

\textbf{Étape 1 — Plan :} (a) Constat : l'école est belle mais sale par moments. (b) Actions : trier les déchets, économiser l'eau, planter des arbres. (c) Conclusion : l'école est notre foyer commun.

\textbf{Étape 2 — Choix linguistiques :} \zh{应该} (devoir), \zh{不要} (interdiction), \zh{把} (objet anticipé), connecteurs \zh{首先}, \zh{然后}, \zh{因为……所以……}.

\textbf{Étape 3 — Texte produit (8 phrases) :}

\zh{我们的学校很漂亮，但是有时候地上有垃圾。我们应该保护学校的环境。首先，我们不要把垃圾扔在地上，要把垃圾放进垃圾桶里。然后，我们应该节约用水，不浪费水。我们还应该多种树，因为树可以让空气更干净。如果人人都努力，我们的学校会更美丽。我爱我们的学校。}

(Wǒmen de xuéxiào hěn piàoliang, dànshì yǒu shíhou dìshang yǒu lājī. Wǒmen yīnggāi bǎohù xuéxiào de huánjìng. Shǒuxiān, wǒmen bú yào bǎ lājī rēng zài dìshang, yào bǎ lājī fàng jìn lājītǒng lǐ. Ránhòu, wǒmen yīnggāi jiéyuē yòngshuǐ, bù làngfèi shuǐ. Wǒmen hái yīnggāi duō zhòng shù, yīnwèi shù kěyǐ ràng kōngqì gèng gānjìng. Rúguǒ rénrén dōu nǔlì, wǒmen de xuéxiào huì gèng měilì. Wǒ ài wǒmen de xuéxiào.)

\textbf{Justifications :} \zh{把垃圾扔在地上} (construction \zh{把} + compl. de lieu résultatif) ; \zh{不浪费水} (négation courte \zh{不} + verbe) ; \zh{更干净} (comparatif \zh{更} + adj.) ; \zh{如果……会……} (condition/conséquence). Le texte compte 8 phrases, 3 connecteurs (\zh{首先}, \zh{然后}, \zh{因为}) et respecte le plan en 3 parties.
\end{exoresolu}

\courssub{Structure et connecteurs logiques}

\begin{methode}[Choisir et placer le bon connecteur]
\begin{enumerate}
\item \textbf{Identifier la relation logique} entre les deux propositions : chronologie, addition, cause, conséquence, opposition, condition.
\item \textbf{Choisir le connecteur adapté} (tableau de référence) :
\end{enumerate}
\begin{center}
\begin{tabularx}{\linewidth}{|X|X|X|X|}
\hline
\rowcolor{popblueL} Relation & Connecteur & Pinyin & Emploi / Exemple \\
\hline
Chronologie & 首先……然后…… & shǒuxiān… ránhòu… & \zh{首先扫地，然后擦桌子} (D'abord balayer, ensuite essuyer) \\
\hline
Addition & 还 (hái) / 而且 (érqiě) & hái / érqiě & \zh{我们要扫地，还要擦桌子} (Balayer \textbf{et aussi} essuyer) \\
\hline
Cause & 因为 (yīnwèi) & yīnwèi & Place en tête de la 1\up{re} prop. : \zh{因为环境重要，所以我们保护它} \\
\hline
Conséquence & 所以 (suǒyǐ) / 因此 (yīncǐ) & suǒyǐ / yīncǐ & Place en tête de la 2\up{e} prop. : \zh{环境重要，所以我们保护它} \\
\hline
Opposition & 但是 (dànshì) / 可是 (kěshì) & dànshì / kěshì & \zh{学校很美，但是有垃圾} (Beau \textbf{mais} des déchets) \\
\hline
Condition & 如果……就…… (rúguǒ… jiù…) & rúguǒ… jiù… & \zh{如果人人努力，环境就好} (Si tous font des efforts, l'env. sera bien) \\
\hline
\end{tabularx}
\end{center}
\begin{enumerate}
\setcounter{enumi}{2}
\item \textbf{Règle de position} : le connecteur se place \textbf{au début de la proposition} (avant le sujet), jamais entre le sujet et le verbe. Ex. : \zh{所以我们要…} (donc nous…), pas \zh{我们所以要…}.
\item \textbf{Particularité \zh{因为……所以……}} : en chinois, les deux marqueurs coexistent dans la phrase complexe ; en français, on ne garde qu'un seul connecteur (« Parce que c'est important, \textbf{nous} protégeons » ou « C'est important, \textbf{donc} nous protégeons »).
\end{enumerate}
\end{methode}

\begin{popfigure}
\begin{tikzpicture}[
  mindmap, grow cyclic, every node/.style=concept, concept color=popblue,
  level 1/.append style={level distance=4cm, sibling angle=72},
  level 2/.append style={level distance=3cm, sibling angle=40},
  text=white, font=\small
]
\node {Connecteurs logiques}
  child[concept color=popgreen] { node {Chronologie} child { node {首先} } child { node {然后} } }
  child[concept color=poporange] { node {Addition} child { node {还} } child { node {而且} } }
  child[concept color=poppurple] { node {Cause / Conséquence} child { node {因为} } child { node {所以} } child { node {因此} } }
  child[concept color=poppink] { node {Opposition} child { node {但是} } child { node {可是} } }
  child[concept color=popgold] { node {Condition} child { node {如果} } child { node {就} } };
\end{tikzpicture}
\legende{Carte mentale des connecteurs logiques (niveau HSK 2-3)}
\end{popfigure}

\begin{exoresolu}[Compléter avec le connecteur correct]
\textbf{Énoncé.} Complète avec : \zh{首先}, \zh{然后}, \zh{因为}, \zh{所以}, \zh{但是}, \zh{如果}.
\tcblower
\textbf{Solution détaillée.}

1. \zh{我们的学校很漂亮，\textbf{但是} 有时候不太干净。} (Opposition)
2. \zh{\textbf{首先} 我们要扫地，\textbf{然后} 要擦黑板。} (Chronologie)
3. \zh{树可以让空气更干净，\textbf{所以} 我们应该多种树。} (Conséquence)
4. \zh{\textbf{因为} 环境很重要，我们要保护它。} (Cause — \zh{所以} omis en 2\up{e} prop., correct)
5. \zh{\textbf{如果} 每个人都不乱扔垃圾，地面就干净了。} (Condition \zh{如果}…\zh{就}… + \zh{了} changement d'état)

\textbf{Conclusion :} analyser la logique (opposition, suite, cause, condition) avant de choisir ; placer le connecteur en tête de proposition.
\end{exoresolu}

\courssub{Écriture correcte des caractères du thème}

\begin{methode}[Traits, radicaux et discrimination graphique]
\begin{enumerate}
\item \textbf{Ordre des traits (règles d'or)} : haut $\rightarrow$ bas ; gauche $\rightarrow$ droite ; horizontal avant vertical ; extérieur avant intérieur ; fermer la boîte en dernier (ex. \zh{口}, \zh{国}).
\item \textbf{Radicaux (clés sémantiques) — Mémoriser par le sens} :
\end{enumerate}
\begin{center}
\begin{tabularx}{\linewidth}{|X|X|X|X|}
\hline
\rowcolor{popblueL} Caractère & Pinyin & Radical (Clé) & Sens du radical / Indice mnémotechnique \\
\hline
保 & bǎo & \zh{亻} (homme) & Un homme \zh{亻} qui veille \zh{呆} → protéger \\
\hline
护 & hù & \zh{扌} (main) & La main \zh{扌} qui garde \zh{戊} → protéger (action) \\
\hline
环 & huán & \zh{王} (jade) & Un jade \zh{王} en anneau → environnement (ce qui entoure) \\
\hline
境 & jìng & \zh{土} (terre) & La terre \zh{土} délimitée \zh{镸} → territoire, environnement \\
\hline
垃 & lā & \zh{土} (terre) & Terre \zh{土} + phonétique \zh{拉} → déchet (sol sale) \\
\hline
圾 & jī & \zh{土} (terre) & Terre \zh{土} + phonétique \zh{及} → déchet (suite de \zh{垃}) \\
\hline
树 & shù & \zh{木} (bois) & Bois \zh{木} + \zh{寸} → arbre (plante ligneuse) \\
\hline
桶 & tǒng & \zh{木} (bois) & Bois \zh{木} + \zh{同} → seau/poubelle (autrefois en bois) \\
\hline
扔 & rēng & \zh{扌} (main) & Main \zh{扌} + \zh{禺} → action de jeter \\
\hline
\end{tabularx}
\end{center}
\begin{enumerate}
\setcounter{enumi}{2}
\item \textbf{Caractères proches (pièges visuels)} — Différencier par le radical :
  \begin{itemize}
  \item \zh{境} (jìng, environnement, radical \zh{土}) $\neq$ \zh{镜} (jìng, miroir, radical \zh{钅} métal).
  \item \zh{保} (bǎo, protéger, radical \zh{亻}) $\neq$ \zh{呆} (dāi, bête, radical \zh{口} bouche).
  \item \zh{扔} (rēng, jeter, radical \zh{扌}) $\neq$ \zh{仍} (réng, encore, radical \zh{頁} tête).
  \item \zh{桶} (tǒng, seau, radical \zh{木}) $\neq$ \zh{筒} (tǒng, tube, radical \zh{竹} bambou).
  \end{itemize}
\item \textbf{Entraînement} : écrire chaque caractère 5 fois en énonçant l'ordre des traits (ex. \zh{环} : 1. horizontale haute, 2. verticale, 3. horizontale basse, 4. \zh{王}, 5. trait de liaison, 6. boucle finale).
\end{enumerate}
\end{methode}

\begin{exoresolu}[Radicaux et choix du bon caractère]
\textbf{Énoncé.} a) Donne le radical et le sens pour \zh{境}, \zh{护}, \zh{桶}. b) Choisis : \zh{我们应该保护学校的环（境 / 镜）。} ; \zh{不要（扔 / 仍）垃圾。} ; \zh{一个垃圾（桶 / 筒）。}
\tcblower
\textbf{Solution détaillée.}

a) \zh{境} : radical \zh{土} (terre) $\rightarrow$ « territoire, environnement » ; \zh{护} : radical \zh{扌} (main) $\rightarrow$ « protéger par action manuelle » ; \zh{桶} : radical \zh{木} (bois) $\rightarrow$ « récipient en bois $\rightarrow$ seau, poubelle ».

b) \zh{我们应该保护学校的环\textbf{境}。} (Choix \zh{境} : clé terre = lieu). \zh{不要\textbf{扔}垃圾。} (Choix \zh{扔} : clé main = action de jeter ; \zh{仍} = encore). \zh{一个垃圾\textbf{桶}。} (Choix \zh{桶} : clé bois = conteneur ; \zh{筒} = tube cylindrique, clé bambou).

\textbf{Conclusion :} le radical oriente le sens (terre/lieu, main/action, bois/objet) ; la prononciation identique ne suffit pas.
\end{exoresolu}

\courssub{Exercices de thème (français $\rightarrow$ chinois) et version (chinois $\rightarrow$ français)}

\begin{methode}[Réussir le thème : FR $\rightarrow$ CN]
\begin{enumerate}
\item \textbf{Segmenter} la phrase française : [Sujet] [Temps/Lieu] [Verbe] [Objet].
\item \textbf{Reclasser} selon l'ordre chinois canonique : \textbf{Sujet + Temps + Lieu (\zh{在} + lieu) + Verbe + Objet}. \emph{Le lieu précède le verbe.}
\item \textbf{Traiter les mots-outils} : \zh{的} (possession/attribut), \zh{应该} (devoir), \zh{把} (objet antéposé si manipulation/résultat), \zh{了} (aspect accompli), classificateurs (\zh{棵} arbre, \zh{个} général, \zh{本} livre).
\item \textbf{Règle \zh{两} vs \zh{二}} : \zh{两} devant classificateur/nombre (\zh{两棵树}) ; \zh{二} pour compter, numéros, décimales (\zh{二}, \zh{十二}, \zh{0.2}).
\item \textbf{Vérifier} les tons à voix haute.
\end{enumerate}
\end{methode}

\begin{exoresolu}[Thème guidé — 3 phrases]
\textbf{Énoncé.} Traduis en chinois :
1. Nous devons protéger l'environnement de notre école.
2. À l'école, les élèves ne doivent pas jeter les déchets par terre.
3. Il y a deux arbres devant la salle de classe.
\tcblower
\textbf{Solution détaillée.}

1. [Nous] [devons] [protéger] [notre école + \zh{的} + environnement] $\rightarrow$ \zh{我们应该保护我们学校的环境。} (Wǒmen yīnggāi bǎohù wǒmen xuéxiào de huánjìng.)

2. [À l'école = \zh{在学校里}] [élèves] [ne doivent pas] [jeter] [déchets] [par terre = \zh{在地上}].
   Objet \zh{垃圾} manipulé $\rightarrow$ construction \zh{把} : \zh{在学校里，学生不要把垃圾扔在地上。} (Zài xuéxiào lǐ, xuésheng bú yào bǎ lājī rēng zài dìshang.)

3. [Devant la classe = \zh{教室前面}] [il y a = \zh{有}] [deux = \zh{两}] [arbres = \zh{棵树}].
   Structure \zh{有} : Lieu + \zh{有} + Objet. $\rightarrow$ \zh{教室前面有两棵树。} (Jiàoshì qiánmian yǒu liǎng kē shù.) \emph{Attention : \zh{两} et non \zh{二}.}

\textbf{Pièges évités :} lieu avant verbe (phr. 2), construction \zh{把} pour manipulation (phr. 2), \zh{两} + classif. (phr. 3).
\end{exoresolu}

\begin{methode}[Réussir la version : CN $\rightarrow$ FR]
\begin{enumerate}
\item \textbf{Lecture globale} : saisir le sens général sans traduire mot à mot.
\item \textbf{Découpage en blocs} : souligner Sujet / Verbe / Compléments ; repérer \zh{的}, \zh{把}, \zh{应该}, \zh{因为……所以……}, \zh{如果……就……}.
\item \textbf{Réordonnancer} pour le français : remettre le lieu \textbf{après} le verbe (ex. \zh{在学校里学习} $\rightarrow$ « étudier \textbf{à l'école} »).
\item \textbf{Rendre \zh{把}} : supprimer la structure, mettre l'objet en COD direct (\zh{把垃圾扔} $\rightarrow$ « jeter les déchets »).
\item \textbf{Français naturel} : ne pas calquer \zh{因为……所以……} par « parce que… donc… » (lourd) ; choisir « Parce que…, … » ou « …, donc… ».
\end{enumerate}
\end{methode}

\begin{exoresolu}[Version guidée — Texte suivi]
\textbf{Énoncé.} Traduis en français :
\zh{因为空气和水对我们很重要，所以我们应该保护环境。首先，我们不要乱扔垃圾。然后，我们应该多种树，节约用水。如果每个人都努力，我们的城市会更干净、更漂亮。}
\tcblower
\textbf{Solution détaillée.}

\textbf{Découpage :}
(1) \zh{因为空气和水对我们很重要，所以我们应该保护环境。} — Cause (\zh{因为}…\zh{对我们} = pour nous) / Conséquence (\zh{所以}).
(2) \zh{首先，我们不要乱扔垃圾。} — Chronologie (\zh{首先}) ; \zh{乱扔} = jeter n'importe où / salir.
(3) \zh{然后，我们应该多种树，节约用水。} — Suite (\zh{然后}) ; deux VP coordonnés.
(4) \zh{如果每个人都努力，我们的城市会更干净、更漂亮。} — Condition (\zh{如果}…\zh{就}… omis) ; comparatif \zh{更} (plus).

\textbf{Traduction :}
« \textbf{Parce que} l'air et l'eau sont très importants \textbf{pour nous}, nous devons protéger l'environnement. \textbf{D'abord}, nous ne devons pas jeter les déchets n'importe où. \textbf{Ensuite}, nous devons planter plus d'arbres et économiser l'eau. \textbf{Si} chacun fait des efforts, notre ville sera \textbf{plus} propre et \textbf{plus} belle. »

\textbf{Note :} \zh{因为……所以……} rendu par « Parce que…, … » (pas de « donc » redondant). \zh{会更…} $\rightarrow$ futur + comparatif.
\end{exoresolu}

\courssub{Grille d'évaluation et auto-correction}

\begin{methode}[Grille d'auto-correction avant de rendre la copie]
Coche chaque critère après relecture :
\begin{enumerate}
\item \textbf{Contenu / Structure} : Introduction (constat) — Développement (2-3 actions avec \zh{应该}/\zh{不要}) — Conclusion (appel). $\square$
\item \textbf{Connecteurs} : au moins \textbf{deux} différents utilisés correctement (\zh{首先}, \zh{然后}, \zh{因为}, \zh{所以}, \zh{但是}, \zh{如果}). $\square$
\item \textbf{Structures obligatoires} : \zh{应该} (1x min), \zh{不要} (1x min), \zh{把} + objet (1x min si action sur objet). $\square$
\item \textbf{Ordre des mots} : Sujet + Temps + \zh{在} Lieu + Verbe + Objet (vérifier \textbf{chaque phrase}). $\square$
\item \textbf{Caractères cibles} : \zh{环境}, \zh{保护}, \zh{垃圾}, \zh{桶}, \zh{树}, \zh{种} — relire, compter les traits, vérifier le radical. $\square$
\item \textbf{Pinyin / Tons} : \zh{huán-jìng} (2-4), \zh{lā-jī} (1-1), \zh{yīng-gāi} (1-1), \zh{bǎo-hù} (3-4), \zh{zhòng-shù} (4-4), \zh{rēng} (1). Lire à voix haute. $\square$
\item \textbf{Forme} : 6 à 8 phrases minimum ; ponctuation chinoise (。 ，) ; pas d'espaces entre caractères. $\square$
\end{enumerate}
\end{methode}

\begin{exoresolu}[Correction d'un texte d'élève — 4 erreurs à trouver]
\textbf{Énoncé.} Le texte ci-dessous contient 4 erreurs (grammaire, vocabulaire, connectique, ponctuation). Corrige-le.

\zh{我们学校的环境很漂亮。我们应该保护它。我们不要扔垃圾在地上。我们应该多植树，因为树让空气更干净。所以我爱我们的学校。}
\tcblower
\textbf{Solution détaillée.}

\textbf{Erreur 1 (Ordre des mots / \zh{把}) :} \zh{我们不要扔垃圾在地上} $\rightarrow$ Incorrect. Complément de lieu résultatif après verbe de déplacement $\rightarrow$ structure \zh{把} obligatoire ou lieu avant verbe.
\emph{Correction :} \zh{我们不要\textbf{把}垃圾扔在地上。} (ou \zh{我们不要在地上扔垃圾。})

\textbf{Erreur 2 (Choix lexical) :} \zh{多植树} — \zh{植} (zhí) est formel (ex. \zh{植树节} Jour de l'arbre). À l'oral/écrit courant : \zh{种树} (zhòng shù).
\emph{Correction :} \zh{我们应该多\textbf{种树}。}

\textbf{Erreur 3 (Cohérence \zh{因为……所以……}) :} \zh{因为树让空气更干净。所以我爱我们的学校。} — Rupture logique. \zh{因为} appelle \zh{soi} dans la même phrase ou la suivante immédiate pour justifier l'action \zh{多种树}.
\emph{Correction :} Fusionner : \zh{\textbf{因为}树可以让空气更干净，\textbf{所以}我们应该多种树。}

\textbf{Erreur 4 (Connecteur final parasite) :} \zh{所以我爱我们的学校} — \zh{所以} (donc) mal employé en conclusion ; pas de cause antérieure pour « aimer l'école ».
\emph{Correction :} Supprimer \zh{所以} : \zh{我爱我们的学校。} (ou \zh{所以我们要爱护它。} si lien avec phrase 1).

\textbf{Texte corrigé (7 phrases) :}
\zh{我们学校的环境很漂亮。我们应该保护它。我们不要把垃圾扔在地上。因为树可以让空气更干净，所以我们应该多种树。我们还应该节约用水。如果大家都努力，学校会更干净。我爱我们的学校。}

\textbf{Bilan :} relecture systématique (grille) = détection ordre des mots, registres lexicaux, logique connecteurs, ponctuation.
\end{exoresolu}

\begin{attention}[Les 5 erreurs fatales au Bac / Évaluation — Check-list finale]
\begin{enumerate}
\item \textbf{Place du complément de lieu} : \textbf{Jamais} \zh{我们学习在学校} (calque FR). \textbf{Toujours} \zh{我们\textbf{在学校}学习。} (Sujet + \zh{在} Lieu + Verbe). Le lieu précède l'action.
\item \textbf{Confusion \zh{两} / \zh{二}} : Devant classificateur = \zh{两} (\zh{两棵树}, \zh{两个人}). \zh{二} seulement pour compter (\zh{一、二、三}), numéros (\zh{二楼}), maths. \zh{二棵树} = \textbf{faute grave}.
\item \textbf{Caractères « jumeaux » mal choisis} : \zh{境} (environnement, \zh{土}) $\neq$ \zh{镜} (miroir, \zh{钅}) ; \zh{扔} (jeter, \zh{扌}, 1\up{er} ton) $\neq$ \zh{仍} (encore, \zh{頁}, 2\up{e} ton) ; \zh{桶} (seau, \zh{木}) $\neq$ \zh{筒} (tube, \zh{竹}). 1 caractère faux = 0 pt sur le mot.
\item \textbf{Calque « Parce que… donc… » en version} : \zh{因为……所以……} $\rightarrow$ FR « Parce que…, … » \textbf{OU} « …, donc… » (un seul connecteur). Écrire « Parce que c'est important, donc nous protégeons » = \textbf{faute de français}.
\item \textbf{Tons approximatifs sur mots-outils fréquents} : \zh{yīnggāi} (1-1, pas 3-4), \zh{huánjìng} (2-4, pas 2-2), \zh{lājī} (1-1), \zh{bǎohù} (3-4), \zh{zhòngshù} (4-4). \textbf{Relire le pinyin à voix haute} avant de rendre la copie.
\end{enumerate}
\end{attention}

\newpage

\coursec{Exercices}

\courssub{Je vérifie mes connaissances}

\begin{exercice}[QCM : le bon choix]
Pour chaque question, entoure la bonne réponse.

\textbf{1.} « Protéger l'environnement » se dit :
\begin{itemize}
\item[a)] \zh{保护环境} (bǎohù huánjìng)
\item[b)] \zh{保护情况} (bǎohù qíngkuàng)
\item[c)] \zh{保险环境} (bǎoxiǎn huánjìng)
\end{itemize}

\textbf{2.} Pour exprimer une obligation morale ou un conseil (« tu dois, tu devrais »), on utilise :
\begin{itemize}
\item[a)] \zh{会} (huì)
\item[b)] \zh{应该} (yīnggāi)
\item[c)] \zh{可以} (kěyǐ)
\end{itemize}

\textbf{3.} Quelle phrase est correcte ?
\begin{itemize}
\item[a)] \zh{我们应该扔垃圾在地上。}
\item[b)] \zh{我们应该把垃圾扔进垃圾箱。}
\item[c)] \zh{我们应该垃圾箱把扔垃圾进。}
\end{itemize}

\textbf{4.} « D'abord..., ensuite..., enfin... » se dit :
\begin{itemize}
\item[a)] \zh{因为……所以……但是……}
\item[b)] \zh{首先……然后……最后……}
\item[c)] \zh{虽然……可是……而且……}
\end{itemize}
\end{exercice}

\begin{exercice}[Vrai ou faux ? Justifie ta réponse]
Dis si chaque affirmation est vraie (V) ou fausse (F), puis justifie en une phrase en français.

\textbf{1.} Le caractère \zh{环} contient le radical « roi / jade » \zh{王} à gauche.

\textbf{2.} Dans la phrase \zh{你应该保护学校的环境}, le mot \zh{学校} est un complément de lieu placé après le verbe.

\textbf{3.} \zh{垃圾} (lājī) signifie « déchets, ordures ».

\textbf{4.} On peut écrire \zh{保护环境很重要} pour dire « Protéger l'environnement est très important ».

\textbf{5.} Le connecteur \zh{所以} (suǒyǐ) introduit une cause.
\end{exercice}

\begin{exercice}[Vocabulaire : complète le tableau]
Complète le tableau suivant : retrouve le caractère, le pinyin ou la traduction manquante.

\begin{center}
\begin{tabularx}{\linewidth}{|X|X|X|X|}
\hline
\rowcolor{popblueL} Caractères & Pinyin & Nature & Traduction \\
\hline
环境 & \ldots\ldots\ldots & nom & \ldots\ldots\ldots \\
\hline
\ldots\ldots\ldots & bǎohù & verbe & \ldots\ldots\ldots \\
\hline
垃圾 & \ldots\ldots\ldots & \ldots\ldots\ldots & déchets, ordures \\
\hline
\ldots\ldots\ldots & gānjìng & adjectif & \ldots\ldots\ldots \\
\hline
节约 & \ldots\ldots\ldots & verbe & \ldots\ldots\ldots \\
\hline
\ldots\ldots\ldots & zhòng shù & \ldots\ldots\ldots & planter des arbres \\
\hline
\end{tabularx}
\end{center}
\end{exercice}

\begin{exercice}[Associe le pinyin aux caractères]
Relie chaque caractère de la colonne A à son pinyin de la colonne B, puis écris la traduction française de chaque caractère.

\textbf{Colonne A :} \zh{环} \quad \zh{境} \quad \zh{保} \quad \zh{护} \quad \zh{校} \quad \zh{圾} \quad \zh{净} \quad \zh{约}

\textbf{Colonne B :} jīng \quad hù \quad jī \quad huán \quad yuē \quad bǎo \quad jìng \quad xiào
\end{exercice}

\courssub{Je m'entraîne}

\begin{methode}[Traduire du français vers le chinois (thème) : la démarche en 4 étapes]
\begin{enumerate}
\item \textbf{Repérer la structure} : qui fait quoi ? (sujet + verbe + objet). Y a-t-il un mot modal (\zh{应该}, \zh{可以}, \zh{要}) ? Il se place \emph{avant} le verbe.
\item \textbf{Choisir la construction} : l'objet est-il déplacé et « traité » ? Alors utilise \zh{把} : Sujet + \zh{把} + objet + verbe + résultat/lieu.
\item \textbf{Placer les compléments} : le temps et le lieu se mettent avant le verbe ; \zh{的} relie un nom à son complément (\zh{学校的环境}).
\item \textbf{Vérifier} : tons du pinyin, classificateur si numéral (\zh{三棵树}), ponctuation chinoise pleine chasse （，。！）.
\end{enumerate}
\end{methode}

\begin{exercice}[Traduction : thème (français → chinois)]
Traduis les phrases suivantes en chinois (caractères + pinyin).

\textbf{1.} Nous devons protéger l'environnement de notre école.

\textbf{2.} Ne jette pas les déchets par terre !

\textbf{3.} D'abord, nous économisons l'eau ; ensuite, nous plantons des arbres.

\textbf{4.} Garder l'école propre, c'est la responsabilité de chacun.
\end{exercice}

\begin{exercice}[Traduction : version (chinois → français)]
Traduis les phrases suivantes en français.

\textbf{1.} \zh{学校的环境很干净，我们都很高兴。}\\
\emph{Xuéxiào de huánjìng hěn gānjìng, wǒmen dōu hěn gāoxìng.}

\textbf{2.} \zh{为了保护环境，我们不应该浪费水。}\\
\emph{Wèile bǎohù huánjìng, wǒmen bù yīnggāi làngfèi shuǐ.}

\textbf{3.} \zh{每天下午，同学们一起打扫教室。}\\
\emph{Měi tiān xiàwǔ, tóngxuémen yìqǐ dǎsǎo jiàoshì.}
\end{exercice}

\begin{exercice}[Texte à trous : les mots-outils]
Complète le texte suivant avec les mots de la liste : \zh{应该} \quad \zh{把} \quad \zh{的} \quad \zh{了} \quad \zh{因为} \quad \zh{所以}

\begin{textebox}[L'environnement de notre école]
我们学校（1）环境不太干净。（2）有些同学扔垃圾在地上，（3）我们（4）一起做点儿什么。昨天，我们（5）垃圾放进垃圾箱，还种（6）三棵小树。现在学校漂亮多了！\\
\emph{Wǒmen xuéxiào (1) huánjìng bú tài gānjìng. (2) yǒuxiē tóngxué rēng lājī zài dìshang, (3) wǒmen (4) yìqǐ zuò diǎnr shénme. Zuótiān, wǒmen (5) lājī fàng jìn lājī xiāng, hái zhòng (6) sān kē xiǎo shù. Xiànzài xuéxiào piàoliang duō le !}\\
« L'environnement de notre école n'est pas très propre. Comme certains élèves jettent des déchets par terre, nous devons faire quelque chose ensemble. Hier, nous avons mis les déchets dans la poubelle et nous avons aussi planté trois petits arbres. Maintenant, l'école est bien plus belle ! »
\end{textebox}
\end{exercice}

\begin{exercice}[Remets les mots dans l'ordre]
Remets les mots dans le bon ordre pour former des phrases correctes, puis traduis-les en français.

\textbf{1.} \zh{保护 / 我们 / 环境 / 应该 / 学校 / 的}

\textbf{2.} \zh{垃圾 / 把 / 请 / 垃圾箱 / 扔进 / 你}

\textbf{3.} \zh{首先 / 水 / 节约 / 我们 / ， / 然后 / 树 / 种}

\textbf{4.} \zh{很 / 环境 / 重要 / 保护 / 是 / 的 / 事 / 一 / 件}
\end{exercice}

\courssub{Je me prépare au Bac}

\begin{exercice}[Compréhension de l'écrit et questions de langue]
Lis le texte suivant, puis réponds aux questions.

\begin{textebox}[Une école verte à Yaoundé]
雅温得的一所学校最近有了大变化。以前，校园里有很多垃圾，水也常常被浪费。今年，老师和学生一起想了很多办法：首先，每个教室外面都放了两个垃圾箱，一个放纸，一个放别的东西；然后，大家种了二十棵树；最后，学校还告诉学生要节约用水、用电。现在，校园又干净又漂亮，学生们学习得也更开心了。校长说：“保护环境是每个人的责任。”\\
\emph{Yǎwēndé de yí suǒ xuéxiào zuìjìn yǒu le dà biànhuà. Yǐqián, xiàoyuán li yǒu hěn duō lājī, shuǐ yě chángcháng bèi làngfèi. Jīnnián, lǎoshī hé xuésheng yìqǐ xiǎng le hěn duō bànfǎ : shǒuxiān, měi ge jiàoshì wàimiàn dōu fàng le liǎng ge lājī xiāng, yí ge fàng zhǐ, yí ge fàng bié de dōngxi ; ránhòu, dàjiā zhòng le èrshí kē shù ; zuìhòu, xuéxiào hái gàosu xuésheng yào jiéyuē yòng shuǐ, yòng diàn. Xiànzài, xiàoyuán yòu gānjìng yòu piàoliang, xuéshengmen xuéxí de yě gèng kāixīn le. Xiàozhǎng shuō : « Bǎohù huánjìng shì měi ge rén de zérèn. »}\\
« Une école de Yaoundé a beaucoup changé récemment. Avant, il y avait beaucoup de déchets dans la cour, et l'eau était souvent gaspillée. Cette année, professeurs et élèves ont trouvé ensemble de nombreuses solutions : d'abord, deux poubelles ont été placées devant chaque salle, une pour le papier, une pour le reste ; ensuite, tout le monde a planté vingt arbres ; enfin, l'école a aussi demandé aux élèves d'économiser l'eau et l'électricité. Maintenant, la cour est à la fois propre et belle, et les élèves étudient avec encore plus de joie. Le directeur dit : "Protéger l'environnement, c'est la responsabilité de chacun." »
\end{textebox}

\textbf{A. Questions de compréhension} (réponds en français) :

\textbf{1.} Quels problèmes l'école avait-elle avant ?

\textbf{2.} Cite les trois mesures prises par l'école cette année.

\textbf{3.} Quel est le résultat de ces mesures ?

\textbf{4.} Que signifie la phrase du directeur ? Es-tu d'accord ? (2 ou 3 phrases)

\textbf{B. Questions de langue} :

\textbf{5.} Releve dans le texte les trois connecteurs d'étapes et traduis-les.

\textbf{6.} Explique la structure de \zh{学习得也更开心了} : quel mot-outil introduit le complément de manière ?

\textbf{7.} Dans \zh{水也常常被浪费}, quel mot marque le passif ? Traduis la phrase.
\end{exercice}

\begin{exercice}[Thème et version type Bac]
\textbf{A. Thème.} Traduis en chinois (caractères + pinyin) :

\textbf{1.} Chaque élève doit jeter les déchets dans la poubelle.

\textbf{2.} Si nous économisons l'eau, nous protégeons aussi l'environnement.

\textbf{3.} Notre école est devenue plus propre et plus belle.

\textbf{B. Version.} Traduis en français le texte suivant (environ 60 caractères). Tu peux t'aider des notes de vocabulaire.

\begin{textebox}[Le tri des déchets]
在中国，很多城市开始垃圾分类。人们把垃圾分成四种：可以回收的、厨房的垃圾、有害的垃圾和其他的垃圾。这样做可以保护环境，也可以节约资源。\\
\emph{Zài Zhōngguó, hěn duō chéngshì kāishǐ lājī fēnlèi. Rénmen bǎ lājī fēn chéng sì zhǒng : kěyǐ huíshōu de, chúfáng de lājī, yǒuhài de lājī hé qítā de lājī. Zhèyàng zuò kěyǐ bǎohù huánjìng, yě kěyǐ jiéyuē zīyuán.}
\end{textebox}

\textbf{Notes de vocabulaire :} \zh{分类} (fēnlèi) : trier, classer ; \zh{回收} (huíshōu) : recycler ; \zh{有害} (yǒuhài) : nuisible, dangereux ; \zh{资源} (zīyuán) : ressources.
\end{exercice}

\begin{methode}[Rédiger une lettre de proposition en chinois : 5 étapes]
\begin{enumerate}
\item \textbf{Salutation} : \zh{校长先生，您好！} (Xiàozhǎng xiānsheng, nín hǎo !) — utilise \zh{您}, le « vous » de politesse.
\item \textbf{Constat} : une phrase sur la situation (\zh{我们学校的环境不太干净。}).
\item \textbf{Propositions} : trois mesures reliées par \zh{首先……然后……最后……}, avec \zh{应该} ou \zh{可以} et au moins une structure en \zh{把}.
\item \textbf{Conclusion} : une phrase de synthèse (\zh{这样，我们的学校会更干净、更漂亮。}).
\item \textbf{Formule finale} : \zh{谢谢您！} puis signature simple (\zh{学生：李明}).
\end{enumerate}
\end{methode}

\begin{exercice}[Expression écrite : lettre au directeur]
\textbf{Sujet.} Tu es élève en Première A au lycée de Douala. Rédige en chinois une lettre au directeur de ton école pour proposer \textbf{trois mesures concrètes} de protection de l'environnement scolaire (par exemple : poubelles de tri, plantation d'arbres, économie d'eau et d'électricité, nettoyage de la cour).

\textbf{Contraintes :}
\begin{itemize}
\item Longueur : 80 à 100 caractères chinois.
\item Utilise au moins une fois les structures suivantes : \zh{应该}, \zh{把}, \zh{首先……然后……最后……}.
\item Commence par \zh{校长先生，您好！} (Xiàozhǎng xiānsheng, nín hǎo !) et termine par une formule de politesse simple, par exemple \zh{谢谢您！} (Xièxie nín !).
\item Soigne l'écriture des caractères du thème : \zh{环}, \zh{境}, \zh{保}, \zh{护}, \zh{学}, \zh{校}, \zh{垃}, \zh{圾}.
\end{itemize}

\textbf{Aide au démarrage :} tu peux suivre ce plan : salutation → constat (l'école n'est pas assez propre) → trois propositions avec connecteurs → conclusion et remerciement.
\end{exercice}

\begin{exoresolu}[Modèle de lettre rédigée et commentée]
Rédige la lettre demandée dans l'exercice précédent, en respectant toutes les contraintes.
\tcblower
\textbf{Modèle de production (environ 95 caractères) :}

\zh{校长先生，您好！我们学校的环境不太干净，我们应该保护它。首先，请同学们把垃圾扔进垃圾箱，不要扔在地上；然后，我们可以一起种树、种花；最后，大家应该节约用水、用电。这样，我们的学校会又干净又漂亮。谢谢您！学生：李明}

\emph{Xiàozhǎng xiānsheng, nín hǎo ! Wǒmen xuéxiào de huánjìng bú tài gānjìng, wǒmen yīnggāi bǎohù tā. Shǒuxiān, qǐng tóngxuémen bǎ lājī rēng jìn lājī xiāng, bú yào rēng zài dìshang ; ránhòu, wǒmen kěyǐ yìqǐ zhòng shù, zhòng huā ; zuìhòu, dàjiā yīnggāi jiéyuē yòng shuǐ, yòng diàn. Zhèyàng, wǒmen de xuéxiào huì yòu gānjìng yòu piàoliang. Xièxie nín ! Xuésheng : Lǐ Míng.}

« Monsieur le Directeur, bonjour ! L'environnement de notre école n'est pas très propre, nous devons le protéger. D'abord, que les camarades jettent les déchets dans la poubelle et non par terre ; ensuite, nous pouvons planter ensemble des arbres et des fleurs ; enfin, tout le monde doit économiser l'eau et l'électricité. Ainsi, notre école sera à la fois propre et belle. Merci ! L'élève : Li Ming. »

\textbf{Points de réussite :} salutation avec \zh{您} ; constat bref ; trois connecteurs d'étapes ; une structure en \zh{把} ; deux emplois de \zh{应该} ; conclusion avec \zh{这样} et \zh{又……又……} ; formule de politesse finale.
\end{exoresolu}

\begin{attention}[Erreurs fréquentes des francophones dans cette production]
\begin{itemize}
\item \textbf{Oublier \zh{的}} : on écrit \zh{学校的环境} (l'environnement \emph{de} l'école), jamais \zh{学校环境} dans une rédaction soignée à ce niveau.
\item \textbf{Mal placer le modal} : \zh{应该} se met \emph{avant} le verbe : \zh{我们应该保护环境}, pas après.
\item \textbf{Structure en \zh{把} incomplète} : après \zh{把} + objet, le verbe ne peut pas rester seul : il faut un complément (\zh{把垃圾扔进垃圾箱}, \zh{把垃圾放好}).
\item \textbf{Confondre les trois « de »} : \zh{的} (nom), \zh{地} (adverbe avant verbe), \zh{得} (complément après verbe : \zh{学习得很开心}).
\item \textbf{Classificateur} : \zh{三棵树} (arbres : \zh{棵}), \zh{一件事} (chose : \zh{件}), \zh{两个垃圾箱} (poubelles : \zh{个}).
\item \textbf{Tons} : \zh{环境} huánjìng (2\ieme{}--4\ieme{}), \zh{保护} bǎohù (3\ieme{}--4\ieme{}), \zh{垃圾} lājī (1\ier{}--1\ier{}), \zh{干净} gānjìng (1\ier{}--4\ieme{}).
\end{itemize}
\end{attention}

\begin{savaistu}[Le tri des déchets en Chine et au Cameroun]
Depuis quelques années, de grandes villes chinoises comme Shanghai ont rendu le tri des déchets obligatoire : on sépare les déchets recyclables, les déchets de cuisine, les déchets dangereux et les autres déchets, avec des poubelles de couleurs différentes. Au Cameroun, à Douala et à Yaoundé, la collecte des ordures est assurée par des entreprises spécialisées, et de nombreuses associations de jeunes organisent des journées de nettoyage des quartiers et des écoles. Dans les deux pays, l'école joue un rôle clé : c'est souvent là que naissent les bonnes habitudes, comme le dit le proverbe chinois \zh{保护环境，人人有责} (bǎohù huánjìng, rén rén yǒu zé) : « Protéger l'environnement, chacun en a la responsabilité. »
\end{savaistu}

\newpage

\coursec{Corrigés des exercices}

\courssub{Corrigés détaillés}

\begin{corrige}[Exercice 1 — QCM : le bon choix]
\textbf{1. a)} \zh{保护环境} (bǎohù huánjìng) = protéger l'environnement. \zh{情况} (qíngkuàng) = situation ; \zh{保险} (bǎoxiǎn) = assurance. Faux ami sonore : \zh{保} seul ne suffit pas à dire « assurer ».

\textbf{2. b)} \zh{应该} (yīnggāi) = devoir, conseil moral. \zh{会} (huì) = savoir faire / futur probable ; \zh{可以} (kěyǐ) = permission, possibilité.

\textbf{3. b)} \zh{我们应该把垃圾扔进垃圾箱。} Structure en \zh{把} : Sujet + \zh{把} + objet + verbe + complément de direction (\zh{进}). La phrase a) est un contresens (jeter par terre !) ; la c) a un ordre des mots impossible : \zh{把} précède toujours l'objet, jamais le verbe.

\textbf{4. b)} \zh{首先……然后……最后……} (shǒuxiān... ránhòu... zuìhòu...) = d'abord, ensuite, enfin. La série a) mélange cause et opposition ; la c) mélange concession et addition.

\textbf{Points de vigilance :} au QCM, lis toujours les trois propositions en entier ; beaucoup de pièges reposent sur des caractères proches (\zh{境} jìng vs \zh{情} qíng) ou sur l'ordre des mots.
\end{corrige}

\begin{corrige}[Exercice 2 — Vrai ou faux ?]
\textbf{1. V.} \zh{环} = radical \zh{王} (clé du jade, à gauche) + \zh{不} (support phonétique). Dans les caractères, \zh{王} à gauche est presque toujours la clé du jade (\zh{玉}), non du « roi ».

\textbf{2. F.} \zh{学校} n'est pas un complément de lieu : grâce à \zh{的}, c'est un complément du nom : \zh{学校的环境} = « l'environnement \emph{de l'école} ». Tout le groupe \zh{学校的环境} est l'objet du verbe \zh{保护}.

\textbf{3. V.} \zh{垃圾} (lājī) = déchets, ordures ; les deux caractères portent le radical de la terre \zh{土}.

\textbf{4. V.} Un groupe verbal peut être sujet : \zh{保护环境很重要} (Bǎohù huánjìng hěn zhòngyào) = « Protéger l'environnement est très important ». Pas besoin de traduire « c'est ».

\textbf{5. F.} \zh{所以} (suǒyǐ) introduit la \emph{conséquence} (« donc ») ; la cause est introduite par \zh{因为} (yīnwèi, « parce que »). La paire complète est \zh{因为……所以……}.
\end{corrige}

\begin{corrige}[Exercice 3 — Vocabulaire : complète le tableau]
\begin{center}
\begin{tabularx}{\linewidth}{|X|X|X|X|}
\hline
\rowcolor{popblueL} Caractères & Pinyin & Nature & Traduction \\
\hline
环境 & huánjìng & nom & environnement \\
\hline
保护 & bǎohù & verbe & protéger \\
\hline
垃圾 & lājī & nom & déchets, ordures \\
\hline
干净 & gānjìng & adjectif & propre \\
\hline
节约 & jiéyuē & verbe & économiser \\
\hline
种树 & zhòng shù & verbe + objet & planter des arbres \\
\hline
\end{tabularx}
\end{center}

\textbf{Points de vigilance :} \zh{种树} est un verbe à objet séparable (离合词) : on peut dire \zh{种了三棵树}. Attention aux tons : \zh{节约} jiéyuē (2e ton -- 1er ton), \zh{干净} gānjìng (1er ton -- 4e ton).
\end{corrige}

\begin{corrige}[Exercice 4 — Associe le pinyin aux caractères]
\zh{环} huán (entourer, anneau) ; \zh{境} jìng (territoire, cadre) ; \zh{保} bǎo (protéger, garder) ; \zh{护} hù (protéger) ; \zh{校} xiào (école) ; \zh{圾} jī (dans \zh{垃圾} lājī, déchets) ; \zh{净} jìng (propre) ; \zh{约} yuē (restreindre, dans \zh{节约} jiéyuē, économiser).

\textbf{Homophones à retenir :} \zh{境} jìng et \zh{净} jìng se prononcent exactement pareil (4e ton) mais s'écrivent différemment : \zh{境} porte le radical de la terre \zh{土}, \zh{净} le radical de la glace \zh{冫}. Seul le contexte (et l'écriture !) les distingue.
\end{corrige}

\begin{corrige}[Exercice 5 — Traduction : thème (français → chinois)]
\textbf{1.} \zh{我们应该保护我们学校的环境。} \emph{Wǒmen yīnggāi bǎohù wǒmen xuéxiào de huánjìng.} Modal \zh{应该} avant le verbe ; deux groupes nominaux imbriqués reliés par \zh{的}.

\textbf{2.} \zh{不要把垃圾扔在地上！} \emph{Bú yào bǎ lājī rēng zài dìshang !} Interdiction : \zh{不要} + structure en \zh{把} ; le complément de lieu \zh{在地上} suit le verbe \zh{扔}.

\textbf{3.} \zh{首先，我们节约用水；然后，我们种树。} \emph{Shǒuxiān, wǒmen jiéyuē yòng shuǐ; ránhòu, wǒmen zhòng shù.} Connecteurs d'étapes en tête de proposition, suivis de la virgule chinoise.

\textbf{4.} \zh{保持学校干净是每个人的责任。} \emph{Bǎochí xuéxiào gānjìng shì měi ge rén de zérèn.} Le groupe verbal \zh{保持学校干净} fait office de sujet ; \zh{每个人} = chacun.

\textbf{Barème indicatif (4 points) :} 1 point par phrase — 0,5 pour la structure correcte, 0,5 pour les caractères et le pinyin exacts.
\end{corrige}

\begin{corrige}[Exercice 6 — Traduction : version (chinois → français)]
\textbf{1.} « L'environnement de l'école est très propre, nous sommes tous très contents. » \zh{都} (dōu) = tous, placé avant l'adjectif ou le verbe.

\textbf{2.} « Pour protéger l'environnement, nous ne devons pas gaspiller l'eau. » \zh{为了} (wèile) = pour, afin de (+ but) ; \zh{浪费} (làngfèi) = gaspiller ; négation du modal : \zh{不应该}.

\textbf{3.} « Chaque après-midi, les camarades nettoient la salle de classe ensemble. » Le complément de temps \zh{每天下午} est en tête de phrase ; \zh{一起} (yìqǐ) = ensemble, avant le verbe.

\textbf{Points de vigilance :} en version, restitue les connecteurs et l'ordre logique en français naturel ; ne traduis pas mot à mot (\zh{打扫} = nettoyer, pas « frapper-balayer »).
\end{corrige}

\begin{corrige}[Exercice 7 — Texte à trous : les mots-outils]
(1) \zh{的} : complément du nom — « l'environnement \emph{de} notre école ».\\
(2) \zh{因为} : cause — « comme certains élèves... ».\\
(3) \zh{所以} : conséquence — paire \zh{因为……所以……}.\\
(4) \zh{应该} : modal de devoir, placé avant le verbe \zh{做}.\\
(5) \zh{把} : construction de disposition — \zh{把垃圾放进垃圾箱} = mettre les déchets dans la poubelle.\\
(6) \zh{了} : accompli après le verbe \zh{种}, action réalisée (\zh{昨天} = hier).

Texte complété : \zh{我们学校的环境不太干净。因为有些同学扔垃圾在地上，所以我们应该一起做点儿什么。昨天，我们把垃圾放进垃圾箱，还种了三棵小树。现在学校漂亮多了！}
\end{corrige}

\begin{corrige}[Exercice 8 — Remets les mots dans l'ordre]
\textbf{1.} \zh{我们应该保护学校的环境。} \emph{Wǒmen yīnggāi bǎohù xuéxiào de huánjìng.} « Nous devons protéger l'environnement de l'école. »

\textbf{2.} \zh{请你把垃圾扔进垃圾箱。} \emph{Qǐng nǐ bǎ lājī rēng jìn lājī xiāng.} « S'il te plaît, jette les déchets dans la poubelle. » \zh{请} en tête pour la politesse ; \zh{把} + objet + verbe de déplacement + \zh{进}.

\textbf{3.} \zh{首先，我们节约水，然后种树。} \emph{Shǒuxiān, wǒmen jiéyuē shuǐ, ránhòu zhòng shù.} « D'abord, nous économisons l'eau, ensuite nous plantons des arbres. »

\textbf{4.} \zh{保护环境是一件很重要的事。} \emph{Bǎohù huánjìng shì yí jiàn hěn zhòngyào de shì.} « Protéger l'environnement est une chose très importante. » Classificateur \zh{件} pour \zh{事} ; sandhi : \zh{一} se prononce \emph{yí} (2e ton) devant un 4e ton.
\end{corrige}

\begin{corrige}[Exercice 9 — Compréhension de l'écrit et questions de langue]
\textbf{A. 1.} Avant, il y avait beaucoup de déchets dans la cour et l'eau était souvent gaspillée.

\textbf{2.} Trois mesures : (a) deux poubelles devant chaque salle (une pour le papier, une pour le reste) ; (b) la plantation de vingt arbres ; (c) la sensibilisation des élèves à économiser l'eau et l'électricité.

\textbf{3.} La cour est maintenant à la fois propre et belle, et les élèves étudient avec encore plus de joie.

\textbf{4.} Le directeur affirme que protéger l'environnement est la responsabilité de chacun : tout le monde doit participer, pas seulement les adultes ou les autorités. Réponse personnelle argumentée attendue, par exemple : « Je suis d'accord, car si chaque élève jette ses déchets dans la poubelle, l'école reste propre. »

\textbf{B. 5.} \zh{首先} (shǒuxiān) = d'abord ; \zh{然后} (ránhòu) = ensuite ; \zh{最后} (zuìhòu) = enfin.

\textbf{6.} Le mot-outil \zh{得} (de), placé après le verbe, introduit le complément de manière ou de degré : \zh{学习得也更开心} = « étudier de façon encore plus joyeuse ». À ne pas confondre avec \zh{的} (complément du nom) ni \zh{地} (adverbe avant le verbe).

\textbf{7.} Le passif est marqué par \zh{被} (bèi) : patient + \zh{被} + (agent) + verbe. Traduction : « l'eau était aussi souvent gaspillée ».

\textbf{Barème indicatif (10 points) :} A : 1 + 2 + 1 + 2 = 6 points ; B : 1,5 + 1,5 + 1 = 4 points.
\end{corrige}

\begin{corrige}[Exercice 10 — Thème et version type Bac]
\textbf{A. Thème.}

\textbf{1.} \zh{每个同学都应该把垃圾扔进垃圾箱。} \emph{Měi ge tóngxué dōu yīnggāi bǎ lājī rēng jìn lājī xiāng.} Après \zh{每个}, on ajoute souvent \zh{都} avant le verbe ; structure en \zh{把} obligatoire ici car l'objet est « traité ».

\textbf{2.} \zh{如果我们节约用水，我们也保护环境。} \emph{Rúguǒ wǒmen jiéyuē yòng shuǐ, wǒmen yě bǎohù huánjìng.} Hypothèse : \zh{如果……，……也/就……}.

\textbf{3.} \zh{我们的学校变得更干净、更漂亮了。} \emph{Wǒmen de xuéxiào biàn de gèng gānjìng, gèng piàoliang le.} \zh{变得} = devenir + adjectif ; \zh{了} final marque le changement d'état.

\textbf{B. Version.} « En Chine, de nombreuses villes ont commencé le tri des déchets. Les gens divisent les déchets en quatre catégories : ceux qui peuvent être recyclés, les déchets de cuisine, les déchets dangereux et les autres déchets. Agir ainsi permet de protéger l'environnement et aussi d'économiser les ressources. »

Points de langue : \zh{把垃圾分成四种} = structure en \zh{把} + \zh{分成} (diviser en) ; \zh{这样做} = « faire ainsi » ; \zh{可以……也可以……} = « permet de... et aussi de... ».

\textbf{Barème indicatif (10 points) :} thème 3 × 1,5 = 4,5 points (structure 0,75 + caractères/pinyin 0,75) ; version 5,5 points (fidélité du sens 3,5, qualité du français 2).

\textbf{Points de vigilance :} ne pas oublier \zh{都} après \zh{每个} ; ne pas laisser le verbe seul après \zh{把} + objet ; en version, \zh{厨房的垃圾} se traduit naturellement par « déchets de cuisine ».
\end{corrige}

\begin{corrige}[Exercice 11 — Expression écrite : lettre au directeur]
\textbf{Proposition de rédaction modèle (environ 95 caractères) :}

\zh{校长先生，您好！我们学校的环境不太干净，我们应该保护它。首先，请同学们把垃圾扔进垃圾箱，不要扔在地上；然后，我们可以一起种树、种花；最后，大家应该节约用水、用电。这样，我们的学校会又干净又漂亮。谢谢您！学生：李明}

\emph{Xiàozhǎng xiānsheng, nín hǎo! Wǒmen xuéxiào de huánjìng bú tài gānjìng, wǒmen yīnggāi bǎohù tā. Shǒuxiān, qǐng tóngxuémen bǎ lājī rēng jìn lājī xiāng, bú yào rēng zài dìshang; ránhòu, wǒmen kěyǐ yìqǐ zhòng shù, zhòng huā; zuìhòu, dàjiā yīnggāi jiéyuē yòng shuǐ, yòng diàn. Zhèyàng, wǒmen de xuéxiào huì yòu gānjìng yòu piàoliang. Xièxie nín! Xuésheng: Lǐ Míng.}

« Monsieur le Directeur, bonjour ! L'environnement de notre école n'est pas très propre, nous devons le protéger. D'abord, que les camarades jettent les déchets dans la poubelle et non par terre ; ensuite, nous pouvons planter ensemble des arbres et des fleurs ; enfin, tout le monde doit économiser l'eau et l'électricité. Ainsi, notre école sera à la fois propre et belle. Merci ! L'élève : Li Ming. »

\textbf{Barème indicatif (10 points) :} respect des contraintes de forme (salutation, formule finale, longueur) : 2 points ; contenu (constat + trois mesures concrètes) : 3 points ; structures imposées (\zh{应该}, \zh{把}, \zh{首先……然后……最后……}) : 3 points ; exactitude des caractères et de la grammaire : 2 points.

\textbf{Points de vigilance :}
\begin{itemize}
\item Salutation avec \zh{您} (vouvoiement de politesse), jamais \zh{你} pour s'adresser au directeur.
\item \zh{应该} toujours \emph{avant} le verbe ; après \zh{把} + objet, le verbe porte un complément (\zh{扔进垃圾箱}, \zh{放好}).
\item Ne pas oublier \zh{的} : \zh{学校的环境}, \zh{我们的学校}.
\item Classificateurs : \zh{三棵树} (\zh{棵} pour les arbres), \zh{两个垃圾箱} (\zh{个}).
\item Conclusion avec \zh{这样} (« ainsi ») et la structure \zh{又……又……} (« à la fois... et... »).
\end{itemize}
\end{corrige}

\newpage

\coursec{Fiche bilan}

\begin{popfigure}
\begin{tikzpicture}[mindmap, grow cyclic, every node/.style={concept, align=center, font=\small}, concept color=popgreen, level 1 concept/.append style={level distance=3.4cm, sibling angle=51.4}, root concept/.append style={concept color=popgreen, font=\bfseries}]
\node[align=center]{保护环境\\ \emph{bǎohù huánjìng}\\ Protéger l'environnement}
  child[concept color=popblue]{ node[align=center]{Lexique\\ 垃圾 回收\\ 节约 污染} }
  child[concept color=poporange]{ node[align=center]{Modaux\\ 应该 / 不应该\\ 要 / 不要} }
  child[concept color=poppurple]{ node[align=center]{Connecteurs\\ 首先 其次 最后\\ 因为 所以 如果} }
  child[concept color=popgold]{ node[align=center]{Structure\\ du texte\\ 开头-正文-结尾} }
  child[concept color=poppink]{ node[align=center]{Caractères\\ 环 境 垃 圾\\ 保 护 节 约} }
  child[concept color=popgreen!70]{ node[align=center]{Traduction\\ thème / version\\ 我们应该…} }
  child[concept color=popblue!60]{ node[align=center]{Méthode\\ brouillon\\ relecture} };
\end{tikzpicture}
\legende{Carte mentale de la leçon : écrire un texte sur la protection de l'environnement de l'école.}
\end{popfigure}

\begin{bilan}
\textbf{Lexique clé du thème (à connaître par cœur)}

\begin{center}
\begin{tabularx}{\linewidth}{|X|X|X|X|}
\hline
\rowcolor{popblueL} Caractères & Pinyin & Nature & Traduction \\
\hline
环境 & huánjìng & n. & environnement \\
\hline
保护 & bǎohù & v. & protéger \\
\hline
垃圾 & lājī & n. & déchets, ordures \\
\hline
回收 & huíshōu & v. & recycler \\
\hline
节约 & jiéyuē & v. & économiser \\
\hline
水 / 电 & shuǐ / diàn & n. & eau / électricité \\
\hline
污染 & wūrǎn & v./n. & polluer ; pollution \\
\hline
干净 / 脏 & gānjìng / zāng & adj. & propre / sale \\
\hline
种树 & zhòng shù & v. compl. & planter des arbres \\
\hline
扔 & rēng & v. & jeter \\
\hline
\end{tabularx}
\end{center}

\textbf{Grammaire à connaître par cœur}

\begin{center}
\begin{tabularx}{\linewidth}{|X|X|X|}
\hline
\rowcolor{popblueL} Structure & Exemple & Traduction \\
\hline
Sujet + 应该 + verbe & 我们应该保护学校的环境。Wǒmen yīnggāi bǎohù xuéxiào de huánjìng. & Nous devons protéger l'environnement de l'école. \\
\hline
Sujet + 不应该 + verbe & 我们不应该乱扔垃圾。Wǒmen bù yīnggāi luàn rēng lājī. & Nous ne devons pas jeter les déchets n'importe où. \\
\hline
要 / 不要 + verbe & 不要浪费水。Bú yào làngfèi shuǐ. & Ne gaspille pas l'eau. \\
\hline
因为 … 所以 … & 因为环境很重要，所以我们要保护它。Yīnwèi huánjìng hěn zhòngyào, suǒyǐ wǒmen yào bǎohù tā. & Parce que l'environnement est important, nous devons le protéger. \\
\hline
如果 … 就 … & 如果每个人都努力，学校就会更干净。Rúguǒ měi ge rén dōu nǔlì, xuéxiào jiù huì gèng gānjìng. & Si chacun fait des efforts, l'école sera plus propre. \\
\hline
\end{tabularx}
\end{center}

\textbf{Méthodes express}
\begin{itemize}
  \item \textbf{Plan en 3 parties} : 开头 (introduction : le problème) → 正文 (2 ou 3 actions avec 首先、其次、最后) → 结尾 (conclusion avec un appel : 让我们一起…).
  \item \textbf{Connecteurs utiles} : 首先 (d'abord), 其次 (ensuite), 最后 (enfin), 而且 (de plus), 但是 (mais), 所以 (donc).
  \item \textbf{Écriture des caractères} : 环 (clé 王 à gauche + 不), 境 (clé 土 + 竟), 保 (clé 亻+ 呆), 护 (clé 扌+ 户) ; ordre : gauche → droite, haut → bas. Ne pas confondre 境 / 镜, 垃 / 拉, 护 / 沪.
  \item \textbf{Traduction} : repérer d'abord sujet + verbe, placer le complément de lieu avant le verbe (在学校), vérifier les tons et la négation (不 avant 应该).
  \item \textbf{Relecture} : vérifier les tons du pinyin, la place de 不, la ponctuation chinoise 。， et la présence d'un connecteur par paragraphe.
\end{itemize}
\end{bilan}

\begin{aretenir}[Checklist avant le Bac]
\begin{itemize}
  \item $\square$ Je connais le lexique clé : 环境、保护、垃圾、回收、节约、污染、干净、种树.
  \item $\square$ Je sais utiliser 应该 / 不应该 avec la négation \textbf{avant} 应该.
  \item $\square$ Je sais construire une phrase avec 因为 … 所以 … et 如果 … 就 ….
  \item $\square$ Je structure mon texte en trois parties : 开头、正文、结尾.
  \item $\square$ J'utilise au moins trois connecteurs : 首先、其次、最后.
  \item $\square$ J'écris correctement 环、境、保、护、垃、圾 (radicaux et ordre des traits).
  \item $\square$ Je ne confonds pas les caractères proches : 境/镜、垃/拉、护/沪.
  \item $\square$ Je place le complément de lieu avant le verbe : 在学校保护环境.
  \item $\square$ J'écris le pinyin avec les tons sur les voyelles, sans chiffres.
  \item $\square$ J'utilise la ponctuation chinoise pleine chasse ：，。！
  \item $\square$ Je termine par un appel : 让我们一起保护我们的学校！Ràng wǒmen yīqǐ bǎohù wǒmen de xuéxiào !
  \item $\square$ Je relis mon texte : sens, tons, ordre des mots, caractères.
\end{itemize}
\end{aretenir}

\findecours

\end{document}
